% Magnetic Fields — Physics Upper Sixth — Cours signé OnBuch+
% Official MINESEC syllabus. Compiler avec : tectonic PHY07-magnetic-fields.tex
\newcommand{\DOCMATIERE}{Physics}
\newcommand{\DOCNIVEAU}{Upper Sixth}
\newcommand{\DOCMODULE}{Upper Sixth — Physics}
\newcommand{\DOCLECON}{7}
\newcommand{\DOCTITRE}{Magnetic Fields}
% OnBuch+ — préambule « cours signés » ENRICHI (compilé avec tectonic / XeLaTeX)
% Système visuel « Pop » d'OnBuch+ : fond crème (jamais blanc), bordures encre
% épaisses, ombres « dures » décalées, titres Archivo Black en capitales.
% Version MATHÉMATIQUES : graphes (pgfplots/tikz), boîtes pédagogiques
% (définition, propriété, méthode, attention, le savais-tu, exercice résolu,
% exercice, TP, bilan), page de garde et signature OnBuch+ ancrée en bas.
%
% Macros attendues dans main.tex AVANT \input{preamble} :
%   \DOCMATIERE  \DOCNIVEAU  \DOCMODULE  \DOCLECON (numéro)  \DOCTITRE
\documentclass[11pt]{article}

\usepackage{fontspec}
\usepackage[english]{babel}
\usepackage[a4paper,margin=2cm,top=3.4cm,bottom=2.8cm,headheight=1.4cm,footskip=1.3cm]{geometry}
\usepackage{amsmath,amssymb,amsfonts}
\usepackage{mathtools} % \xmapsto, \coloneqq... souvent utilisés par les modèles
\usepackage{cancel} % simplifications barrées (bilans de forces)
% \abs{z} (module, valeur absolue) et \norm{u} (norme), utilisés par les modèles
\providecommand{\abs}{}\let\abs\relax\DeclarePairedDelimiter{\abs}{\lvert}{\rvert}
\providecommand{\norm}{}\let\norm\relax\DeclarePairedDelimiter{\norm}{\lVert}{\rVert}
\usepackage[table]{xcolor} % [table] : fournit aussi \rowcolors (alternance de lignes)
\usepackage{pagecolor}
\usepackage{tikz}
\usetikzlibrary{arrows.meta,positioning,calc,shapes.geometric,shapes.misc,shapes.arrows,decorations.pathmorphing,decorations.pathreplacing,decorations.markings,patterns,shadows,mindmap,trees,fit,backgrounds,matrix,intersections,angles,quotes}
\usepackage{pgfplots}
\usepackage[version=4]{mhchem} % équations nucléaires \ce{^{235}_{92}U}
\usepackage[european,siunitx]{circuitikz} % schémas électriques
\pgfplotsset{compat=1.18}
\usepgfplotslibrary{fillbetween,groupplots} % aires entre courbes (calcul intégral)
\usepackage{enumitem}
\usepackage{fancyhdr}
% [most] charge toutes les bibliothèques tcolorbox (listings, minted, raster,
% documentation...) et gonfle inutilement la mémoire TeX sur un document long
% (~300 boîtes) : on ne charge que ce qui est réellement utilisé.
\usepackage[skins,breakable]{tcolorbox}
\usepackage{graphicx}
\usepackage{colortbl}
\usepackage{booktabs}
\usepackage{tabularx}
\usepackage{diagbox} % case d'en-tête diagonale des tableaux à double entrée
% Colonnes extensibles centrée / gauche / droite (C, L, R), souvent utilisées par les modèles
\newcolumntype{C}{>{\centering\arraybackslash}X}
\newcolumntype{L}{>{\raggedright\arraybackslash}X}
\newcolumntype{R}{>{\raggedleft\arraybackslash}X}
\usepackage{array}
\usepackage{multicol}
\usepackage{siunitx}
% parse-numbers=false : accepte aussi \SI{2,5 \times 10^{-3}}{...}
\sisetup{locale=UK,output-decimal-marker={.},group-minimum-digits=4,parse-numbers=false}
\DeclareSIUnit\ohm{\ensuremath{\symup{\Omega}}} % U+2126 absent de la police
\DeclareSIUnit\division{div} % réglages d'oscilloscope (ms/div, V/div)
\DeclareSIUnit\tr{tr} % tours (vitesse de rotation, ex. \SI{1500}{\tr\per\minute})
\usepackage{qrcode}
% \degree : commande fréquemment utilisée par les modèles pour un angle ou une
% température en dehors de \SI{}{} (non fournie par les paquets chargés).
\providecommand{\degree}{^{\circ}}
% \qed (fin de démonstration), utilisé par les modèles sans amsthm
\providecommand{\qed}{\hfill$\square$}
% \barre : double barre des valeurs interdites dans un tableau de variations (tkz-tab)
\providecommand{\barre}{\ensuremath{\|}}
% environnement proof (amsthm non chargé) utilisé par les modèles
\ifdefined\proof\else\newenvironment{proof}[1][Proof]{\par\noindent\textit{#1.}\ }{\qed\par}\fi
\usepackage{lastpage}
\usepackage{hyperref}
\hypersetup{hidelinks}

% ============================================================
% Palette « Pop » OnBuch+ — mêmes hex que la classe P (plus_ui.dart)
% ============================================================
\definecolor{poporange}{HTML}{F59321}
\definecolor{popdark}{HTML}{E07A0C}
\definecolor{popcream}{HTML}{FFF6E8}
\definecolor{popink}{HTML}{1C1714}
\definecolor{poppurple}{HTML}{7A5AE0}
\definecolor{popgreen}{HTML}{2E9E6A}
\definecolor{poppink}{HTML}{E0457B}
\definecolor{popblue}{HTML}{2F7DE1}
\definecolor{popgold}{HTML}{C98A12}
\definecolor{popmuted}{HTML}{8A7F76}
\definecolor{popline}{HTML}{EADFD2}
\definecolor{popgray}{HTML}{8A7F76}
\colorlet{popgrey}{popgray}
% Alias tolérants (le modèle nomme parfois les couleurs différemment)
\colorlet{popwhite}{white}
\colorlet{popblack}{popink}
\colorlet{popred}{poppink}
\colorlet{popyellow}{popgold}
% Teintes claires pour fonds de tableaux / zones
\colorlet{popblueL}{popblue!10!white}
\colorlet{poporangeL}{poporange!14!white}
\colorlet{popgreenL}{popgreen!12!white}
\colorlet{poppurpleL}{poppurple!12!white}
\colorlet{poppinkL}{poppink!12!white}
\colorlet{popgoldL}{popgold!14!white}

\pagecolor{popcream}

% ============================================================
% Typographie
% ============================================================
\defaultfontfeatures{Ligatures=TeX}
\setmainfont{PlusJakartaSans-Medium.ttf}[Path=../../../fonts/,BoldFont=PlusJakartaSans-Bold.ttf]
% Maths en sans-serif (Fira Math) avec chiffres et lettres droites de Plus
% Jakarta, pour que les formules se fondent dans le texte.
\usepackage[math-style=ISO,bold-style=ISO]{unicode-math}
\setmathfont{FiraMath-Regular.otf}[Path=../../../fonts/]
\setmathfont{PlusJakartaSans-Medium.ttf}[Path=../../../fonts/,range={up/num,up/{Latin,latin},\mathup}]
% (icomma retiré : décimales avec point en anglais)
% Caractères mathématiques Unicode absents des polices de texte (Plus Jakarta,
% Archivo Black) : rendus par leur équivalent mathématique, en texte comme en
% formule — sinon ils s'affichent en carré vide (ex. « Arithmétique dans ℤ »).
\usepackage{newunicodechar}
\newunicodechar{ℝ}{\ensuremath{\mathbb{R}}}
\newunicodechar{ℤ}{\ensuremath{\mathbb{Z}}}
\newunicodechar{ℂ}{\ensuremath{\mathbb{C}}}
\newunicodechar{ℕ}{\ensuremath{\mathbb{N}}}
\newunicodechar{ℚ}{\ensuremath{\mathbb{Q}}}
\newunicodechar{𝔻}{\ensuremath{\mathbb{D}}}
\newunicodechar{α}{\ensuremath{\alpha}}
\newunicodechar{β}{\ensuremath{\beta}}
\newunicodechar{γ}{\ensuremath{\gamma}}
\newunicodechar{δ}{\ensuremath{\delta}}
\newunicodechar{ε}{\ensuremath{\varepsilon}}
\newunicodechar{θ}{\ensuremath{\theta}}
\newunicodechar{λ}{\ensuremath{\lambda}}
\newunicodechar{μ}{\ensuremath{\mu}}
\newunicodechar{σ}{\ensuremath{\sigma}}
\newunicodechar{τ}{\ensuremath{\tau}}
\newunicodechar{φ}{\ensuremath{\varphi}}
\newunicodechar{ψ}{\ensuremath{\psi}}
\newunicodechar{ω}{\ensuremath{\omega}}
\newunicodechar{Σ}{\ensuremath{\Sigma}}
% majuscules produites par \MakeUppercase dans les titres (α-aminés -> Α)
\newunicodechar{Α}{\ensuremath{\alpha}}
\newunicodechar{Β}{\ensuremath{\beta}}
\newunicodechar{Γ}{\ensuremath{\Gamma}}
\newunicodechar{Δ}{\ensuremath{\Delta}}
\newunicodechar{Ε}{\ensuremath{\varepsilon}}
\newunicodechar{Θ}{\ensuremath{\Theta}}
\newunicodechar{Λ}{\ensuremath{\Lambda}}
\newunicodechar{Μ}{\ensuremath{\mu}}
\newunicodechar{Τ}{\ensuremath{\tau}}
\newunicodechar{Φ}{\ensuremath{\Phi}}
\newunicodechar{Ψ}{\ensuremath{\Psi}}
\newunicodechar{Ω}{\ensuremath{\Omega}}
\newunicodechar{↦}{\ensuremath{\mapsto}}
\newunicodechar{∈}{\ensuremath{\in}}
\newunicodechar{∉}{\ensuremath{\notin}}
\newunicodechar{∧}{\ensuremath{\wedge}}
\newunicodechar{⇔}{\ensuremath{\Leftrightarrow}}
\newunicodechar{⟺}{\ensuremath{\Longleftrightarrow}}
\newunicodechar{⇒}{\ensuremath{\Rightarrow}}
\newunicodechar{⟹}{\ensuremath{\Longrightarrow}}
\newunicodechar{∘}{\ensuremath{\circ}}
\newunicodechar{∪}{\ensuremath{\cup}}
\newunicodechar{∩}{\ensuremath{\cap}}
\newunicodechar{≡}{\ensuremath{\equiv}}
\newunicodechar{⊂}{\ensuremath{\subset}}
\newunicodechar{⊥}{\ensuremath{\perp}}
\newunicodechar{∃}{\ensuremath{\exists}}
\newunicodechar{∀}{\ensuremath{\forall}}
\newunicodechar{∛}{\ensuremath{\sqrt[3]{\,}}}
\newunicodechar{∜}{\ensuremath{\sqrt[4]{\,}}}
\newunicodechar{⃗}{\ensuremath{{}^{\to}}}
\newunicodechar{ⁿ}{\ifmmode^{n}\else\textsuperscript{\ensuremath{n}}\fi}
\newunicodechar{ˣ}{\ifmmode^{x}\else\textsuperscript{\ensuremath{x}}\fi}
\newunicodechar{ᵇ}{\ifmmode^{b}\else\textsuperscript{\ensuremath{b}}\fi}
\newunicodechar{ʳ}{\ifmmode^{r}\else\textsuperscript{\ensuremath{r}}\fi}
\newunicodechar{ᵘ}{\ifmmode^{u}\else\textsuperscript{\ensuremath{u}}\fi}
\newunicodechar{ʸ}{\ifmmode^{y}\else\textsuperscript{\ensuremath{y}}\fi}
\newunicodechar{ᵃ}{\ifmmode^{a}\else\textsuperscript{\ensuremath{a}}\fi}
\newunicodechar{ᵉ}{\ifmmode^{e}\else\textsuperscript{\ensuremath{e}}\fi}
\newunicodechar{ᵏ}{\ifmmode^{k}\else\textsuperscript{\ensuremath{k}}\fi}
\newunicodechar{ᵖ}{\ifmmode^{p}\else\textsuperscript{\ensuremath{p}}\fi}
\newunicodechar{ᴺ}{\ifmmode^{N}\else\textsuperscript{\ensuremath{N}}\fi}
\newunicodechar{ᶜ}{\ifmmode^{c}\else\textsuperscript{\ensuremath{c}}\fi}
\newunicodechar{ʰ}{\ifmmode^{h}\else\textsuperscript{\ensuremath{h}}\fi}
\newunicodechar{ᵠ}{\ifmmode^{q}\else\textsuperscript{\ensuremath{q}}\fi}
\newunicodechar{ₙ}{\ifmmode_{n}\else\textsubscript{\ensuremath{n}}\fi}
\newunicodechar{ₐ}{\ifmmode_{a}\else\textsubscript{\ensuremath{a}}\fi}
\newunicodechar{ᵢ}{\ifmmode_{i}\else\textsubscript{\ensuremath{i}}\fi}
\newunicodechar{ᵣ}{\ifmmode_{r}\else\textsubscript{\ensuremath{r}}\fi}
\newunicodechar{ₖ}{\ifmmode_{k}\else\textsubscript{\ensuremath{k}}\fi}
\newunicodechar{ᵦ}{\ifmmode_{\beta}\else\textsubscript{\ensuremath{\beta}}\fi}
\newunicodechar{ₓ}{\ifmmode_{x}\else\textsubscript{\ensuremath{x}}\fi}
\newfontfamily\popdisplay{ArchivoBlack-Regular.ttf}[Path=../../../fonts/]
\newfontfamily\poplogo{SpaceGrotesk-Bold.ttf}[Path=../../../fonts/]
\newfontfamily\popsemi{PlusJakartaSans-SemiBold.ttf}[Path=../../../fonts/]
\newfontfamily\popxbold{PlusJakartaSans-ExtraBold.ttf}[Path=../../../fonts/]
\newfontfamily\popxboldls{PlusJakartaSans-ExtraBold.ttf}[Path=../../../fonts/,LetterSpace=3.0]

\setlist[enumerate]{leftmargin=1.7em,itemsep=5pt,topsep=4pt}
\setlist[itemize]{leftmargin=1.7em,itemsep=4pt,topsep=4pt,label={\color{poporange}\textbullet}}
\setlength{\parindent}{0pt}
\setlength{\parskip}{5pt}
\renewcommand{\arraystretch}{1.25}

% pgfplots : style Pop par défaut
\pgfplotsset{
  every axis/.append style={axis line style={popink,line width=1pt},tick style={popink},
    grid style={popline,line width=0.5pt},label style={font=\small},tick label style={font=\footnotesize}},
  popcurve/.style={poporange,line width=1.8pt,no markers,smooth},
}
% Même style disponible dans un \draw TikZ simple (hors pgfplots)
\tikzset{popcurve/.style={poporange,line width=1.8pt}}
\tikzset{popgrid/.style={popline,very thin}}
\colorlet{popcurve}{poporange} % le nom du style est parfois utilisé comme couleur

% ============================================================
% Pill / squiggle
% ============================================================
\newcommand{\poppill}[3]{%
  \tikz[baseline=-0.55ex]\node[draw=popink,line width=1.2pt,rounded corners=2.6pt,%
    fill=#2,inner xsep=6.5pt,inner ysep=3pt]{%
    \popxboldls\fontsize{7}{7}\selectfont\color{#3}\MakeUppercase{#1}};%
}
\newcommand{\popsquiggle}[1]{%
  \tikz[baseline=-0.4ex]{
    \draw[#1,line width=2.6pt,line cap=round]
      plot [smooth] coordinates {(0,0.09) (0.34,0.01) (0.68,0.21) (1.02,0.01) (1.36,0.10)};
  }%
}
% Mot-clé surligné (marqueur orange sous le texte)
\newcommand{\cle}[1]{{\popxbold\color{popdark}#1}}

% ============================================================
% En-tête / pied
% ============================================================
\pagestyle{fancy}
\fancyhf{}
\renewcommand{\headrulewidth}{0pt}
\renewcommand{\footrulewidth}{0pt}
\lhead{\raisebox{-0.15ex}{\poplogo\fontsize{13}{13}\selectfont\color{popink}OnBuch\color{poporange}+}}
\rhead{\poppill{\DOCMATIERE\ \textbullet\ \DOCNIVEAU\ \textbullet\ Chapter \DOCLECON}{white}{popink}}
\lfoot{\popxbold\fontsize{7.6}{7.6}\selectfont\color{popmuted}\MakeUppercase{Signed course OnBuch+}\ \textbullet\ {\popsemi\DOCTITRE}}
\rfoot{\tikz[baseline=-0.6ex]\node[rounded corners=8.5pt,fill=popink,minimum height=17pt,minimum width=17pt,inner xsep=5pt,inner sep=0]{\popxbold\fontsize{8}{8}\selectfont\color{popcream}\thepage\,/\,\pageref*{LastPage}};}

% ============================================================
% Titres
% ============================================================
\newcommand{\chaptitre}[2]{%
  \begin{tcolorbox}[enhanced,colback=poporange,colframe=popink,boxrule=2.6pt,arc=11pt,
      drop shadow=popink,left=15pt,right=15pt,top=12pt,bottom=11pt,before skip=3pt,after skip=18pt]
    \poppill{#2}{popink}{popcream}\par\vspace{9pt}
    {\raggedright\hyphenpenalty=10000\exhyphenpenalty=10000\popdisplay\fontsize{22}{24}\selectfont\color{popcream}\MakeUppercase{#1}\par}
  \end{tcolorbox}
}
% Section numérotée : pastille orange avec le numéro + titre Archivo
\newcounter{popsec}
\newcommand{\coursec}[1]{%
  \stepcounter{popsec}\setcounter{popsub}{0}%
  \par\vspace{16pt}\noindent
  \begin{minipage}{\linewidth}
  \tikz[baseline=-0.7ex]\node[draw=popink,line width=1.6pt,fill=poporange,rounded corners=4pt,minimum size=22pt,inner sep=0]{\popdisplay\fontsize{12}{12}\selectfont\color{popcream}\thepopsec};%
  \hspace{8pt}{\popdisplay\fontsize{15}{17}\selectfont\color{popink}\MakeUppercase{#1}}\par
  \vspace{4pt}\noindent{\color{poporange}\rule{36pt}{3.6pt}}
  \end{minipage}\par\nopagebreak\vspace{8pt}
}
\newcounter{popsub}
\newcommand{\courssub}[1]{%
  \stepcounter{popsub}%
  \par\vspace{9pt}\noindent{\popxbold\fontsize{11.5}{13}\selectfont\color{popink}{\color{poporange}\thepopsec.\thepopsub}\hspace{6pt}#1}\par\nopagebreak\vspace{4pt}
}

% ============================================================
% Boîtes pédagogiques — même recette : fond blanc, bordure épaisse colorée,
% ombre dure encre, badge plein. Titre optionnel [#1].
% ============================================================
\tcbset{popbase/.style={breakable,enhanced,colback=white,boxrule=2.3pt,arc=9pt,
  shadow={3pt}{-3pt}{0pt}{popink},left=14pt,right=14pt,top=11pt,bottom=11pt,before skip=11pt,after skip=15pt,
  fontupper=\popsemi\fontsize{10.6}{14.5}\selectfont\color{popink},
  fontlower=\popsemi\fontsize{10.6}{14.5}\selectfont\color{popink}}}
% Coche dessinée (le glyphe ✓ n'existe pas dans les polices du document)
\newcommand{\popcheck}[1]{\tikz[baseline=-0.5ex]\draw[#1,line width=1.6pt,line cap=round,line join=round](-0.13,0.0)--(-0.04,-0.09)--(0.13,0.1);}
\providecommand{\coche}{\popcheck{popgreen}}
\providecommand{\resultat}[1]{\par\smallskip\noindent\fcolorbox{popgreen}{popgreenL}{\parbox{\dimexpr\linewidth-2\fboxsep-2\fboxrule\relax}{\textbf{Result.} #1}}\par\smallskip}
\newcommand{\popboxhead}[3]{\poppill{#1}{#2}{white}\ifx\relax#3\relax\else\hspace{6pt}{\popxbold\fontsize{10.6}{12}\selectfont\color{popink}#3}\fi\par\vspace{7pt}}

\newtcolorbox{aretenir}[1][]{popbase,colframe=popgreen,before upper={\popboxhead{Key points}{popgreen}{#1}}}
\newtcolorbox{exemplebox}[1][]{popbase,colframe=poporange,before upper={\popboxhead{Exemple}{poporange}{#1}}}
\newtcolorbox{definition}[1][]{popbase,colframe=popblue,before upper={\popboxhead{Definition}{popblue}{#1}}}
\newtcolorbox{propriete}[1][]{popbase,colframe=poppurple,before upper={\popboxhead{Property}{poppurple}{#1}}}
\newtcolorbox{methode}[1][]{popbase,colframe=popgold,before upper={\popboxhead{Method}{popgold}{#1}}}
\newtcolorbox{attention}[1][]{popbase,colframe=poppink,before upper={\popboxhead{Warning}{poppink}{#1}}}
\newtcolorbox{experience}[1][]{popbase,colframe=popblue,colback=popblueL,before upper={\popboxhead{Experiment}{popblue}{#1}}}
% « Le savais-tu ? » : carte violette pleine (contexte, culture, Cameroun)
\newtcolorbox{savaistu}[1][]{popbase,colback=poppurple,colframe=popink,
  fontupper=\popsemi\fontsize{10.4}{14.2}\selectfont\color{white},
  before upper={\poppill{Did you know?}{white}{poppurple}\ifx\relax#1\relax\else\hspace{6pt}{\popxbold\color{white}#1}\fi\par\vspace{7pt}}}
% Exercice résolu : énoncé en haut, \tcblower puis la solution sur fond crème
\newcounter{popexo}\newcounter{popexr}
\newtcolorbox{exoresolu}[1][]{popbase,colframe=poporange,colbacklower=popcream,
  segmentation style={popink,line width=1.2pt,dashed},
  before upper={\stepcounter{popexr}\popboxhead{Worked example \thepopexr}{poporange}{#1}},
  before lower={\poppill{Solution}{popink}{popcream}\par\vspace{6pt}}}
% Exercice (non résolu en place) — la correction est dans la partie Corrigés
\newtcolorbox{exercice}[1][]{popbase,colframe=popink,
  before upper={\stepcounter{popexo}\popboxhead{Exercise \thepopexo}{popink}{#1}}}
% Corrigé
\newtcolorbox{corrige}[1][]{popbase,colframe=popgreen,colback=popgreenL,
  before upper={\popboxhead{Solution}{popgreen}{#1}}}
% Objectifs / prérequis / bilan
\newtcolorbox{objectifs}{popbase,colframe=popink,colback=poporangeL,
  before upper={\popboxhead{Chapter objectives}{popink}{}}}
\newtcolorbox{prerequis}{popbase,colframe=popink,colback=popblueL,
  before upper={\popboxhead{Prerequisites}{popblue}{}}}
\newtcolorbox{bilan}{popbase,colframe=popink,colback=white,boxrule=2.8pt,
  before upper={\popboxhead{Fiche bilan}{poporange}{L'essentiel à connaître pour l'examen}}}
% Figure encadrée avec légende
\newtcolorbox{popfigure}[1][]{enhanced,breakable=false,colback=white,colframe=popline,boxrule=1.4pt,arc=8pt,
  left=10pt,right=10pt,top=10pt,bottom=8pt,before skip=10pt,after skip=12pt,halign=center,
  lower separated=false,halign lower=center,
  fontlower=\popsemi\fontsize{9}{11.5}\selectfont\color{popmuted},
  #1}
\newcommand{\legende}[1]{\par\vspace{6pt}{\popsemi\fontsize{9}{11.5}\selectfont\color{popmuted}#1}}

% ============================================================
% Signature OnBuch+
% ============================================================
\newcommand{\poplogoinline}[1]{{\poplogo\fontsize{#1}{#1}\selectfont\color{popink}OnBuch\color{poporange}+}}
\newcommand{\signatureonbuch}{%
  \begin{tcolorbox}[enhanced,colback=popink,colframe=popink,boxrule=2.2pt,arc=10pt,
      drop shadow=poporange,left=14pt,right=14pt,top=10pt,bottom=10pt,before skip=0pt,after skip=0pt]
    \tikz[baseline=-0.7ex]\node[circle,fill=popgreen,draw=popcream,line width=1.3pt,minimum size=22pt,inner sep=0]{\popcheck{white}};%
    \hspace{9pt}\begin{minipage}[c]{0.62\linewidth}
      {\popxbold\fontsize{12}{13}\selectfont\color{popcream}Signed\ \textbullet\ {\poplogo OnBuch\color{poporange}+}}\par\vspace{2pt}
      {\raggedright\popsemi\fontsize{8}{10.5}\selectfont\color{popline}\MakeUppercase{OnBuch+ teaching team}\\\MakeUppercase{Follows the official MINESEC syllabus}\par}
    \end{minipage}\hfill
    {\poplogo\fontsize{22}{22}\selectfont\color{popcream}OnBuch\color{poporange}+}
  \end{tcolorbox}%
}

% ============================================================
% Page de garde
% #1 titre · #2 matière · #3 niveau · #4 module · #5 n° leçon · #6 durée ·
% #7 description / objectifs (texte ou itemize)
% ============================================================
\newcommand{\pagegarde}[7]{%
  \thispagestyle{empty}
  \newgeometry{a4paper,margin=1.7cm,top=1.6cm,bottom=1.4cm}%
  \begin{tikzpicture}[remember picture,overlay]
    \fill[poporange!18] ([xshift=-3.2cm,yshift=-3.6cm]current page.north east) circle (4.6cm);
    \fill[poppurple!14] ([xshift=2.4cm,yshift=7.6cm]current page.south west) circle (3.8cm);
  \end{tikzpicture}%
  {\poplogo\fontsize{30}{30}\selectfont\color{popink}OnBuch\color{poporange}+}\hfill
  \raisebox{0.6ex}{\poppill{Signed course \textbullet\ Premium}{popink}{popcream}}\par
  \vspace{1pt}\popsquiggle{poppurple}\par
  \vspace{8pt}
  \begin{tcolorbox}[enhanced,colback=poporange,colframe=popink,boxrule=2.8pt,arc=13pt,
      drop shadow=popink,left=18pt,right=18pt,top=12pt,bottom=13pt,after skip=10pt]
    \poppill{#2 \textbullet\ #3}{popink}{popcream}\hspace{5pt}\poppill{#4}{white}{popink}\par\vspace{8pt}
    {\popdisplay\fontsize{14}{14}\selectfont\color{popink}CHAPTER #5}\par\vspace{4pt}
    {\raggedright\hyphenpenalty=10000\exhyphenpenalty=10000\popdisplay\fontsize{25}{27}\selectfont\color{popcream}\MakeUppercase{#1}\par}
    \vspace{8pt}
    {\popxbold\fontsize{9.5}{9.5}\selectfont\color{popink}\MakeUppercase{Suggested time: #6}}
  \end{tcolorbox}
  \begin{tcolorbox}[enhanced,colback=white,colframe=popink,boxrule=2.2pt,arc=10pt,
      drop shadow=popink,left=16pt,right=16pt,top=10pt,bottom=10pt,after skip=8pt]
    \poppill{In this chapter}{poppurple}{white}\par\vspace{5pt}
    \popsemi\fontsize{9.3}{12.6}\selectfont\color{popink}#7
  \end{tcolorbox}
  \noindent\begin{minipage}[t]{0.487\linewidth}
    \begin{tcolorbox}[enhanced,colback=popgreen,colframe=popink,boxrule=2.2pt,arc=10pt,
        drop shadow=popink,left=11pt,right=11pt,top=10pt,bottom=10pt,height=3.1cm,valign=center]
      \tcbox[colback=white,colframe=popink,boxrule=1.3pt,arc=5pt,boxsep=0pt,nobeforeafter,
        left=3pt,right=3pt,top=3pt,bottom=3pt,valign=center]{\qrcode[height=1.9cm]{https://play.google.com/store/apps/details?id=com.luvvix.onbuch&hl=en}}%
      \hspace{8pt}\begin{minipage}[c]{0.5\linewidth}
        {\popdisplay\fontsize{11}{12.5}\selectfont\color{white}\MakeUppercase{Download OnBuch}}\par\vspace{4pt}
        {\raggedright\popsemi\fontsize{8.4}{11}\selectfont\color{white}Free \textbullet\ Google Play\\AI tutor, past papers, results\par}
      \end{minipage}
    \end{tcolorbox}
  \end{minipage}\hfill
  \begin{minipage}[t]{0.487\linewidth}
    \begin{tcolorbox}[enhanced,colback=poppurple,colframe=popink,boxrule=2.2pt,arc=10pt,
        drop shadow=popink,left=11pt,right=11pt,top=10pt,bottom=10pt,height=3.1cm,valign=center]
      \tikz[baseline=-0.7ex]\node[circle,fill=white,draw=popink,line width=1.3pt,minimum size=1.6cm,inner sep=0]{\popdisplay\fontsize{24}{24}\selectfont\color{poppurple}+};%
      \hspace{8pt}\begin{minipage}[c]{0.6\linewidth}
        {\popdisplay\fontsize{11}{12.5}\selectfont\color{white}\MakeUppercase{Go OnBuch+}}\par\vspace{4pt}
        {\raggedright\popsemi\fontsize{8.4}{11}\selectfont\color{white}All signed courses, unlimited AI tutor, PDF sheets — OnBuch+ tab in the app\par}
      \end{minipage}
    \end{tcolorbox}
  \end{minipage}
  \vspace{8pt}
  \popcontenu
  \vfill
  \signatureonbuch
  \restoregeometry
  \newpage
}

% Rangée « Ce cours contient » (page de garde)
\newcommand{\poptile}[3]{% #1 couleur #2 gros texte #3 légende
  \begin{tcolorbox}[enhanced,colback=#1,colframe=popink,boxrule=2pt,arc=9pt,shadow={2.5pt}{-2.5pt}{0pt}{popink},
      left=6pt,right=6pt,top=9pt,bottom=9pt,halign=center,valign=center,height=1.75cm,width=\linewidth,nobeforeafter]
    {\popdisplay\fontsize{12.5}{13}\selectfont\color{white}\MakeUppercase{#2}}\par\vspace{3pt}
    {\centering\popsemi\fontsize{7.8}{9.5}\selectfont\color{white}#3\par}
  \end{tcolorbox}}
\newcommand{\popcontenu}{%
  \par\noindent{\popxboldls\fontsize{7.5}{7.5}\selectfont\color{popmuted}THIS SIGNED COURSE CONTAINS}\par\vspace{7pt}
  \noindent\begin{minipage}[t]{0.235\linewidth}\poptile{popblue}{Course}{complete, illustrated, step by step}\end{minipage}\hfill
  \begin{minipage}[t]{0.235\linewidth}\poptile{poporange}{Methods}{and worked examples}\end{minipage}\hfill
  \begin{minipage}[t]{0.235\linewidth}\poptile{poppink}{Exercises}{exam style + solutions}\end{minipage}\hfill
  \begin{minipage}[t]{0.235\linewidth}\poptile{popgreen}{Summary sheet}{the essentials to remember}\end{minipage}\par}

% Sommaire
\newtcolorbox{sommaire}{enhanced,colback=white,colframe=popink,boxrule=2.2pt,arc=10pt,
  drop shadow=popink,left=18pt,right=18pt,top=13pt,bottom=13pt,before skip=6pt,after skip=20pt,
  before upper={\poppill{Contents}{poppurple}{white}\par\vspace{9pt}\popsemi\fontsize{10.6}{14.5}\selectfont\color{popink}}}

% Signature de fin de cours — poussée tout en bas de la dernière page
\newcommand{\findecours}{%
  \par\vfill\nopagebreak
  \begin{center}\popsquiggle{poporange}\end{center}\vspace{4pt}
  \signatureonbuch
}
\AtBeginDocument{\def\llbracket{\mathopen{[\mkern-3.5mu[}}\def\rrbracket{\mathclose{]\mkern-3.5mu]}}} % intervalles d'entiers (glyphe ⟦ absent de la police)
% Glyphes absents de Fira Math : redéfinis à partir de glyphes disponibles
\AtBeginDocument{%
  \def\setminus{\mathbin{\backslash}}\let\smallsetminus\setminus
  \def\vdots{\mathord{\vbox{\baselineskip3.2pt\lineskiplimit0pt\kern4pt\hbox{.}\hbox{.}\hbox{.}}}}%
  \def\ddots{\mathinner{\mkern1mu\raise6.5pt\hbox{.}\mkern2mu\raise3.6pt\hbox{.}\mkern2mu\raise0.7pt\hbox{.}\mkern1mu}}%
  \def\triangle{\mathord{\scalebox{0.9}{$\symup{\Delta}$}}}%
  \def\nabla{\mathord{\raisebox{\depth}{\scalebox{1}[-1]{$\symup{\Delta}$}}}}%
  \def\checkmark{\popcheck{popdark}}%
  \def\star{\mathbin{\ast}}\def\bigstar{\mathord{\ast}}%
  \def\diamond{\mathbin{\scalebox{0.6}{$\Diamond$}}}%
  \def\bigcup{\mathop{\mathchoice{\raisebox{-0.3em}{\scalebox{1.7}{$\cup$}}}{\scalebox{1.25}{$\cup$}}{\cup}{\cup}}\displaylimits}%
  \def\bigcap{\mathop{\mathchoice{\raisebox{-0.3em}{\scalebox{1.7}{$\cap$}}}{\scalebox{1.25}{$\cap$}}{\cap}{\cap}}\displaylimits}%
  \def\lesseqqgtr{\mathrel{\lessgtr}}\def\bigoplus{\mathop{\oplus}\displaylimits}%
}
\providecommand{\ssi}{if and only if} % abréviation fréquente
\AtBeginDocument{\providecommand{\wideparen}{\overparen}} % arc de cercle

%% FIGFIX-BEGIN v4 — correctifs globaux des figures (docs/onbuch-generation/tools/figfix)
\usepackage{adjustbox}\usepackage{etoolbox}
% toute figure TikZ/pgfplots plus large que la ligne est réduite à la largeur disponible
\BeforeBeginEnvironment{tikzpicture}{\begin{adjustbox}{max width=\linewidth}}
\AfterEndEnvironment{tikzpicture}{\end{adjustbox}}
% légendes pgfplots lisibles par-dessus les courbes
\pgfplotsset{every axis/.append style={legend style={fill=white,fill opacity=0.92,text opacity=1,draw=black!15,inner sep=3pt}}}
% étiquettes d'arêtes (syntaxe edge["..."]) avec fond blanc
\tikzset{every edge quotes/.append style={fill=white,inner sep=1.2pt}}
%% FIGFIX-END
\begin{document}

\pagegarde{Magnetic Fields}{Physics}{Upper Sixth}{Upper Sixth — Physics}{7}{4 periods}{This chapter introduces the magnetic field, its sources and its effects: field patterns of current-carrying conductors, forces on conductors and moving charges, torque on coils and the electric motor, the laws of Biot-Savart and Ampère, the Hall effect and magnetic materials. It provides the theoretical foundation for electromagnetism and electromagnetic induction studied later in Upper Sixth.
\vspace{4pt}
\begin{multicols}{2}\small
\begin{itemize}
\item By the end of this chapter I can define magnetic flux density B, state its unit (the tesla) and represent magnetic field lines.
\item By the end of this chapter I can sketch and describe the magnetic field patterns of a long straight wire, a plane circular coil and a solenoid.
\item By the end of this chapter I can apply F = BIL sin θ and Fleming's left-hand rule to determine the magnitude and direction of the force on a current-carrying conductor.
\item By the end of this chapter I can derive and use T = NAIB for the torque on a rectangular coil and explain the working of dc and ac electric motors.
\item By the end of this chapter I can state and apply the Biot-Savart law and Ampère's law to calculate B for a long straight wire, a long solenoid and a plane circular coil.
\item By the end of this chapter I can calculate the force per unit length between two parallel current-carrying conductors and link it to the definition of the ampere.
\end{itemize}\end{multicols}}

\begin{sommaire}
\begin{multicols}{2}
\begin{enumerate}
\item Magnetic Fields and Forces on Current-Carrying Conductors
\item Torque on a Coil and the Electric Motor
\item Sources of Magnetic Fields: Biot-Savart, Ampère and Forces Between Conductors
\item Moving Charges, Hall Effect and Magnetic Materials
\item Integration activity
\item Methods and worked examples
\item Exercises
\item Solutions to the exercises
\item Summary sheet
\end{enumerate}
\end{multicols}
\end{sommaire}

\begin{objectifs}
\begin{itemize}
\item By the end of this chapter I can define magnetic flux density B, state its unit (the tesla) and represent magnetic field lines.
\item By the end of this chapter I can sketch and describe the magnetic field patterns of a long straight wire, a plane circular coil and a solenoid.
\item By the end of this chapter I can apply F = BIL sin θ and Fleming's left-hand rule to determine the magnitude and direction of the force on a current-carrying conductor.
\item By the end of this chapter I can derive and use T = NAIB for the torque on a rectangular coil and explain the working of dc and ac electric motors.
\item By the end of this chapter I can state and apply the Biot-Savart law and Ampère's law to calculate B for a long straight wire, a long solenoid and a plane circular coil.
\item By the end of this chapter I can calculate the force per unit length between two parallel current-carrying conductors and link it to the definition of the ampere.
\item By the end of this chapter I can apply F = Bqv sin θ to a moving charge, describe circular and helical paths, and explain the measurement of the specific charge e/m.
\item By the end of this chapter I can explain the Hall effect, derive the Hall voltage, and use it to determine charge carrier density.
\item By the end of this chapter I can distinguish dia-, para- and ferromagnetic materials and explain the principle of magnetic shielding.
\item By the end of this chapter I can combine electric and magnetic forces into the Lorentz force and analyse crossed-field (velocity selector) situations.
\end{itemize}
\end{objectifs}

\begin{prerequis}
\begin{itemize}
\item Electric current, potential difference and Ohm's law (Lower Sixth, current electricity)
\item Uniform circular motion: centripetal acceleration and force F = mv²/r (Lower Sixth, mechanics)
\item Vectors: resolution into components, cross-product direction via right-hand rules (Lower Sixth, mechanics/mathematics)
\item Electric field concept and force on a charge in an electric field F = qE (Lower Sixth, electrostatics)
\item Trigonometry (sin θ, cos θ) and use of scientific notation with powers of ten
\end{itemize}
\end{prerequis}

\newpage

\coursec{Magnetic Fields and Forces on Current-Carrying Conductors}

In Douala and Yaoundé, every \emph{petit générateur} that lights a home during a power cut, every loudspeaker at a makossa concert, and every hospital MRI scanner in our reference hospitals relies on one invisible actor: the \cle{magnetic field}. Yet the same field that spins a motor can also wipe the magnetic strip of a bank card or disturb a compass used by surveyors in the bush. How can a simple coil of copper wire carrying current lift a car in a scrapyard electromagnet? And how do engineers protect sensitive equipment from stray magnetic fields? This chapter gives you the laws and tools to answer these questions \emph{quantitatively}.

We begin with the most fundamental question: what is a magnetic field, how do we describe it, and what force does it exert on a conductor carrying a current?

\courssub{Magnets, Poles and the Notion of Magnetic Field}

\textbf{Observation.} Bring a small compass near a bar magnet: the needle turns, even though nothing touches it. Suspend a bar magnet freely and it settles roughly along the north--south direction of the Earth --- this is why surveyors and travellers have used compasses for centuries. Clearly, a magnet modifies the space around it.

Every magnet has two regions where its effect is strongest, called \cle{poles}: a \textbf{north pole} (N) and a \textbf{south pole} (S). Two facts are easily verified with two bar magnets:
\begin{itemize}
\item Like poles (N--N or S--S) \textbf{repel}; unlike poles (N--S) \textbf{attract}.
\item Poles always come in pairs: cutting a magnet in half produces two complete magnets, each with its own N and S. Isolated magnetic poles (``magnetic monopoles'') have never been observed.
\end{itemize}

\begin{definition}[Magnetic field]
A \cle{magnetic field} is a region of space in which a magnetic pole, a moving charge or a current-carrying conductor experiences a \textbf{magnetic force}. A small compass needle placed in the field aligns itself with the field: the direction of the field at a point is the direction in which the \textbf{north pole of a compass needle points}.
\end{definition}

\courssub{Magnetic Field Lines}

\textbf{Observation.} Sprinkle iron filings on a card placed over a bar magnet and tap the card gently: the filings arrange themselves into smooth curves running from one pole to the other. These curves reveal the structure of the field.

\begin{definition}[Field line]
A \cle{magnetic field line} is a line whose tangent at every point gives the \textbf{direction of the magnetic field} at that point. By convention, field lines are directed from the \textbf{north pole to the south pole} outside a magnet, and they form \textbf{closed loops} (they continue from S to N \emph{inside} the magnet).
\end{definition}

\begin{propriete}[Characteristics of magnetic field lines]
\begin{enumerate}
\item The tangent to a field line at a point gives the direction of $\vec{B}$ at that point.
\item The \textbf{density} of the lines (number of lines per unit area perpendicular to them) is proportional to the \textbf{magnitude} of the field: where lines crowd together, the field is strong.
\item Field lines \textbf{never cross}: at any point the field has one unique direction, so two tangents cannot exist at the same point.
\item Field lines form \textbf{closed loops} without beginning or end (a consequence of the non-existence of isolated poles).
\end{enumerate}
\end{propriete}

\courssub{Magnetic Flux Density and the Tesla}

A compass tells us the \emph{direction} of the field, but physics needs a \emph{quantitative} description: how strong is the field, and in which direction does it act? For this we introduce a vector quantity.

\begin{definition}[Magnetic flux density]
The \cle{magnetic flux density} $\vec{B}$ is the vector quantity that characterises the magnetic field at a point. Its direction is that of the field (the direction pointed to by the north pole of a compass needle); its magnitude $B$ measures the strength of the field. The SI unit of $B$ is the \cle{tesla} (symbol T). As we shall justify below, a conductor of length $1\ \mathrm{m}$ carrying a current of $1\ \mathrm{A}$ placed perpendicular to a field of $1\ \mathrm{T}$ experiences a force of $1\ \mathrm{N}$:
\[
1\ \mathrm{T} = 1\ \mathrm{N\,A^{-1}\,m^{-1}}.
\]
\end{definition}

\begin{aretenir}[Orders of magnitude]
\begin{center}
\begin{tabularx}{\linewidth}{|l|X|}\hline
\rowcolor{popblueL} \textbf{Source of field} & \textbf{Approximate magnitude of $B$} \\ \hline
Earth's magnetic field (in Cameroon) & $\approx 5\times 10^{-5}\ \mathrm{T}$ \\ \hline
Small bar magnet (at its surface) & $\approx 10^{-2}\ \mathrm{T}$ \\ \hline
Strong laboratory electromagnet & $1$ to $2\ \mathrm{T}$ \\ \hline
MRI scanner (medical imaging) & $1.5$ to $3\ \mathrm{T}$ \\ \hline
\end{tabularx}
\end{center}
The tesla is a \emph{large} unit: the Earth's field is about twenty thousand times weaker than $1\ \mathrm{T}$.
\end{aretenir}

\begin{savaistu}[Nikola Tesla]
The unit honours Nikola Tesla (1856--1943), the engineer whose work on alternating current made modern power distribution possible --- the very principle behind the transformers you see on electricity poles along our roads. A field of $1\ \mathrm{T}$ is strong enough to make a $10\ \mathrm{cm}$ wire carrying $10\ \mathrm{A}$ feel a force of $1\ \mathrm{N}$, roughly the weight of an orange.
\end{savaistu}

\courssub{Oersted's Experiment: Currents Create Magnetic Fields}

\begin{experience}[Oersted's experiment (1820)]
Place a straight wire above a compass needle, parallel to the needle, with no current flowing. The needle points north as usual. Now pass a current through the wire: \textbf{the needle deflects}. Reverse the current: the needle deflects the \emph{other way}. Switch the current off: the needle returns to north.
\end{experience}

\textbf{Conclusion.} An electric current produces a magnetic field around it. The direction of the field depends on the direction of the current. This discovery by the Danish physicist Hans Christian Oersted was the birth of \cle{electromagnetism}: electricity and magnetism, long thought to be separate, are two faces of one phenomenon. Every application mentioned in our opening --- generators, loudspeakers, electromagnets --- rests on this single fact.

\courssub{Field Pattern of a Long Straight Conductor}

\textbf{Observation.} Pass a vertical wire through a horizontal card sprinkled with iron filings and send a strong current through it: the filings form \textbf{concentric circles} centred on the wire. A compass placed on the card is always tangent to one of these circles.

\begin{propriete}[Field of a long straight wire]
The magnetic field lines around a long straight current-carrying conductor are \textbf{concentric circles} in planes perpendicular to the wire, centred on the wire. The field is \textbf{stronger closer to the wire} (the circles are more tightly packed near the wire) and its direction is given by the \textbf{right-hand grip rule}.
\end{propriete}

\begin{methode}[The right-hand grip rule (straight wire)]
\begin{enumerate}
\item Grasp the wire with your \textbf{right hand}.
\item Point your \textbf{thumb} along the direction of the \textbf{conventional current} $I$.
\item Your \textbf{curled fingers} then show the direction of the magnetic field lines $\vec{B}$.
\end{enumerate}
\end{methode}

For diagrams we use the \textbf{dot--cross convention} for the wire seen in cross-section: a \textbf{dot} $\odot$ means current flowing \emph{out of} the page (towards you, like the tip of an approaching arrow); a \textbf{cross} $\otimes$ means current flowing \emph{into} the page (like the feathers of a receding arrow).

\begin{popfigure}
\begin{center}
\begin{tikzpicture}[scale=1.0]
% Left: current out of page
\draw[popline, dashed] (0,0) circle (0.55);
\draw[popline, dashed] (0,0) circle (1.05);
\draw[popline, dashed] (0,0) circle (1.55);
\draw[popblue, thick, ->] (0.55,0) arc (0:70:0.55);
\draw[popblue, thick, ->] (-1.05,0) arc (180:250:1.05);
\draw[popblue, thick, ->] (0,1.55) arc (90:160:1.55);
\filldraw[poporange] (0,0) circle (0.14);
\filldraw[popdark] (0,0) circle (0.045);
\node[popdark] at (0,-2.0) {$I$ out of page: $B$ anticlockwise};
% Right: current into page
\begin{scope}[xshift=6.5cm]
\draw[popline, dashed] (0,0) circle (0.55);
\draw[popline, dashed] (0,0) circle (1.05);
\draw[popline, dashed] (0,0) circle (1.55);
\draw[popgreen, thick, ->] (0.55,0) arc (0:-70:0.55);
\draw[popgreen, thick, ->] (-1.05,0) arc (180:110:1.05);
\draw[popgreen, thick, ->] (0,1.55) arc (90:20:1.55);
\filldraw[poporange] (0,0) circle (0.14);
\draw[popdark, thick] (-0.05,-0.05) -- (0.05,0.05);
\draw[popdark, thick] (-0.05,0.05) -- (0.05,-0.05);
\node[popdark] at (0,-2.0) {$I$ into page: $B$ clockwise};
\end{scope}
\end{tikzpicture}
\end{center}
\legende{Magnetic field lines around a long straight conductor seen in cross-section. The circles are closer together near the wire, where the field is stronger.}
\end{popfigure}

\courssub{Field Pattern of a Plane Circular Coil (Loop)}

Bend the straight wire into a circular loop and pass a current through it. Applying the right-hand grip rule to \emph{each small portion} of the loop shows that inside the loop all the contributions point the \emph{same way} through the loop, while outside they loop back around. The result:

\begin{propriete}[Field of a circular loop]
The field lines of a plane circular coil pass \textbf{through the loop} (nearly straight and close together near the centre, where the field is almost uniform) and close back \textbf{around the outside}. The face from which the lines \textbf{emerge} behaves as a \textbf{north pole}; the face into which they \textbf{enter} behaves as a \textbf{south pole}.
\end{propriete}

\textbf{Polarity rule for a face of the loop.} Look at one face of the loop: if the current circulates \textbf{anticlockwise} around that face, the face is a \textbf{north} pole; if \textbf{clockwise}, it is a \textbf{south} pole. (Memory aid: anticlockwise $\to$ N, clockwise $\to$ S.)

\courssub{Field Pattern of a Solenoid}

A \cle{solenoid} is a long coil of many turns of insulated wire wound on a cylindrical former. It is the most important source of controlled magnetic fields in technology (electromagnets, relays, MRI coils).

\begin{propriete}[Field of a solenoid carrying a current]
\begin{itemize}
\item \textbf{Inside} the solenoid (far from the ends), the field is \textbf{uniform}: the lines are parallel, equally spaced and directed along the axis.
\item \textbf{Outside}, the pattern is identical to that of a \textbf{bar magnet}: lines emerge from one end (the \textbf{north pole}) and re-enter at the other end (the \textbf{south pole}).
\end{itemize}
\textbf{Pole rule:} grasp the solenoid with the right hand, fingers curled in the direction of the current in the turns; the \textbf{thumb points to the north pole} (the direction of $\vec{B}$ inside).
\end{propriete}

\begin{popfigure}
\begin{center}
\begin{tikzpicture}[scale=0.95]
% Circular loop
\begin{scope}[xshift=-4.2cm]
\draw[popdark, very thick] (0,0) ellipse (0.35 and 1.1);
\draw[popblue, thick, ->] (-2.2,0) -- (-0.5,0);
\draw[popblue, thick] (-0.5,0) -- (0.35,0);
\draw[popblue, thick, ->] (0.35,0) -- (2.2,0);
\draw[popblue, thick] (2.2,0.7) arc (60:300:0.8);
\draw[popblue, thick, ->] (2.2,-0.7) arc (-60:-120:0.8);
\draw[popblue, thick] (-2.2,-0.7) arc (240:120:0.8);
\node[popdark] at (0,1.5) {\textbf{N}};
\node[popdark] at (-2.3,1.2) {\textbf{S}};
\node[popmuted] at (0,-1.6) {circular loop};
\end{scope}
% Solenoid
\begin{scope}[xshift=3.2cm]
\foreach \x in {-1.2,-0.8,-0.4,0,0.4,0.8,1.2}{
  \draw[popdark, thick] (\x,0) ellipse (0.16 and 0.62);}
\draw[popgreen, thick, ->] (-2.6,0) -- (-1.4,0);
\draw[popgreen, thick, ->] (1.4,0) -- (2.6,0);
\draw[popgreen, thick] (2.6,0) .. controls (3.0,1.4) and (-3.0,1.4) .. (-2.6,0);
\draw[popgreen, thick, ->] (0,1.42) -- (0.05,1.42);
\draw[popgreen, thick] (2.6,0) .. controls (3.0,-1.4) and (-3.0,-1.4) .. (-2.6,0);
\draw[popgreen, thick, ->] (-0.05,-1.42) -- (0,-1.42);
\node[popdark] at (1.7,0.35) {\textbf{N}};
\node[popdark] at (-1.7,0.35) {\textbf{S}};
\node[popmuted] at (0,-1.7) {solenoid};
\end{scope}
\end{tikzpicture}
\end{center}
\legende{Field patterns of a plane circular loop (left) and a solenoid (right). Outside, the solenoid behaves exactly like a bar magnet; inside, its field is uniform.}
\end{popfigure}

\courssub{Force on a Current-Carrying Conductor in a Uniform Field}

\textbf{Observation.} Place a flexible wire between the poles of a strong horseshoe magnet so that the wire is perpendicular to the field. With no current, nothing happens. Pass a current: the wire \emph{jumps} out of the gap. Reverse the current (or the magnet): the wire jumps the \emph{other way}. Turn the wire until it lies \emph{along} the field: no motion at all, whatever the current.

A magnetic field therefore exerts a \cle{force} on a current-carrying conductor --- this is the \textbf{motor effect}, the working principle of every electric motor and loudspeaker.

Careful measurements (varying $I$, the length $L$ in the field, $B$, and the angle $\theta$ between the conductor and $\vec{B}$) lead to the following law.

\begin{propriete}[Force on a current-carrying conductor (proof not examinable)]
A straight conductor of length $L$ carrying a current $I$, placed in a \textbf{uniform} magnetic field $\vec{B}$ and making an angle $\theta$ with $\vec{B}$, experiences a force of magnitude
\[
\boxed{\,F = BIL\sin\theta\,}
\]
where $F$ is in newtons (N), $B$ in tesla (T), $I$ in amperes (A) and $L$ in metres (m). The force is perpendicular to \textbf{both} the conductor and the field.
\begin{itemize}
\item \textbf{Maximum force} $F_{\max} = BIL$ when the conductor is \textbf{perpendicular} to the field ($\theta = 90^{\circ}$, $\sin\theta = 1$).
\item \textbf{Zero force} when the conductor is \textbf{parallel} to the field ($\theta = 0^{\circ}$ or $180^{\circ}$, $\sin\theta = 0$).
\end{itemize}
\textbf{Conditions of validity:} uniform field, straight conductor, $L$ = length of conductor actually lying in the field.
\end{propriete}

Notice that this law gives operational meaning to the definition of the tesla: with $\theta = 90^{\circ}$, $B = \dfrac{F}{IL}$, so $1\ \mathrm{T} = 1\ \mathrm{N\,A^{-1}\,m^{-1}}$.

\begin{attention}[Zero, not maximum, when parallel!]
A very common GCE error is to write that the force is \emph{maximum} when the conductor is parallel to $\vec{B}$. It is the opposite: parallel $\Rightarrow$ $\sin\theta = 0 \Rightarrow$ $F = 0$. The force is maximum when the conductor is \textbf{perpendicular} to the field. Sketch the angle $\theta$ between the wire and $\vec{B}$ before substituting into $F = BIL\sin\theta$.
\end{attention}

\begin{popfigure}
\begin{center}
\begin{tikzpicture}
\begin{axis}[
  width=9.5cm, height=5.5cm,
  xlabel={$\theta$ (degrees)}, ylabel={$F$ (N)},
  xmin=0, xmax=180, ymin=0, ymax=1.15,
  xtick={0,30,60,90,120,150,180},
  ytick={0,0.5,1.0},
  grid=both, grid style={popline, dashed},
  axis lines=left, thick]
\addplot[popblue, very thick, domain=0:180, samples=100] {sin(x)};
\addplot[poporange, dashed, domain=0:180] {1};
\node[poporange] at (axis cs:150,1.06) {\small $F_{\max}=BIL$};
\draw[popdark, dashed] (axis cs:90,0) -- (axis cs:90,1);
\node[popdark] at (axis cs:90,-0.13) {\small max};
\end{axis}
\end{tikzpicture}
\end{center}
\legende{Variation of the force with the angle $\theta$ between the conductor and the field (here $BIL = 1\ \mathrm{N}$): $F = BIL\sin\theta$ is maximum at $90^{\circ}$ and zero at $0^{\circ}$ and $180^{\circ}$.}
\end{popfigure}

\courssub{Direction of the Force: Fleming's Left-Hand Rule}

The magnitude of the force comes from $F = BIL\sin\theta$; its \textbf{direction} is given by \cle{Fleming's left-hand rule}.

\begin{methode}[Fleming's left-hand rule, step by step]
\begin{enumerate}
\item Spread the \textbf{thumb}, \textbf{first finger} and \textbf{second finger} of your \textbf{left hand} so that all three are mutually perpendicular (like the three edges of a box meeting at one corner).
\item Point the \textbf{first finger} along the \textbf{F}ield $\vec{B}$ (north to south).
\item Point the \textbf{second finger} along the \textbf{C}urrent $I$ (conventional current, $+$ to $-$).
\item Your \textbf{thumb} then points along the \textbf{M}otion --- the direction of the \textbf{force} $\vec{F}$.
\end{enumerate}
Memory aid: \textbf{F}irst = \textbf{F}ield, se\textbf{C}ond = \textbf{C}urrent, thu\textbf{M}b = \textbf{M}otion.
\end{methode}

\begin{popfigure}
\begin{center}
\begin{tikzpicture}[scale=1.1, line cap=round]
% axes
\draw[popblue, ultra thick, ->] (0,0) -- (3.0,0) node[right] {$\vec{B}$ (first finger)};
\draw[popgreen, ultra thick, ->] (0,0) -- (0,2.6) node[above] {$\vec{F}$ (thumb)};
\draw[poporange, ultra thick, ->] (0,0) -- (-1.7,-1.5) node[below left] {$I$ (second finger)};
% right-angle marks
\draw[popmuted] (0.4,0) -- (0.4,0.4) -- (0,0.4);
\draw[popmuted] (0,-0.35) -- (-0.28,-0.60);
\node[popdark] at (1.5,0.9) {$\vec{F}\perp\vec{B}$,\ \ $\vec{F}\perp I$};
\end{tikzpicture}
\end{center}
\legende{Fleming's left-hand rule: field, current and force are mutually perpendicular.}
\end{popfigure}

\begin{attention}[Do not confuse the two hand rules!]
\begin{itemize}
\item \textbf{Right-hand grip rule}: finds the \emph{field produced by a current} (source of $\vec{B}$).
\item \textbf{Fleming's left-hand rule}: finds the \emph{force on a current} placed in an \emph{external} field (effect of $\vec{B}$).
\end{itemize}
In an exam, ask yourself: ``Am I looking for the field \emph{created} by the wire (right hand) or the force \emph{suffered} by the wire (left hand)?'' Mixing them up is the most frequent source of lost marks in this chapter.
\end{attention}

\begin{exoresolu}[Force on a power line conductor]
A straight section of an aluminium conductor of length $L = 25\ \mathrm{m}$ runs horizontally in the east--west direction near Bamenda and carries a current $I = 300\ \mathrm{A}$. At this location the Earth's magnetic field has magnitude $B = 4.0\times 10^{-5}\ \mathrm{T}$ and is horizontal, pointing due north.
\begin{enumerate}
\item Calculate the magnitude of the magnetic force on the conductor.
\item Give its direction.
\item The line is rotated so that it runs north--south. What is the force now?
\end{enumerate}
\tcblower
\textbf{Solution.}
\begin{enumerate}
\item The conductor (east--west) is perpendicular to the field (north), so $\theta = 90^{\circ}$ and $\sin\theta = 1$:
\[
F = BIL\sin\theta = (4.0\times 10^{-5})(300)(25)\times 1 = 0.30\ \mathrm{N}.
\]
\item Fleming's left-hand rule: first finger north ($\vec{B}$), second finger along the current (say eastwards); the thumb points \textbf{vertically upwards}. (If the current flows westwards, the force is downwards.)
\item Now the conductor is parallel to the field: $\theta = 0^{\circ}$, $\sin\theta = 0$, hence
\[
F = 0\ \mathrm{N}.
\]
\end{enumerate}
\textbf{Comment.} The force is small because the Earth's field is weak --- this is why power lines do not visibly move, but in a $1\ \mathrm{T}$ field the same conductor would feel $7.5\ \mathrm{kN}$, the weight of a small car!
\end{exoresolu}

\courssub{Applications}

\textbf{The moving-coil galvanometer.} A rectangular coil carrying the current to be measured sits in the radial field of a permanent magnet. The forces on the two sides of the coil form a couple that rotates the coil against a hairspring; the deflection is proportional to the current. This is the heart of every analogue ammeter and voltmeter --- we shall analyse its torque in the next section.

\begin{experience}[The current balance: measuring $B$ (GCE favourite)]
A rigid rectangular wire frame is balanced on a pivot with one horizontal arm of length $L$ placed perpendicular to the field $\vec{B}$ to be measured (between magnet poles). With no current, the frame is balanced with a counterweight. A current $I$ is then passed: the force $F = BIL$ on the arm unbalances the frame, and a small mass $m$ is added to restore balance.
At equilibrium, $BIL = mg$, so
\[
B = \frac{mg}{IL}.
\]
With $m = 0.20\ \mathrm{g} = 2.0\times 10^{-4}\ \mathrm{kg}$, $I = 2.0\ \mathrm{A}$, $L = 5.0\times 10^{-2}\ \mathrm{m}$ and $g = 9.81\ \mathrm{m\,s^{-2}}$:
\[
B = \frac{(2.0\times 10^{-4})(9.81)}{(2.0)(5.0\times 10^{-2})} \approx 2.0\times 10^{-2}\ \mathrm{T}.
\]
This simple experiment is a classic GCE practical question: know the setup, the balance condition and the formula.
\end{experience}

\begin{aretenir}[Key points of this section]
\begin{itemize}
\item A magnetic field is described by the vector \textbf{magnetic flux density} $\vec{B}$, measured in \textbf{tesla}: $1\ \mathrm{T} = 1\ \mathrm{N\,A^{-1}\,m^{-1}}$.
\item Field lines point from N to S outside a magnet, never cross, and their density shows the field strength.
\item Oersted showed that \textbf{currents create magnetic fields}.
\item Straight wire: concentric circles (right-hand grip rule). Loop: N face = anticlockwise current. Solenoid: uniform field inside, bar-magnet field outside; right-hand thumb gives the north pole.
\item Force on a conductor: $F = BIL\sin\theta$; maximum when $\perp$ to $\vec{B}$, zero when parallel; direction from \textbf{Fleming's left-hand rule}.
\end{itemize}
\end{aretenir}

\begin{exercice}[Check your understanding]
\begin{enumerate}
\item A long straight wire carries a current vertically upwards. Seen from above, do the field lines circulate clockwise or anticlockwise?
\item A solenoid is connected to a d.c. supply and the current in its turns circulates clockwise as seen from its left end. Which end is the north pole?
\item A conductor of length $0.40\ \mathrm{m}$ carrying $5.0\ \mathrm{A}$ makes an angle of $30^{\circ}$ with a uniform field of $0.25\ \mathrm{T}$. Calculate the force on it, and state the angle for which this force would be maximum.
\end{enumerate}
\end{exercice}

\begin{corrige}[Exercise 1]
\begin{enumerate}
\item Right-hand grip rule, thumb upwards: seen from above the fingers curl \textbf{anticlockwise}.
\item Clockwise current as seen from the left end $\Rightarrow$ the \textbf{left end is a south pole}; the \textbf{right end is the north pole} (right-hand grip rule: thumb points to the right).
\item $F = BIL\sin\theta = (0.25)(5.0)(0.40)\sin 30^{\circ} = 0.50\times 0.5 = 0.25\ \mathrm{N}$. The force is maximum for $\theta = 90^{\circ}$ (conductor perpendicular to $\vec{B}$), giving $F_{\max} = 0.50\ \mathrm{N}$.
\end{enumerate}
\end{corrige}

\coursec{Torque on a Coil and the Electric Motor}

In the previous section we established that a straight conductor carrying a current $I$ in a uniform magnetic field of flux density $B$ experiences a force $F = BIl\sin\theta$, whose direction is given by Fleming's left-hand rule. A single straight wire, however, can only be pushed sideways. If we bend the wire into a rectangular loop, something far more interesting happens: the forces on opposite sides combine into a \cle{couple}, and the loop \emph{rotates}. This is the physical principle behind the electric motor, the moving-coil galvanometer, and countless devices that convert electrical energy into mechanical work --- from the ceiling fans of Douala to the water pumps of the Far North.

\courssub{A Rectangular Coil in a Uniform Magnetic Field}

\textbf{Setting up the situation.}\par
Consider a rectangular coil $PQRS$ of $N$ turns, with sides of length $a$ (the ``vertical'' sides $PQ$ and $RS$) and width $b$ (the ``horizontal'' sides $QR$ and $SP$). The coil carries a current $I$ and is placed in a uniform magnetic field $\vec{B}$ directed horizontally. The coil is free to rotate about a vertical axis passing through the centres of sides $QR$ and $SP$. We begin with the plane of the coil \emph{parallel} to $\vec{B}$.

\begin{popfigure}
\begin{tikzpicture}[scale=1.0,>=stealth]
% Magnetic field lines
\foreach \y in {-0.9,-0.3,0.3,0.9} {
  \draw[->,popblue,thick] (-3.4,\y) -- (3.4,\y);
}
\node[popblue] at (-3.4,1.2) {$\vec{B}$};
% Coil in perspective
\draw[very thick,poporange] (-1.1,-1.1) -- (1.1,-1.1) -- (1.7,1.1) -- (-1.7,1.1) -- cycle;
% Axis of rotation
\draw[dashed,popdark,thick] (0,-1.7) -- (0,1.7);
\node[popdark,align=center,text width=2.2cm] at (0.1,2.0) {axis of rotation};
% Labels of sides
\node[poporange] at (-1.55,0) {$PQ$};
\node[poporange] at (1.55,0) {$RS$};
\node[poporange] at (0,-1.3) {$QR$};
\node[poporange] at (0,1.35) {$SP$};
% Current directions
\draw[->,popgreen,very thick] (-1.45,-0.6) -- (-1.45,0.5);
\draw[->,popgreen,very thick] (1.45,0.6) -- (1.45,-0.5);
\node[popgreen] at (-1.9,0.6) {$I$};
\node[popgreen] at (1.9,-0.6) {$I$};
% Forces
\draw[->,popred,very thick] (-1.45,0.2) -- (-2.6,0.2);
\node[popred] at (-2.9,0.2) {$\vec{F}$};
\draw[->,popred,very thick] (1.45,-0.2) -- (2.6,-0.2);
\node[popred] at (2.9,-0.2) {$\vec{F}$};
\end{tikzpicture}
\legende{A rectangular coil carrying a current $I$ in a uniform field $\vec{B}$. The forces on sides $PQ$ and $RS$ form a couple that rotates the coil about the dashed axis.}
\end{popfigure}

\textbf{Forces on each of the four sides.}\par
Apply $F = BIl\sin\theta$ and Fleming's left-hand rule to each side, with the plane of the coil parallel to $\vec{B}$:

\begin{itemize}
\item \textbf{Side $PQ$ (length $a$):} the current flows perpendicular to $\vec{B}$, so $\theta = 90^{\circ}$ and the force has magnitude $F = BIa$. Fleming's left-hand rule shows it acts \emph{out of the plane} of the coil (say, towards the reader).
\item \textbf{Side $RS$ (length $a$):} the current here flows in the opposite direction to that in $PQ$ (the current goes round the loop), still perpendicular to $\vec{B}$. The force again has magnitude $F = BIa$ but acts in the \emph{opposite} direction (away from the reader).
\item \textbf{Sides $QR$ and $SP$ (length $b$):} here the current flows \emph{parallel or antiparallel} to $\vec{B}$, so $\sin\theta = 0$ and these sides experience \textbf{no force at all}.
\end{itemize}

\begin{attention}[Which sides feel a force?]
Only the sides whose current has a component \emph{perpendicular} to $\vec{B}$ experience a force. When the plane of the coil is parallel to $\vec{B}$, the two ``vertical'' sides carry current perpendicular to the field and feel the full force $F = BIa$; the two ``horizontal'' sides carry current parallel to the field and feel nothing. Do not assume all four sides are pushed!
\end{attention}

\courssub{The Couple and the Torque: Derivation of $T = NAIB$}

The two forces on $PQ$ and $RS$ are equal in magnitude, opposite in direction, and act along \emph{different lines of action}. They therefore constitute a \cle{couple}, whose turning effect is measured by its \cle{torque} (moment of the couple).

\begin{definition}[Torque of a couple]
The torque $T$ of a couple formed by two equal and opposite forces of magnitude $F$ whose lines of action are separated by a perpendicular distance $d$ (the \cle{lever arm}) is
\[ T = F \times d. \]
Its SI unit is the newton metre, $\mathrm{N\,m}$. The torque tends to rotate the body about the axis midway between the two forces.
\end{definition}

\begin{popfigure}
\begin{tikzpicture}[scale=1.0,>=stealth]
% Field
\foreach \y in {-1.3,-0.65,0,0.65,1.3} {
  \draw[->,popblue] (-3.2,\y) -- (3.2,\y);
}
\node[popblue] at (3.2,1.6) {$\vec{B}$};
% Coil seen edge-on: vertical sides as cross/dot
\draw[very thick,poporange] (-1.4,-1.5) -- (-1.4,1.5);
\draw[very thick,poporange] (1.4,-1.5) -- (1.4,1.5);
\draw[very thick,poporange] (-1.4,1.5) -- (1.4,1.5);
\draw[very thick,poporange] (-1.4,-1.5) -- (1.4,-1.5);
% Axis
\draw[dashed,popdark,thick] (0,-1.9) -- (0,1.9);
\fill[popdark] (0,0) circle (0.06);
\node[popdark] at (0.35,-1.75) {axis};
% Current symbols
\node[popgreen] at (-1.4,0) {\Large $\odot$};
\node[popgreen] at (1.4,0) {\Large $\otimes$};
\node[popgreen] at (-2.0,0.35) {$I$ out};
\node[popgreen] at (2.0,0.35) {$I$ in};
% Forces
\draw[->,popred,very thick] (-1.4,0.9) -- (-1.4,1.9);
\node[popred] at (-1.4,2.1) {$\vec{F}$};
\draw[->,popred,very thick] (1.4,-0.9) -- (1.4,-1.9);
\node[popred] at (1.4,-2.1) {$\vec{F}$};
% Lever arm
\draw[<->,poppurple,thick] (-1.4,-0.55) -- (1.4,-0.55);
\node[poppurple] at (0,-0.8) {lever arm $d = b$};
% side labels
\node[popmuted] at (0,1.75) {width $b$};
\end{tikzpicture}
\legende{Edge-on view of the coil. The current leaves the page on the left side ($\odot$) and enters on the right ($\otimes$). The two forces $F = BIa$, separated by the lever arm $b$, form a couple of torque $T = Fb$.}
\end{popfigure}

\textbf{Derivation (plane of coil parallel to $\vec{B}$).}\par
For one turn of the coil:
\begin{itemize}
\item Force on each vertical side: $F = BIa$ (since the current is perpendicular to $\vec{B}$).
\item The lines of action of the two forces are separated by the width of the coil, so the lever arm is $d = b$.
\item Torque of the couple: $T_1 = F \times b = BIa \times b = BI(ab)$.
\end{itemize}
But $ab = A$ is the \emph{area} of the coil. For $N$ identical turns, each carrying the same current $I$, the torques add:
\[ T = N \times BIA = NAIB. \]

\begin{propriete}[Torque on a rectangular coil --- examinable derivation]
A plane rectangular coil of $N$ turns, area $A$, carrying a current $I$, placed in a uniform magnetic field of flux density $B$ with its plane parallel to the field, experiences a torque
\[ \boxed{T = NAIB} \]
\begin{center}
\begin{tabularx}{\linewidth}{|l|X|l|}
\hline
\rowcolor{popblueL} Symbol & Meaning & SI unit \\
\hline
$T$ & torque (moment of the couple) on the coil & $\mathrm{N\,m}$ \\
\hline
$N$ & number of turns of the coil & (dimensionless) \\
\hline
$A$ & area enclosed by one turn, $A = ab$ & $\mathrm{m^2}$ \\
\hline
$I$ & current in the coil & $\mathrm{A}$ \\
\hline
$B$ & magnetic flux density & $\mathrm{T}$ \\
\hline
\end{tabularx}
\end{center}
The derivation above (force on each side $\rightarrow$ couple $\rightarrow$ $T = NAIB$) is a classic GCE question and must be reproducible with a labelled diagram.
\end{propriete}

\courssub{Coil at an Angle: $T = NAIB\cos\theta$ and Equilibrium Positions}

Suppose the coil has rotated so that its \emph{plane} makes an angle $\theta$ with the field $\vec{B}$. The forces on the vertical sides still have magnitude $F = BIa$ (the current is still perpendicular to $\vec{B}$), but the \emph{perpendicular distance} between their lines of action is now reduced to $d = b\cos\theta$. Hence
\[ T = F \times b\cos\theta = NAIB\cos\theta. \]

\begin{attention}[The angle convention trap]
Here $\theta$ is the angle between the \textbf{plane of the coil} and $\vec{B}$. Some textbooks define $\alpha$ as the angle between the \textbf{normal to the coil} and $\vec{B}$; then $\alpha = 90^{\circ} - \theta$ and the same physics is written $T = NAIB\sin\alpha$. Read the question carefully! The safest statement to remember physically: \textbf{the torque is maximum when the plane of the coil is parallel to $\vec{B}$, and zero when the plane of the coil is perpendicular to $\vec{B}$} (normal parallel to $\vec{B}$).
\end{attention}

\begin{popfigure}
\begin{tikzpicture}
\begin{axis}[
  width=11cm, height=6.5cm,
  xlabel={$\theta$ (angle between plane of coil and $\vec{B}$) / degrees},
  ylabel={$T / (NAIB)$},
  xmin=0, xmax=360, ymin=-1.25, ymax=1.25,
  xtick={0,90,180,270,360},
  ytick={-1,-0.5,0,0.5,1},
  grid=both, grid style={popline!40},
  axis lines=middle,
  domain=0:360, samples=200,
]
\addplot[very thick,poporange] {cos(x)};
\addplot[mark=*,mark size=1.5pt,popdark,only marks] coordinates {(0,1) (90,0) (180,-1) (270,0) (360,1)};
\node[popdark,align=center,text width=3.2cm] at (axis cs:45,1.13) {max: plane $\parallel \vec{B}$};
\node[popdark,align=center,text width=3.4cm] at (axis cs:90,-0.35) {zero: plane $\perp \vec{B}$};
\node[popdark,align=center,text width=3.2cm] at (axis cs:225,-1.12) {max, reversed sense};
\end{axis}
\end{tikzpicture}
\legende{Torque on the coil as a function of the angle $\theta$ between the plane of the coil and the field: $T = NAIB\cos\theta$.}
\end{popfigure}

\textbf{Equilibrium positions.}\par
The torque vanishes at $\theta = 90^{\circ}$ and $\theta = 270^{\circ}$, i.e.\ when the plane of the coil is perpendicular to $\vec{B}$:
\begin{itemize}
\item \textbf{Stable equilibrium} ($\theta = 90^{\circ}$): if the coil is displaced slightly, the torque acts to bring it back. The magnetic flux through the coil is maximum here.
\item \textbf{Unstable equilibrium} ($\theta = 270^{\circ}$): the torque is also zero, but the slightest displacement produces a torque that turns the coil \emph{away}, through half a revolution, towards the stable position.
\end{itemize}

\textbf{Reversing the current.}\par
The force on each side is proportional to the current and its direction is given by Fleming's left-hand rule. Reversing $I$ reverses \emph{every} force, hence reverses the sense of the torque: the coil rotates the other way. In the formula, $T = NAIB\cos\theta$ changes sign. This is exactly why a direct current would make a simple coil oscillate rather than spin continuously --- a difficulty the commutator will solve.

\begin{methode}[Finding the forces on each side and the sense of rotation]
\begin{enumerate}
\item Draw the coil and mark the direction of $\vec{B}$ and of the current $I$ around the loop.
\item For each side, check the angle between the current and $\vec{B}$: if they are parallel, the force is zero; if perpendicular, $F = BIl$.
\item Apply Fleming's left-hand rule to each current-carrying side that is perpendicular to $\vec{B}$: first finger = Field, second finger = Current, thumb = Motion (force).
\item Identify the pair of equal, opposite, non-collinear forces: they form the couple.
\item Look along the axis of rotation and state whether the coil turns clockwise or anticlockwise.
\end{enumerate}
\end{methode}

\begin{exoresolu}[Torque on a galvanometer coil]
A rectangular coil has $N = 200$ turns, each of area $A = 2.0 \times 10^{-4}\ \mathrm{m^2}$. It carries a current $I = 5.0\ \mathrm{mA}$ in a uniform field $B = 0.15\ \mathrm{T}$. (a) Calculate the maximum torque on the coil. (b) Calculate the torque when the plane of the coil makes $30^{\circ}$ with the field.
\tcblower
\textbf{(a)} The torque is maximum when the plane of the coil is parallel to $\vec{B}$ ($\theta = 0^{\circ}$, $\cos\theta = 1$):
\[ T_{\max} = NAIB = 200 \times (2.0\times10^{-4}) \times (5.0\times10^{-3}) \times 0.15 \]
\[ T_{\max} = 3.0\times10^{-5}\ \mathrm{N\,m}. \]
\textbf{(b)} With $\theta = 30^{\circ}$ between the plane and the field:
\[ T = NAIB\cos\theta = 3.0\times10^{-5}\times\cos 30^{\circ} = 3.0\times10^{-5}\times 0.866 \approx 2.6\times10^{-5}\ \mathrm{N\,m}. \]
Note the units: $B$ in tesla, $A$ in $\mathrm{m^2}$, $I$ in amp\`eres give $T$ in $\mathrm{N\,m}$, since $1\ \mathrm{T} = 1\ \mathrm{N\,A^{-1}\,m^{-1}}$.
\end{exoresolu}

\courssub{The Moving-Coil Galvanometer Revisited}

In the chapter on measuring instruments you met the moving-coil galvanometer. We can now understand \emph{why} it works and why its scale is linear.

The coil, of $N$ turns and area $A$, sits in the field of a permanent magnet whose pole pieces are shaped (with a soft-iron cylinder at the centre) so that the field is \cle{radial}: whatever the position of the coil, its plane is always parallel to the field lines it cuts. The magnetic torque is therefore always
\[ T_{\text{mag}} = NAIB, \]
independent of the angle of rotation. The coil is controlled by a \cle{hairspring} which exerts a restoring torque proportional to the angle of deflection:
\[ T_{\text{spring}} = k\phi, \]
where $k$ is the torsion constant of the spring and $\phi$ the angular deflection. The pointer stops when the two torques balance:
\[ NAIB = k\phi \quad\Longrightarrow\quad \phi = \frac{NAB}{k}\,I. \]
The deflection $\phi$ is \emph{directly proportional} to the current $I$: the scale is \textbf{linear}. This proportionality is also why the same movement, fitted with a shunt or a multiplier, becomes an ammeter or a voltmeter with a uniform scale.

\begin{aretenir}[Why the radial field matters]
In an ordinary uniform field the torque would vary as $\cos\theta$ and the scale would be cramped at the ends. The radial field keeps the plane of the coil parallel to $\vec{B}$ at all times, so $T_{\text{mag}} = NAIB$ throughout the rotation and $\phi \propto I$: a perfectly linear scale.
\end{aretenir}

\courssub{The Direct-Current (dc) Motor}

\begin{popfigure}
\begin{tikzpicture}[scale=1.0,>=stealth]
% Magnet poles
\fill[popblueL] (-3.4,-1.6) rectangle (-2.2,1.6);
\fill[popblueL] (2.2,-1.6) rectangle (3.4,1.6);
\node[popdark] at (-2.8,1.9) {\textbf{N}};
\node[popdark] at (2.8,1.9) {\textbf{S}};
% Field lines
\foreach \y in {-1.0,-0.33,0.33,1.0} {
  \draw[->,popblue] (-2.15,\y) -- (2.15,\y);
}
% Coil (edge on)
\draw[very thick,poporange] (-0.9,-1.1) rectangle (0.9,1.1);
\node[popgreen] at (-0.9,0) {\Large $\odot$};
\node[popgreen] at (0.9,0) {\Large $\otimes$};
\draw[->,popred,very thick] (-0.9,0.55) -- (-0.9,1.55);
\draw[->,popred,very thick] (0.9,-0.55) -- (0.9,-1.55);
\node[popred] at (-0.9,1.8) {$\vec{F}$};
\node[popred] at (0.9,-1.8) {$\vec{F}$};
% Axle
\fill[popdark] (0,0) circle (0.09);
% Commutator
\draw[very thick,popgold] (0,-2.6) circle (0.45);
\fill[white] (0.12,-3.05) rectangle (-0.12,-2.15);
\draw[very thick,popgold] (0.45,-2.6) arc (0:180:0.45);
\draw[very thick,popgold] (-0.45,-2.6) arc (180:360:0.45);
\node[popdark,align=center,text width=3.4cm] at (2.6,-2.6) {split-ring commutator (two half-rings)};
% Brushes
\fill[popdark] (-0.75,-2.95) rectangle (-0.45,-2.6);
\fill[popdark] (0.45,-2.95) rectangle (0.75,-2.6);
\node[popdark,align=center,text width=2.4cm] at (-2.3,-3.1) {carbon brushes};
% Battery
\draw[thick] (-0.6,-4.2) -- (0.6,-4.2);
\draw[thick] (-0.25,-4.35) -- (0.25,-4.35);
\draw (-0.6,-2.95) -- (-0.6,-4.2);
\draw (0.6,-2.95) -- (0.6,-4.35);
\node[popdark] at (1.5,-4.3) {dc supply};
% wires coil to commutator
\draw[popmuted] (-0.9,-1.1) -- (-0.35,-2.25);
\draw[popmuted] (0.9,-1.1) -- (0.35,-2.25);
\end{tikzpicture}
\legende{Essential parts of a simple dc motor: permanent-magnet poles, current-carrying coil, split-ring commutator rotating with the coil, and fixed carbon brushes connected to the dc supply.}
\end{popfigure}

\textbf{Construction.}\par
A coil of many turns is wound on a soft-iron core (the \emph{armature}) mounted on an axle between the poles of a magnet. Current reaches the rotating coil through two fixed carbon \cle{brushes} which press against a \cle{split-ring commutator}: a copper ring cut into two insulated halves, each half connected to one end of the coil and rotating with it.

\textbf{Operation.}\par
With the plane of the coil parallel to $\vec{B}$, the forces on the two active sides form a couple $T = NAIB$ and the coil starts to rotate. As the coil passes through the perpendicular position (where the torque would vanish), its momentum carries it past --- and at that very instant each half-ring swaps brushes, so the current in the coil \emph{reverses}. The forces on the two sides therefore reverse too, and the torque keeps acting in the \emph{same sense of rotation}. The current reverses every half-turn, and the coil spins continuously.

\begin{propriete}[Role of the split-ring commutator]
The split-ring commutator reverses the current in the coil every half-revolution, exactly when the coil passes through the position where its plane is perpendicular to the field. Consequently the torque, although pulsating in magnitude, always acts in the \emph{same direction}, and the coil rotates continuously instead of oscillating about the stable equilibrium position. (In real motors, many coils and a multi-segment commutator smooth the torque.)
\end{propriete}

\begin{attention}[Motor or generator?]
A \textbf{motor} converts \emph{electrical} energy into \emph{mechanical} energy: you supply current and the shaft turns. A \textbf{generator} does the opposite: you turn the shaft mechanically and an e.m.f.\ (current) is produced by electromagnetic induction. The two machines have essentially the same construction --- the same ``petit g\'en\'erateur'' engine-alternator used in many Cameroonian neighbourhoods contains a coil rotating in a magnetic field --- but the direction of energy conversion is opposite. Never write that a motor ``produces electricity''.
\end{attention}

\courssub{The Alternating-Current (ac) Motor}

In a simple ac motor, no commutator is needed. In the \emph{induction} motor --- the commonest type --- sets of fixed coils (the \emph{stator}) are fed with alternating current and arranged so that their combined magnetic field \emph{rotates} in space. This rotating field induces currents in the rotor conductors, and the forces on these induced currents drag the rotor round, following the field. Because the alternating supply itself provides the periodic reversal, there are no brushes and no split ring, which makes ac motors rugged, cheap and almost maintenance-free.

\begin{center}
\begin{tabularx}{\linewidth}{|l|X|X|}
\hline
\rowcolor{popblueL} Feature & dc motor & ac (induction) motor \\
\hline
Supply & direct current (battery, rectifier) & alternating current (mains, generator) \\
\hline
Current reversal & split-ring commutator + brushes & provided by the ac supply itself \\
\hline
Wear parts & brushes wear out, need replacement & no brushes: very robust \\
\hline
Speed control & easy, by varying the voltage & less direct (frequency control) \\
\hline
Typical uses & small appliances, toys, car starters & fans, pumps, industrial machines \\
\hline
\end{tabularx}
\end{center}

\begin{savaistu}[Motors around us in Cameroon]
The ceiling and standing fans that make life bearable in Douala and Yaound\'e, the electric water pumps that lift water from boreholes in the Far North, and the alternators of the ``petits g\'en\'erateurs'' used during power cuts all rely on the same physics: forces on current-carrying coils in magnetic fields. Motor and generator are dual machines --- feed electricity in and you get rotation out; force rotation in and you get electricity out.
\end{savaistu}

\begin{exercice}[Torque and motor action]
A coil of $150$ turns, of dimensions $4.0\ \mathrm{cm} \times 2.5\ \mathrm{cm}$, carries a current of $2.0\ \mathrm{A}$ in a uniform field of flux density $0.30\ \mathrm{T}$.
\begin{enumerate}
\item Calculate the maximum torque on the coil and state the orientation of the coil at which it occurs.
\item Calculate the torque when the normal to the plane of the coil makes $60^{\circ}$ with $\vec{B}$.
\item Explain, with reference to the commutator, why the coil of a dc motor does not stop when its plane becomes perpendicular to the field.
\end{enumerate}
\end{exercice}

\begin{corrige}[Exercise 1]
\textbf{1.} $A = 0.040 \times 0.025 = 1.0\times10^{-3}\ \mathrm{m^2}$. Maximum torque (plane of coil parallel to $\vec{B}$):
\[ T_{\max} = NAIB = 150 \times 1.0\times10^{-3} \times 2.0 \times 0.30 = 9.0\times10^{-2}\ \mathrm{N\,m}. \]
\textbf{2.} The normal making $60^{\circ}$ with $\vec{B}$ means the plane makes $\theta = 30^{\circ}$ with $\vec{B}$:
\[ T = T_{\max}\cos 30^{\circ} = 9.0\times10^{-2}\times 0.866 \approx 7.8\times10^{-2}\ \mathrm{N\,m}. \]
(Equivalently $T = NAIB\sin 60^{\circ}$, same result.)
\textbf{3.} At the perpendicular position the torque is zero, but the coil's rotational kinetic energy carries it past. As it crosses this position, the split-ring commutator swaps the connections of the coil ends to the supply, reversing the current. The forces on the active sides reverse accordingly, so the torque again acts in the original sense of rotation, and continuous rotation is maintained.
\end{corrige}

\begin{aretenir}[Key points of the section]
\begin{itemize}
\item Only the sides of the coil whose current is perpendicular to $\vec{B}$ experience forces; these two equal and opposite forces form a couple.
\item Torque on the coil: $T = NAIB\cos\theta$, with $\theta$ the angle between the \emph{plane} of the coil and $\vec{B}$; maximum when plane $\parallel \vec{B}$, zero when plane $\perp \vec{B}$.
\item Reversing the current reverses the torque.
\item Moving-coil galvanometer: radial field $\Rightarrow T_{\text{mag}} = NAIB$; balance with the hairspring $k\phi$ gives $\phi \propto I$ (linear scale).
\item dc motor: the split-ring commutator reverses the current every half-turn so the torque keeps the same sense.
\item ac motor: the alternating supply (or a rotating field) provides the reversal; no commutator.
\item Motor: electrical $\rightarrow$ mechanical; generator: mechanical $\rightarrow$ electrical.
\end{itemize}
\end{aretenir}

\coursec{Sources of Magnetic Fields: Biot-Savart, Ampère and Forces Between Conductors}

In the previous section we saw that a current-carrying coil placed in a magnetic field experiences a torque, and we exploited this to build the electric motor. But a natural question arises: \emph{where does the magnetic field itself come from?} Oersted's famous observation — a compass needle deflecting near a current-carrying wire — tells us that \cle{electric currents produce magnetic fields}. In this section we quantify this statement. Two fundamental laws, the \cle{Biot-Savart law} and \cle{Ampère's circuital law}, allow us to calculate the magnetic flux density produced by conductors of standard shapes: the long straight wire, the plane circular coil and the long solenoid. We finish by combining these results with the motor effect to derive the force between two parallel conductors — the very phenomenon that historically defined the ampere.

\courssub{The Biot-Savart Law}

\textbf{From observation to a law.} Experiments show that the magnetic field produced by a current decreases as we move away from the wire, and increases with the current. To make this precise, we consider a tiny piece of conductor, a \cle{current element} $I\,\mathrm{d}\vec{l}$, and ask what contribution $\mathrm{d}\vec{B}$ it makes to the field at a point $P$ located at a distance $r$ from it.

\begin{propriete}[Biot-Savart law]
The contribution $\mathrm{d}B$ of a current element $I\,\mathrm{d}l$ to the magnetic flux density at a point $P$ situated at distance $r$, where $\theta$ is the angle between the element $\mathrm{d}\vec{l}$ and the line joining the element to $P$, is
\[ \mathrm{d}B = \frac{\mu_0}{4\pi}\,\frac{I\,\mathrm{d}l\,\sin\theta}{r^{2}} \]
where $\mu_0 = 4\pi\times 10^{-7}\ \mathrm{H\,m^{-1}}$ (equivalently $\mathrm{T\,m\,A^{-1}}$) is the \cle{permeability of free space}. The direction of $\mathrm{d}\vec{B}$ is given by the right-hand grip rule: it is perpendicular to the plane containing $\mathrm{d}\vec{l}$ and $\vec{r}$, in the sense of $\mathrm{d}\vec{l}\times\vec{r}$. The total field is obtained by summing (integrating) the contributions of all elements of the circuit. \emph{(The statement is examinable; the integrations below should be understood and reproducible.)}
\end{propriete}

Note the structure of the law: it is an \emph{inverse-square} law in $r$ (like Coulomb's law), it is proportional to $I$, and the factor $\sin\theta$ means that a point on the axis of the element ($\theta = 0$) receives \emph{no} contribution from that element.

\begin{popfigure}
\begin{tikzpicture}[scale=1.1, >=stealth]
\draw[very thick, popblue] (0,-1.6) -- (0,1.6);
\draw[->, very thick, popblue] (0,0.55) -- (0,1.15);
\node[popblue, left] at (0,1.1) {$I$};
\draw[->, ultra thick, poporange] (0,-0.35) -- (0,0.35);
\node[poporange, left] at (0,0) {$I\,\mathrm{d}\vec{l}$};
\fill[popdark] (3.4,1.0) circle (2pt);
\node[popdark, above right] at (3.4,1.0) {$P$};
\draw[->, thick, popdark] (0,0) -- (3.4,1.0);
\node[popdark, below, rotate=16] at (1.7,0.5) {$\vec{r}$};
\draw[thick, popgreen] (0.85,0) arc (0:16:0.85);
\node[popgreen] at (1.15,0.18) {$\theta$};
\draw[dashed, gray] (0,0) -- (3.4,0);
\draw[->, thick, poppurple] (3.4,1.0) -- (3.4,0.35);
\node[poppurple, right] at (3.4,0.6) {$\mathrm{d}\vec{B}$ (into page)};
\end{tikzpicture}
\legende{Biot-Savart geometry: the current element $I\,\mathrm{d}\vec{l}$, the position vector $\vec{r}$ to the point $P$, and the angle $\theta$ between them.}
\end{popfigure}

\courssub{Field at the Centre of a Plane Circular Coil}

Consider a circular loop of radius $r$ carrying a current $I$. Every current element $I\,\mathrm{d}l$ is at the same distance $r$ from the centre $O$, and the radius is always perpendicular to the element, so $\theta = 90^{\circ}$ and $\sin\theta = 1$ everywhere. Moreover, by the right-hand rule, all the contributions $\mathrm{d}\vec{B}$ point in the \emph{same} direction along the axis of the coil. The integration is therefore a simple sum:
\[ B = \oint \mathrm{d}B = \frac{\mu_0}{4\pi}\,\frac{I}{r^{2}}\oint \mathrm{d}l = \frac{\mu_0}{4\pi}\,\frac{I}{r^{2}}\,(2\pi r) = \frac{\mu_0 I}{2r}. \]
For a flat coil of $N$ identical turns wound together, each turn contributes the same field, so:

\begin{propriete}[Field at the centre of a plane circular coil]
\[ B = \frac{\mu_0 N I}{2r} \]
directed along the axis of the coil, with the sense given by the right-hand grip rule (fingers curled with the current, thumb gives $\vec{B}$).
\end{propriete}

\begin{exemplebox}[Laboratory coil]
A student in a physics laboratory in Bamenda winds a flat coil of $N = 50$ turns and mean radius $r = \SI{10}{cm}$, and passes a current $I = \SI{2.0}{A}$. The field at the centre is
\[ B = \frac{(4\pi\times 10^{-7})(50)(2.0)}{2(0.10)} = 6.3\times 10^{-4}\ \mathrm{T} = \SI{0.63}{mT}, \]
about ten times the Earth's magnetic field — easily detected with a compass needle.
\end{exemplebox}

\courssub{Field due to a Long Straight Wire}

For a very long straight wire carrying a current $I$, summing the Biot-Savart contributions of all elements along the wire (an integration over the whole wire, which requires trigonometric substitution and is \emph{not} examinable) gives a beautifully simple result:

\begin{propriete}[Field of a long straight wire]
At a perpendicular distance $r$ from a long straight conductor,
\[ B = \frac{\mu_0 I}{2\pi r}. \]
The field lines are concentric circles centred on the wire, with the sense given by the right-hand grip rule (thumb along the current, fingers curl with $\vec{B}$).
\end{propriete}

The field decays as $1/r$: twice as far from the wire, the field is half as strong. This hyperbolic decay is illustrated below.

\begin{popfigure}
\begin{tikzpicture}
\begin{axis}[
  width=0.85\linewidth, height=5.5cm,
  xlabel={$r$ (cm)}, ylabel={$B$ ($\mu$T)},
  xmin=0, xmax=20, ymin=0, ymax=220,
  xtick={0,5,10,15,20}, ytick={0,50,100,150,200},
  grid=both, grid style={popline, very thin},
  axis line style={popdark}, label style={popdark},
  samples=100, domain=1:20]
\addplot[very thick, popblue] {200/x};
\addplot[only marks, mark=*, poporange] coordinates {(1,200) (2,100) (4,50) (5,40) (10,20)};
\end{axis}
\end{tikzpicture}
\legende{Magnetic flux density $B = \mu_0 I/2\pi r$ versus distance $r$ from a long straight wire carrying $I = \SI{10}{A}$ (here $B = 200/r$ with $B$ in $\mu$T and $r$ in cm). Note the $1/r$ decay.}
\end{popfigure}

\courssub{Ampère's Circuital Law}

The Biot-Savart law always works, but the integration can be heavy. When the current distribution is \emph{highly symmetric}, a second law gives the answer in a few lines.

\begin{propriete}[Ampère's circuital law]
For any closed loop $\mathcal{L}$ (called an \cle{Amperian loop}),
\[ \oint_{\mathcal{L}} \vec{B}\cdot\mathrm{d}\vec{l} = \mu_0\, I_{\text{enclosed}}, \]
where $I_{\text{enclosed}}$ is the total current passing through the surface bounded by the loop, counted with a sign given by the right-hand rule relative to the direction of travel around the loop. \emph{(Statement examinable.)}
\end{propriete}

Ampère's law is always true, but it is only \emph{useful for calculation} when symmetry lets us choose a loop on which $B$ is constant and parallel (or perpendicular) to $\mathrm{d}\vec{l}$, so that the integral reduces to $B\times(\text{length})$.

\begin{attention}[When does Ampère's law help?]
Use Ampère's law only when the symmetry is high enough: infinite straight wire, infinite solenoid, toroid, infinite plane of current. For a finite wire or an arbitrary loop, the integral cannot be simplified and you must return to Biot-Savart.
\end{attention}

\textbf{Check: the straight wire revisited.} Choose a circular Amperian loop of radius $r$ centred on the wire. By symmetry $B$ has the same magnitude everywhere on the circle and is tangential to it, so
\[ \oint \vec{B}\cdot\mathrm{d}\vec{l} = B\,(2\pi r) = \mu_0 I \quad\Longrightarrow\quad B = \frac{\mu_0 I}{2\pi r}, \]
recovering the Biot-Savart result in two lines.

\courssub{Field Inside a Long Solenoid}

A \cle{solenoid} is a long helix of wire. For an \emph{ideal} (infinitely long, tightly wound) solenoid with $n$ turns per unit length carrying a current $I$, symmetry arguments show that the field inside is \cle{uniform} and parallel to the axis, while the field outside is essentially zero.

Choose a rectangular Amperian loop $abcd$ of side $\ell$, with the side $ab$ inside the solenoid parallel to the axis and the side $cd$ outside. Along $cd$, $B \approx 0$; along $bc$ and $da$, $\vec{B}\perp\mathrm{d}\vec{l}$, so these contribute nothing. Hence
\[ \oint \vec{B}\cdot\mathrm{d}\vec{l} = B\ell = \mu_0\,(n\ell) I, \]
since the loop encloses $n\ell$ turns each carrying $I$. Therefore:

\begin{propriete}[Field inside a long solenoid]
\[ B = \mu_0 n I \]
where $n = N/L$ is the number of turns per unit length ($\mathrm{m^{-1}}$). The field is uniform throughout the interior and directed along the axis (right-hand grip rule).
\end{propriete}

\begin{attention}[$n$ is not $N$!]
A very common examination error: $n$ is the number of turns \emph{per metre}, $n = N/L$, not the total number of turns $N$. A solenoid of $N = 2000$ turns and length $L = \SI{40}{cm}$ has $n = 2000/0.40 = 5000\ \mathrm{turns\,m^{-1}}$. Substituting $N$ instead of $n$ in $B = \mu_0 n I$ gives an answer 400 times too small here — and costs marks.
\end{attention}

\begin{attention}[A surprising uniformity]
For the ideal long solenoid, $B = \mu_0 n I$ contains \emph{no} distance variable: the field is the same on the axis, near the winding, and everywhere between. Do not invent a $1/r$ dependence inside a solenoid — that behaviour belongs to the straight wire. (Near the ends of a real solenoid the field does fall to about half its central value.)
\end{attention}

\begin{popfigure}
\begin{tikzpicture}[scale=1.0, >=stealth]
% Straight wire with Amperian loop
\draw[very thick, popblue] (0,-1.7) -- (0,1.7);
\draw[->, very thick, popblue] (0,0.9) -- (0,1.5);
\node[popblue, left] at (0,1.4) {$I$};
\draw[thick, poporange, ->] (1.1,0) arc (0:340:1.1);
\node[poporange] at (1.55,0.55) {$\mathcal{L}$};
\draw[<->, popdark] (0,-0.75) -- (1.1,-0.75);
\node[popdark, below] at (0.55,-0.75) {$r$};
\node[popdark, align=center] at (0,-2.3) {circular loop\\ around a wire};
% Solenoid with rectangular loop
\begin{scope}[xshift=6.2cm]
\foreach \x in {0,0.45,0.9,1.35,1.8,2.25,2.7,3.15}{
  \draw[thick, popblue] (\x,0.55) ellipse (0.14 and 0.55);}
\draw[thick, popblue] (0,1.1) -- (3.15,1.1);
\draw[thick, popblue] (0,0) -- (3.15,0);
\foreach \x in {0.4,1.2,2.0,2.8}{
  \draw[->, thick, popgreen] (\x,0.55) -- (\x+0.55,0.55);}
\node[popgreen, above] at (1.6,0.62) {$\vec{B}$};
\draw[thick, poporange] (0.5,-0.5) rectangle (2.7,0.55);
\draw[->, thick, poporange] (0.5,-0.5) -- (1.6,-0.5);
\node[poporange, below] at (1.6,-0.5) {$\mathcal{L}$};
\node[poporange, left] at (0.5,0.0) {$\ell$};
\node[popdark, align=center] at (1.6,-1.15) {rectangular loop\\ through a solenoid};
\end{scope}
\end{tikzpicture}
\legende{Amperian loops: a circle of radius $r$ around a straight wire, and a rectangle straddling the side of a long solenoid (one side inside, one outside).}
\end{popfigure}

\begin{propriete}[Summary of the three standard results]
\begin{center}
\begin{tabularx}{\linewidth}{|l|X|X|}
\hline
\rowcolor{popblueL} \textbf{Source} & \textbf{Magnetic flux density} & \textbf{Key feature} \\
\hline
Long straight wire & $B = \dfrac{\mu_0 I}{2\pi r}$ & Decays as $1/r$; circular field lines \\
\hline
Centre of plane coil ($N$ turns) & $B = \dfrac{\mu_0 N I}{2r}$ & Uniform only near the centre \\
\hline
Long solenoid ($n$ turns/m) & $B = \mu_0 n I$ & Uniform throughout the interior \\
\hline
\end{tabularx}
\end{center}
When several conductors are present, the total field at a point is the \emph{vector sum} of the individual fields (\cle{superposition principle}): $\vec{B}_{\text{total}} = \vec{B}_1 + \vec{B}_2 + \cdots$
\end{propriete}

\begin{exoresolu}[Field near household wiring]
A cable in a house in Douala carries a current $I = \SI{15}{A}$. Calculate the magnetic flux density at a distance $r = \SI{2.0}{cm}$ from the wire, and compare with the Earth's field ($B_E \approx 5\times 10^{-5}\ \mathrm{T}$).
\tcblower
Treating the cable as a long straight wire:
\[ B = \frac{\mu_0 I}{2\pi r} = \frac{(4\pi\times 10^{-7})(15)}{2\pi(0.020)} = \frac{2\times 10^{-7}\times 15}{0.020} = 1.5\times 10^{-4}\ \mathrm{T}. \]
So $B \approx \SI{0.15}{mT}$, about three times the Earth's field: a compass held \SI{2}{cm} from such a cable would be noticeably disturbed. Note how quickly the field falls: at \SI{20}{cm} it is only \SI{15}{\micro T}.
\end{exoresolu}

\courssub{Force Between Two Parallel Current-Carrying Conductors}

We now combine two facts already established: (i) a current produces a magnetic field, and (ii) a current-carrying conductor in a magnetic field experiences a force $F = BIL\sin\theta$. Two parallel wires therefore exert forces on each other.

\begin{methode}[Deriving the force per unit length — four steps]
\begin{enumerate}
\item Wire 1 carries $I_1$ and creates, at the position of wire 2 (distance $d$ away), the field $B_1 = \dfrac{\mu_0 I_1}{2\pi d}$, perpendicular to wire 2.
\item Wire 2 carries $I_2$ through this field; since $B_1 \perp I_2$, the force on a length $L$ of wire 2 is $F = B_1 I_2 L$.
\item Substitute: $F = \dfrac{\mu_0 I_1 I_2}{2\pi d}\,L$.
\item Divide by $L$:
\[ \frac{F}{L} = \frac{\mu_0 I_1 I_2}{2\pi d}. \]
\end{enumerate}
Fleming's left-hand rule (or the grip rule plus $\vec{F} = I\vec{L}\times\vec{B}$) shows that currents in the \emph{same} direction \textbf{attract}, and currents in \emph{opposite} directions \textbf{repel}.
\end{methode}

\begin{popfigure}
\begin{tikzpicture}[scale=1.0, >=stealth]
% Same direction: attraction
\draw[very thick, popblue] (0,-1.4) -- (0,1.4);
\draw[very thick, popblue] (2.2,-1.4) -- (2.2,1.4);
\draw[->, very thick, popblue] (0,0.7) -- (0,1.3);
\draw[->, very thick, popblue] (2.2,0.7) -- (2.2,1.3);
\node[popblue, left] at (0,1.2) {$I_1$};
\node[popblue, right] at (2.2,1.2) {$I_2$};
\draw[->, ultra thick, poporange] (0.06,-0.5) -- (0.75,-0.5);
\draw[->, ultra thick, poporange] (2.14,-0.5) -- (1.45,-0.5);
\node[poporange, above] at (1.1,-0.5) {$F$ \quad $F$};
\draw[<->, popdark] (0,-1.15) -- (2.2,-1.15);
\node[popdark, below] at (1.1,-1.15) {$d$};
\node[popdark, align=center] at (1.1,-1.9) {same direction:\\ attraction};
% Opposite directions: repulsion
\begin{scope}[xshift=6.5cm]
\draw[very thick, popblue] (0,-1.4) -- (0,1.4);
\draw[very thick, popblue] (2.2,-1.4) -- (2.2,1.4);
\draw[->, very thick, popblue] (0,0.7) -- (0,1.3);
\draw[->, very thick, popblue] (2.2,-0.7) -- (2.2,-1.3);
\node[popblue, left] at (0,1.2) {$I_1$};
\node[popblue, right] at (2.2,-1.2) {$I_2$};
\draw[->, ultra thick, poporange] (-0.06,-0.2) -- (-0.75,-0.2);
\draw[->, ultra thick, poporange] (2.26,-0.2) -- (2.95,-0.2);
\node[poporange, above] at (1.1,0.1) {$F$ \quad\quad $F$};
\draw[<->, popdark] (0,-1.15) -- (2.2,-1.15);
\node[popdark, below] at (1.1,-1.15) {$d$};
\node[popdark, align=center] at (1.1,-1.9) {opposite directions:\\ repulsion};
\end{scope}
\end{tikzpicture}
\legende{Forces between two parallel wires: attraction when the currents flow in the same direction, repulsion when they flow in opposite directions.}
\end{popfigure}

\begin{exoresolu}[Force between two bus-bars]
Two parallel copper bus-bars in a power installation are separated by $d = \SI{5.0}{cm}$ and each carries $I = \SI{200}{A}$ in the same direction. Find the force per unit length between them and state its nature.
\tcblower
\[ \frac{F}{L} = \frac{\mu_0 I_1 I_2}{2\pi d} = \frac{(4\pi\times 10^{-7})(200)(200)}{2\pi(0.050)} = \frac{2\times 10^{-7}\times 4.0\times 10^{4}}{0.050} = 0.16\ \mathrm{N\,m^{-1}}. \]
The currents are in the same direction, so the bars \emph{attract} each other with $\SI{0.16}{N}$ per metre of length. During a short-circuit, currents can reach thousands of amperes and the force grows with the \emph{product} $I_1 I_2$ — bus-bars must therefore be firmly braced.
\end{exoresolu}

\begin{savaistu}[The ampere and the 2019 redefinition]
Until 2019, the SI ampere was \emph{defined} through this very force: one ampere was the constant current which, maintained in two straight parallel conductors of infinite length and negligible circular cross-section, placed one metre apart in vacuum, would produce between them a force of $2\times 10^{-7}\ \mathrm{N}$ per metre of length — exactly what our formula gives, since $\mu_0/2\pi = 2\times 10^{-7}\ \mathrm{N\,A^{-2}}$. Since May 2019 (extra, not examinable), the ampere is instead defined by fixing the value of the elementary charge $e = 1.60\times 10^{-19}\ \mathrm{C}$, so $\mu_0$ is now a measured quantity rather than an exact one.
\end{savaistu}

\begin{aretenir}[Key points of this section]
\begin{itemize}
\item Biot-Savart law: $\mathrm{d}B = \dfrac{\mu_0}{4\pi}\dfrac{I\,\mathrm{d}l\sin\theta}{r^{2}}$, with $\mu_0 = 4\pi\times 10^{-7}\ \mathrm{H\,m^{-1}}$; direction from the right-hand rule.
\item Standard results: straight wire $B = \dfrac{\mu_0 I}{2\pi r}$; centre of a flat coil $B = \dfrac{\mu_0 N I}{2r}$; long solenoid $B = \mu_0 n I$ with $n = N/L$.
\item Ampère's law $\oint\vec{B}\cdot\mathrm{d}\vec{l} = \mu_0 I_{\text{enclosed}}$ gives these results quickly when symmetry allows a suitable Amperian loop.
\item Fields from several conductors add as vectors (superposition).
\item Two parallel conductors: $\dfrac{F}{L} = \dfrac{\mu_0 I_1 I_2}{2\pi d}$; same direction $\Rightarrow$ attraction, opposite directions $\Rightarrow$ repulsion. This force historically defined the ampere.
\end{itemize}
\end{aretenir}

\begin{exercice}[Superposition of fields]
Two long parallel wires, \SI{10}{cm} apart, carry currents $I_1 = \SI{8.0}{A}$ and $I_2 = \SI{8.0}{A}$ in opposite directions. Calculate the magnitude of the magnetic flux density at the point midway between the wires, and state its direction.
\end{exercice}

\begin{corrige}[Exercise 1]
Each wire produces at the midpoint ($r = \SI{5.0}{cm}$) a field
\[ B_1 = B_2 = \frac{\mu_0 I}{2\pi r} = \frac{(4\pi\times 10^{-7})(8.0)}{2\pi(0.050)} = 3.2\times 10^{-5}\ \mathrm{T}. \]
With opposite currents, the right-hand grip rule shows the two fields point in the \emph{same} direction at the midpoint (perpendicular to the plane of the wires). By superposition,
\[ B = B_1 + B_2 = 6.4\times 10^{-5}\ \mathrm{T} = \SI{64}{\micro T}, \]
comparable to the Earth's field. \emph{Trap:} had the currents been in the same direction, the fields would cancel at the midpoint, giving $B = 0$.
\end{corrige}

\coursec{Moving Charges, Hall Effect and Magnetic Materials}

In the previous section we saw that a \emph{current} produces a magnetic field, and we computed the flux density near wires, coils and solenoids. But an electric current is nothing more than a stream of moving charges. It is therefore natural to ask: what happens to a \emph{single charged particle} moving through a magnetic field? The answer to that question explains the Northern Lights, the working of the cathode-ray tube, mass spectrometers, the Hall probes used in every physics laboratory, and even why your bank card is kept away from strong magnets. That is the subject of this final section.

\courssub{Force on a Moving Charge in a Uniform Magnetic Field}

\textbf{From a current to a single charge.}\par
We already know that a conductor of length $L$ carrying a current $I$ in a field $B$ experiences a force $F = BIL\sin\theta$. A current $I$ means that charge $Q = It$ passes in time $t$, travelling a distance $L = vt$. Hence $IL = Qv$, and the force on the total moving charge is $F = BQv\sin\theta$. For a \emph{single} particle of charge $q$ moving with velocity $\vec{v}$:

\begin{propriete}[Magnetic force on a moving charge]
A particle of charge $q$ moving with speed $v$ at an angle $\theta$ to a uniform magnetic field of flux density $B$ experiences a force
\[ F = Bqv\sin\theta \qquad\text{or, in vector form,}\qquad \vec{F} = q\,\vec{v}\wedge\vec{B}. \]
The direction of $\vec{F}$ is given by \cle{Fleming's left-hand rule} applied to the \emph{conventional current}, i.e.\ to the direction of motion of a \emph{positive} charge: first finger = Field, second finger = current (direction of $\vec v$ for $q>0$), thumb = thrust (force).
\end{propriete}

\begin{attention}[Negative charges: reverse the direction!]
The left-hand rule gives the force on a \emph{positive} charge. For an \emph{electron} (or any negative particle), apply the rule as usual and then \textbf{reverse} the resulting direction. Equivalently, point the second finger \emph{opposite} to the electron's velocity. This is one of the most frequent sources of lost marks at the GCE.
\end{attention}

\begin{attention}[The magnetic force does no work]
Because $\vec F \perp \vec v$ at every instant, the power developed is $P = \vec F\cdot\vec v = 0$. A magnetic force can therefore \textbf{change the direction} of a moving charge but \textbf{never its speed, and never its kinetic energy}. Only an electric field can speed a charge up or slow it down.
\end{attention}

\courssub{Circular and Helical Motion in a Uniform Field}

\textbf{Observation.}\par
In the fine-beam tube used in school laboratories, an electron beam entering a region of uniform field (produced by Helmholtz coils) is seen to bend into a perfect glowing circle. Why a circle? Because the force is always perpendicular to the velocity and of constant magnitude: exactly the conditions for uniform circular motion, with the magnetic force playing the role of the centripetal force.

\begin{propriete}[Radius and period of the circular path — proof examinable]
For a charge $q$ of mass $m$ moving with speed $v$ \emph{perpendicular} to a uniform field $B$, the magnetic force supplies the centripetal force:
\[ Bqv = \frac{mv^{2}}{r} \quad\Longrightarrow\quad \boxed{\,r = \frac{mv}{Bq}\,}. \]
The time for one revolution is $T = \dfrac{2\pi r}{v}$, hence
\[ \boxed{\,T = \frac{2\pi m}{Bq}\,} \qquad\text{(independent of } v\text{)}. \]
\end{propriete}

The independence of $T$ from $v$ is remarkable and is the key to the cyclotron: whatever the particle's speed, it always takes the same time to complete a half-circle.

If the velocity has a component $v_{\parallel}$ \emph{parallel} to $\vec B$, that component is unaffected (no force acts along $\vec B$), while the perpendicular component $v_{\perp}$ produces circular motion. The superposition is a \cle{helix} of radius $r = mv_{\perp}/Bq$ and pitch $p = v_{\parallel}T$.

\begin{popfigure}
\begin{center}
\begin{tikzpicture}[scale=1.0]
% field crosses (into page)
\foreach \x in {-3.2,-1.6,0,1.6,3.2}{
  \foreach \y in {-2.6,-1.3,1.3,2.6}{
    \node[popmuted] at (\x,\y) {\small $\times$};
  }
}
\node[popmuted, anchor=west] at (2.1,-2.9) {$\vec B$ into page};
% circular path
\draw[popblue, very thick] (0,0) circle (1.9);
\draw[popblue, very thick, -latex] (1.9,0) arc (0:-40:1.9);
% charge
\fill[poporange] (1.9,0) circle (0.09);
\node[popdark, anchor=west] at (2.0,0.25) {$+q$};
% velocity (tangent)
\draw[-latex, very thick, popgreen] (1.9,0) -- (1.9,-1.4) node[below] {$\vec v$};
% force (towards centre)
\draw[-latex, very thick, poppink] (1.9,0) -- (0.7,0) node[above] {$\vec F$};
% radius
\draw[-latex, thick, popdark, dashed] (0,0) -- (1.9,0) node[midway, above] {$r$};
\fill (0,0) circle (0.05);
\end{tikzpicture}
\end{center}
\legende{A positive charge moving perpendicular to a uniform field directed into the page follows a circle. The magnetic force $\vec F$ is always directed towards the centre; the speed never changes.}
\end{popfigure}

\begin{exoresolu}[Electron circling in a magnetic field]
An electron ($e = \SI{1.60e-19}{C}$, $m_e = \SI{9.11e-31}{kg}$) is fired with speed $v = \SI{2.0e7}{m.s^{-1}}$ perpendicular to a uniform field $B = \SI{2.0e-3}{T}$. Find (a) the radius of its path, (b) the period of revolution.
\tcblower
\textbf{Solution.}
\begin{enumerate}
\item The force is perpendicular to $\vec v$, so the motion is circular with $r = \dfrac{m_e v}{eB}$.
\item (a) $r = \dfrac{9.11\times10^{-31}\times 2.0\times10^{7}}{1.60\times10^{-19}\times 2.0\times10^{-3}} = \dfrac{1.822\times10^{-23}}{3.20\times10^{-22}} = \SI{5.7e-2}{m} = \SI{5.7}{cm}$.
\item (b) $T = \dfrac{2\pi m_e}{eB} = \dfrac{2\pi\times 9.11\times10^{-31}}{3.20\times10^{-22}} = \SI{1.8e-8}{s}$. Note that $T$ does not depend on $v$.
\end{enumerate}
\end{exoresolu}

\courssub{The Lorentz Force and Crossed Fields}

When a charge moves in a region containing \emph{both} an electric field $\vec E$ and a magnetic field $\vec B$, the total force is the sum of the two contributions.

\begin{propriete}[The Lorentz force]
\[ \vec F = q\left(\vec E + \vec v \wedge \vec B\right). \]
The electric part $q\vec E$ acts whether the charge moves or not and can change its speed; the magnetic part $q\vec v\wedge\vec B$ acts only on a moving charge and never changes its speed.
\end{propriete}

\textbf{The velocity selector.}\par
Suppose $\vec E$ and $\vec B$ are perpendicular to each other and to the beam, arranged so that the electric and magnetic forces on a positive charge point in \emph{opposite} directions. A particle passes through undeflected only when the two forces balance:
\[ qE = qvB \quad\Longrightarrow\quad v = \frac{E}{B}. \]
Particles with any other speed are deflected onto the walls. The device therefore \emph{selects} one velocity, independent of the charge and mass of the particle.

\begin{methode}[Analysing crossed-field devices and $e/m$ experiments]
\begin{enumerate}
\item Draw $\vec E$, $\vec B$ and $\vec v$; mark the direction of $q\vec E$ (along $\vec E$ for $q>0$).
\item Use the left-hand rule for $q\vec v\wedge\vec B$; \textbf{reverse it for electrons}.
\item If the beam is undeflected, write $qE = qvB$, i.e.\ $v = E/B$.
\item If only $\vec B$ acts, write $Bqv = mv^{2}/r$ and use $r = mv/Bq$.
\item Combine the two equations to eliminate $v$ when finding $e/m$.
\end{enumerate}
\end{methode}

\begin{popfigure}
\begin{center}
\begin{tikzpicture}[scale=1.0]
% plates
\draw[very thick, popblue] (-2.4,1.4) -- (2.4,1.4) node[midway, above, popdark] {$+$ plate};
\draw[very thick, popblue] (-2.4,-1.4) -- (2.4,-1.4) node[midway, below, popdark] {$-$ plate};
% E field arrows
\foreach \x in {-1.6,-0.8,0,0.8,1.6}{
  \draw[-latex, popgreen] (\x,1.15) -- (\x,0.45);
}
\node[popgreen] at (2.7,0.8) {$\vec E$};
% B crosses
\foreach \x in {-1.6,-0.8,0,0.8,1.6}{
  \node[popmuted] at (\x,-0.6) {$\times$};
}
\node[popmuted] at (2.7,-0.6) {$\vec B$ in};
% beam
\draw[-latex, very thick, poporange] (-3.2,0) -- (3.2,0) node[right, popdark] {beam, $v = E/B$};
\fill[poporange] (-2.6,0) circle (0.07);
% forces
\draw[-latex, thick, poppink] (-2.6,0) -- (-2.6,0.8) node[above] {$q\vec v\wedge\vec B$};
\draw[-latex, thick, poppurple] (-2.6,0) -- (-2.6,-0.8) node[below] {$q\vec E$};
% slit
\draw[very thick, popdark] (2.4,0.5) -- (2.4,1.4);
\draw[very thick, popdark] (2.4,-0.5) -- (2.4,-1.4);
\node[popdark, align=center] at (3.6,1.1) {exit\\slit};
\end{tikzpicture}
\end{center}
\legende{Velocity selector: with $\vec E$ downwards and $\vec B$ into the page, the two forces on a positive ion cancel only when $v = E/B$; those ions pass straight through the exit slit.}
\end{popfigure}

\courssub{Measuring the Specific Charge $e/m$}

\textbf{Thomson's crossed-field method (historical, 1897).}\par
J.\,J.\ Thomson passed a cathode-ray (electron) beam through crossed $\vec E$ and $\vec B$ fields.
\begin{enumerate}
\item With both fields on, adjust them until the beam is \emph{undeflected}: then $v = E/B$.
\item Switch off $\vec E$. The beam curves in the magnetic field alone; measure the radius $r$ of the circular arc (from the deflection on the screen). Then $evB = m_ev^{2}/r$, so
\[ \frac{e}{m_e} = \frac{v}{Br} = \frac{E}{B^{2}r}. \]
\end{enumerate}
Thomson found the same value of $e/m$ whatever gas or electrode material he used — the first evidence that the electron is a universal constituent of matter.

\textbf{The fine-beam tube method (school laboratory).}\par
Electrons are accelerated from rest through a known p.d.\ $V$, then enter a uniform field $B$ (Helmholtz coils) perpendicularly and describe a visible circle of measured radius $r$.
\begin{enumerate}
\item Energy gained: $eV = \tfrac{1}{2}m_ev^{2}$, so $v = \sqrt{2eV/m_e}$.
\item Circular motion: $r = \dfrac{m_ev}{eB}$.
\item Eliminate $v$: squaring, $r^{2} = \dfrac{m_e^{2}}{e^{2}B^{2}}\cdot\dfrac{2eV}{m_e} = \dfrac{2m_eV}{eB^{2}}$, hence
\[ \frac{e}{m_e} = \frac{2V}{B^{2}r^{2}}. \]
\end{enumerate}
The accepted value is $e/m_e \approx \SI{1.76e11}{C.kg^{-1}}$.

\begin{exoresolu}[Fine-beam tube]
In a fine-beam tube experiment, electrons are accelerated through $V = \SI{250}{V}$ and describe a circle of radius $r = \SI{4.0}{cm}$ in a field $B = \SI{1.2e-3}{T}$. Calculate $e/m_e$.
\tcblower
\textbf{Solution.} Using the result derived above:
\[ \frac{e}{m_e} = \frac{2V}{B^{2}r^{2}} = \frac{2\times 250}{(1.2\times10^{-3})^{2}\times(0.040)^{2}} = \frac{500}{1.44\times10^{-6}\times 1.6\times10^{-3}} = \frac{500}{2.304\times10^{-9}} \approx \SI{2.2e11}{C.kg^{-1}}. \]
This is of the right order of magnitude; the discrepancy with the accepted value comes from uncertainties in measuring $r$ and the non-uniformity of $B$ at the edges of the coils.
\end{exoresolu}

\begin{savaistu}[The cyclotron — a GCE-useful extra]
Because $T = 2\pi m/Bq$ is independent of speed, a particle inside a cyclotron can be accelerated by an alternating p.d.\ of \emph{fixed} frequency $f = 1/T = Bq/2\pi m$, applied between two D-shaped electrodes. Each crossing of the gap adds kinetic energy $qV$; the radius grows as $r = mv/Bq$ while the revolution time stays constant. Cyclotrons produce high-energy beams for medicine (radioisotope production) and research.
\end{savaistu}

\courssub{The Hall Effect}

\textbf{Observation.}\par
Take a flat strip of conductor carrying a current $I$ and place it in a magnetic field perpendicular to its face. A small voltage appears across the \emph{width} of the strip — a direction in which no current flows! This is the \cle{Hall effect}.

\textbf{Explanation.}\par
The moving charge carriers are pushed sideways by the magnetic force $Bqv$. They accumulate on one face, leaving the opposite face charged with the opposite sign. This charge separation builds up a transverse electric field $E_H$ which opposes further drift. Equilibrium is reached when
\[ qE_H = qvB \quad\Longrightarrow\quad E_H = vB. \]
If $w$ is the width of the strip, the \cle{Hall voltage} is $V_H = E_H\,w = vBw$. Now write the current in terms of the drift velocity: $I = nqAv = nq(wt)v$, where $n$ is the number density of carriers and $t$ the thickness of the strip. Thus $v = I/(nqwt)$, and substituting:

\begin{definition}[Hall voltage and Hall coefficient]
\[ V_H = \frac{BI}{nqt} \qquad\text{and the Hall coefficient is}\qquad R_H = \frac{1}{nq}, \]
so that $V_H = R_H\,BI/t$. The \textbf{sign} of $V_H$ reveals the \textbf{sign of the charge carriers}: for positive carriers the face towards which the magnetic force pushes them becomes positive; for electrons, the polarity is reversed.
\end{definition}

\begin{popfigure}
\begin{center}
\begin{tikzpicture}[scale=1.0]
% slab
\draw[very thick, popblue, fill=popblueL] (0,0) rectangle (5,2.2);
\node[popmuted] at (2.5,1.1) {$\times\;\times\;\times$};
\node[popmuted] at (2.5,1.75) {$\vec B$ into page};
% current
\draw[-latex, very thick, poporange] (-1.2,1.1) -- (0,1.1);
\draw[-latex, very thick, poporange] (5,1.1) -- (6.2,1.1);
\node[popdark] at (-1.2,1.45) {$I$};
% charges
\node[popgreen] at (2.5,0.25) {$+\ +\ +\ +\ +$};
\node[poppink] at (2.5,1.98) {$-\ -\ -\ -\ -$};
% Hall voltage
\draw[-latex, thick, poppurple] (5.35,0.1) -- (5.35,2.1);
\node[poppurple, anchor=west] at (5.45,1.1) {$V_H$};
\node[popdark, align=center] at (2.5,-0.55) {positive carriers};
\end{tikzpicture}
\hspace{0.8cm}
\begin{tikzpicture}[scale=1.0]
\draw[very thick, popblue, fill=popblueL] (0,0) rectangle (5,2.2);
\node[popmuted] at (2.5,1.1) {$\times\;\times\;\times$};
\node[popmuted] at (2.5,1.75) {$\vec B$ into page};
\draw[-latex, very thick, poporange] (-1.2,1.1) -- (0,1.1);
\draw[-latex, very thick, poporange] (5,1.1) -- (6.2,1.1);
\node[popdark] at (-1.2,1.45) {$I$};
\node[poppink] at (2.5,0.25) {$-\ -\ -\ -\ -$};
\node[popgreen] at (2.5,1.98) {$+\ +\ +\ +\ +$};
\draw[-latex, thick, poppurple] (5.35,2.1) -- (5.35,0.1);
\node[poppurple, anchor=west] at (5.45,1.1) {$V_H$};
\node[popdark, align=center] at (2.5,-0.55) {negative carriers (electrons)};
\end{tikzpicture}
\end{center}
\legende{Hall effect in a rectangular slab. The same conventional current $I$ and field $\vec B$ give \emph{opposite} Hall-voltage polarities for positive and negative carriers, because the electrons actually drift opposite to $I$.}
\end{popfigure}

\begin{attention}[Same rule, opposite result]
Apply the left-hand rule to the \emph{conventional} current to find where carriers are pushed — that face is positive for positive carriers. But electrons drift \emph{opposite} to $I$, so they are pushed to the \emph{same} face, which then becomes \emph{negative}. The polarity of $V_H$ is therefore reversed. The Hall effect proved that in metals the carriers are negative, while in some semiconductors they behave as positive ``holes''.
\end{attention}

\textbf{Applications of the Hall effect.}
\begin{itemize}
\item \textbf{Hall probe for measuring $B$:} with $I$, $n$, $q$, $t$ fixed, $V_H \propto B$. Calibrated Hall probes are the standard school and industrial instruments for measuring magnetic flux density.
\item \textbf{Measuring carrier density $n$:} from $n = BI/(qtV_H)$; this is how the very small drift velocities in metals are confirmed.
\item \textbf{Hall sensors} are found in car ignition systems, brushless motors, smartphones (detecting flip-cover magnets) and contactless current meters.
\end{itemize}

\begin{exoresolu}[Hall voltage in a copper strip]
A copper strip of thickness $t = \SI{0.10}{mm}$ carries $I = \SI{5.0}{A}$ in a field $B = \SI{0.50}{T}$ perpendicular to its face. Copper has $n = \SI{8.5e28}{m^{-3}}$ free electrons per unit volume. Calculate the Hall voltage.
\tcblower
\textbf{Solution.}
\[ V_H = \frac{BI}{nqt} = \frac{0.50\times 5.0}{8.5\times10^{28}\times 1.60\times10^{-19}\times 1.0\times10^{-4}} = \frac{2.5}{1.36\times10^{6}} \approx \SI{1.8e-6}{V} = 1.8\ \mu\mathrm{V}. \]
Hall voltages in metals are tiny (microvolts) because $n$ is enormous; in semiconductors $n$ is much smaller, so $V_H$ is far larger and easily measured — which is why commercial Hall probes use semiconductor chips.
\end{exoresolu}

\courssub{Magnetic Materials: Dia-, Para- and Ferromagnetism}

All materials respond, weakly or strongly, to a magnetic field, because atoms contain moving electrons. We distinguish three behaviours, characterised by the \cle{relative permeability} $\mu_r = \mu/\mu_0$, which compares the field inside the material with that in vacuum.

\begin{center}
\begin{tabularx}{\linewidth}{|l|X|X|X|}
\hline
\rowcolor{popblueL} \textbf{Type} & \textbf{Origin} & \textbf{$\mu_r$} & \textbf{Examples} \\
\hline
Diamagnetic & Induced atomic currents oppose the applied field (Lenz's law at atomic scale); weakly repelled & slightly $< 1$ & copper, water, bismuth \\
\hline
Paramagnetic & Atomic magnetic moments align weakly with the field; weakly attracted & slightly $> 1$ & aluminium, platinum, oxygen \\
\hline
Ferromagnetic & Strong coupling groups moments into \emph{domains}; very strongly attracted, can stay magnetised & $\gg 1$ (up to $10^{5}$) & iron, nickel, cobalt, steel \\
\hline
\end{tabularx}
\end{center}

\textbf{Domains and hysteresis (qualitative).}\par
In an unmagnetised ferromagnet, the \cle{domains} — microscopic regions in which millions of atomic moments point the same way — are oriented at random, so their effects cancel. In an applied field, favourably oriented domains grow and others rotate, producing a very strong magnetisation. When the applied field is removed, some alignment remains: the material is \emph{permanently magnetised}. Cycling the field produces a \emph{hysteresis loop}: the magnetisation lags behind the applied field. \emph{Soft} magnetic materials (soft iron, mu-metal) have narrow loops — easy to magnetise and demagnetise, ideal for transformer cores and electromagnets. \emph{Hard} materials (steel, ferrites) have wide loops — they retain magnetisation, ideal for permanent magnets.

\courssub{Magnetic Shielding}

A high-permeability material offers field lines an ``easy path'': flux is \emph{channelled through the material} rather than crossing the space it surrounds. Enclosing a region in a shell of soft iron or mu-metal therefore diverts external field lines around the interior, leaving it nearly field-free. This is \cle{magnetic shielding}.

\begin{popfigure}
\begin{center}
\begin{tikzpicture}[scale=1.0]
% shield ring
\draw[very thick, popdark, fill=popblueL] (0,0) circle (1.1);
\draw[very thick, popdark, fill=white] (0,0) circle (0.75);
\node[popdark] at (0,0) {protected};
% field lines bending around
\foreach \y in {-1.6,-1.15,1.15,1.6}{
  \draw[-latex, thick, popblue] (-3.2,\y) .. controls (-1.4,\y) and (-1.3,{0.55*\y}) .. (0,{0.98*\y}) .. controls (1.3,{0.55*\y}) and (1.4,\y) .. (3.2,\y);
}
\draw[-latex, thick, popblue] (-3.2,-0.5) .. controls (-1.6,-0.5) and (-1.5,-0.95) .. (-0.78,-0.78);
\draw[-latex, thick, popblue] (0.78,-0.78) .. controls (1.5,-0.95) and (1.6,-0.5) .. (3.2,-0.5);
\draw[-latex, thick, popblue] (-3.2,0.5) .. controls (-1.6,0.5) and (-1.5,0.95) .. (-0.78,0.78);
\draw[-latex, thick, popblue] (0.78,0.78) .. controls (1.5,0.95) and (1.6,0.5) .. (3.2,0.5);
\node[popblue] at (-3.3,1.95) {$\vec B$};
\node[popdark, align=center] at (0,-1.7) {high-$\mu_r$ shell (mu-metal)};
\end{tikzpicture}
\end{center}
\legende{Magnetic shielding: field lines are drawn into the high-permeability enclosure and guided around the cavity, which remains almost free of field.}
\end{popfigure}

\textbf{Applications.}
\begin{itemize}
\item Shielding sensitive instruments (cathode-ray tubes, photomultipliers, electron microscopes) from the Earth's field and from nearby motors.
\item Mu-metal cases around hard drives and magnetic sensors; keeping strong magnets away from bank cards and mechanical watches.
\item Shielded rooms for precision measurements in biomagnetism (e.g.\ recording the tiny fields of the human heart).
\end{itemize}

\begin{aretenir}[Key points of the section]
\begin{itemize}
\item Force on a moving charge: $F = Bqv\sin\theta$; left-hand rule for positive charges, \textbf{reversed for electrons}.
\item The magnetic force does \textbf{no work}: speed and kinetic energy are constant; only the direction changes.
\item Perpendicular motion: $r = \dfrac{mv}{Bq}$, $T = \dfrac{2\pi m}{Bq}$ (independent of $v$); oblique motion gives a helix.
\item Lorentz force: $\vec F = q(\vec E + \vec v\wedge\vec B)$; velocity selector: $v = E/B$.
\item Specific charge: Thomson ($e/m = E/B^{2}r$) and fine-beam tube ($e/m = 2V/B^{2}r^{2}$).
\item Hall effect: $V_H = \dfrac{BI}{nqt}$; its sign reveals the sign of the carriers; used in Hall probes and sensors.
\item Dia- ($\mu_r \lesssim 1$), para- ($\mu_r \gtrsim 1$), ferromagnetic ($\mu_r \gg 1$, domains, hysteresis).
\item Magnetic shielding channels field lines through a high-$\mu_r$ enclosure, protecting the interior.
\end{itemize}
\end{aretenir}

\begin{exercice}[Helical motion and Hall probe]
\begin{enumerate}
\item A proton ($q = \SI{1.60e-19}{C}$, $m = \SI{1.67e-27}{kg}$) enters a uniform field $B = \SI{0.40}{T}$ with velocity $\SI{3.0e6}{m.s^{-1}}$ at $30^{\circ}$ to the field. Calculate the radius and the pitch of its helical path.
\item A Hall probe chip has $n = \SI{1.0e23}{m^{-3}}$, $t = \SI{0.20}{mm}$ and carries $I = \SI{10}{mA}$. It reads $V_H = \SI{6.3}{mV}$. Deduce $B$.
\end{enumerate}
\end{exercice}

\begin{corrige}[Exercise 1]
\begin{enumerate}
\item Split the velocity: $v_{\perp} = v\sin 30^{\circ} = \SI{1.5e6}{m.s^{-1}}$, $v_{\parallel} = v\cos 30^{\circ} = \SI{2.6e6}{m.s^{-1}}$.\\
Radius: $r = \dfrac{mv_{\perp}}{Bq} = \dfrac{1.67\times10^{-27}\times 1.5\times10^{6}}{0.40\times 1.60\times10^{-19}} = \SI{3.9e-2}{m}$.\\
Period: $T = \dfrac{2\pi m}{Bq} = \dfrac{2\pi\times 1.67\times10^{-27}}{6.4\times10^{-20}} = \SI{1.6e-7}{s}$.\\
Pitch: $p = v_{\parallel}T = 2.6\times10^{6}\times 1.6\times10^{-7} = \SI{0.43}{m}$.
\item From $V_H = \dfrac{BI}{nqt}$: $B = \dfrac{V_H\,nqt}{I} = \dfrac{6.3\times10^{-3}\times 1.0\times10^{23}\times 1.60\times10^{-19}\times 2.0\times10^{-4}}{1.0\times10^{-2}} = \dfrac{2.02\times10^{-2}}{1.0\times10^{-2}} \approx \SI{2.0}{T}$.
\end{enumerate}
\end{corrige}

\newpage

\coursec{Integration activity}

\coursec{Integration Activity: The Current Balance}

\courssub{Aims of the activity}

This integration activity brings together the four lessons of the chapter in a single, realistic engineering task: the design of a \cle{current balance}. By the end of the activity you should be able to:

\begin{itemize}
\item calculate the magnetic flux density inside a long solenoid using $B = \mu_0 n I$;
\item apply the torque relation $\tau = NAIB$ for a rectangular coil in a uniform magnetic field;
\item set up and solve a \cle{torque equilibrium} condition (principle of moments) to deduce an unknown current;
\item evaluate the magnetic disturbance produced by the Earth's field and by a nearby straight cable using $B = \dfrac{\mu_0 I}{2\pi r}$;
\item propose and justify a \cle{magnetic shielding} strategy, and discuss experimental uncertainties like a real physicist.
\end{itemize}

The activity is designed for group work (3--4 students). Each group presents its design, calculations and error discussion on a poster, as at a science fair.

\begin{methode}[General approach for the current balance]
\begin{enumerate}
\item \textbf{Create the field:} Use a long solenoid to produce a known, uniform magnetic field $B = \mu_0 n I_s$.
\item \textbf{Let the field act on a current:} Place a rectangular coil of $N$ turns, area $A$, carrying current $I_c$ in the field. The magnetic torque is $\tau = NAIB$ (maximum when the coil plane contains $\vec{B}$).
\item \textbf{Balance the torque:} Mount the coil on a frictionless pivot with a counterweight arm. A rider of mass $m$ at distance $b$ provides a gravitational torque $\tau_g = mgb$.
\item \textbf{Equate torques:} At equilibrium $NAIB = mgb$. This gives a direct relation between the measured position $b$ and the unknown current $I_c$.
\item \textbf{Assess and correct disturbances:} Quantify external fields (Earth, nearby cables) and use reversal techniques or shielding to minimise systematic errors.
\end{enumerate}
\end{methode}

\courssub{Situation}

\begin{experience}[The physics club of Lyc\'ee de Nkolbisson]
The science club of a lyc\'ee in Yaound\'e wants to build a device that measures electric current \emph{without any electronic ammeter}, using only magnetism and a balance. The idea is the classic \cle{current balance}: a rectangular coil is mounted on a horizontal pivot like the beam of a balance; a magnetic field exerts a torque on the coil, and this torque is balanced by the weight of a small rider (a piece of plasticine) sliding along the arm.

\medskip
\textbf{The apparatus.} A rectangular coil of $N = 50$ turns of insulated copper wire has a horizontal width $l = 4.0\ \mathrm{cm}$ (the side perpendicular to the field, on which the magnetic forces act) and a length $d = 6.0\ \mathrm{cm}$ along the pivot axis. The coil is pivoted about a horizontal axis passing through its centre, so that it can tilt like a balance beam. It carries a current $I_c$ supplied by a battery and a rheostat.

The coil sits at the centre of a \textbf{long solenoid} of length $L = 40\ \mathrm{cm}$ wound with $N_s = 800$ turns, carrying an adjustable current $I_s$. The solenoid's axis is horizontal and perpendicular to the plane of the coil when the coil is horizontal, so that the field $\vec{B}$ inside the solenoid lies in the plane of the coil and is perpendicular to the sides of length $l$.

A light rigid arm of length $b = 10.0\ \mathrm{cm}$ is fixed to the coil frame, on the opposite side of the pivot from the coil's centre of mass. A rider of mass $m = 2.0\ \mathrm{g}$ can be placed at the end of this arm.

\medskip
\textbf{The environment.} The laboratory is in Yaound\'e, where the horizontal component of the Earth's magnetic field is about $B_E \approx 3.0\times 10^{-5}\ \mathrm{T}$. A straight cable feeding the school's photocopier, carrying $I = 20\ \mathrm{A}$, runs in the wall at a perpendicular distance $r = 50\ \mathrm{cm}$ from the apparatus.

\medskip
\textbf{Your mission.} Design the measurement procedure, predict the current needed for balance, assess the environmental disturbances, and recommend improvements. Take $\mu_0 = 4\pi \times 10^{-7}\ \mathrm{H\,m^{-1}}$ and $g = 9.81\ \mathrm{m\,s^{-2}}$.
\end{experience}

\begin{popfigure}
\begin{tikzpicture}[scale=0.85, every node/.style={font=\small}]
% Solenoid cross-section (side view)
\draw[popblue, thick, fill=popblueL!30] (-3.5,1.5) rectangle (3.5,2.5);
\draw[popblue, thick, fill=popblueL!30] (-3.5,-2.5) rectangle (3.5,-1.5);
\draw[popblue, thick] (-3.5,-1.5) -- (-3.5,1.5);
\draw[popblue, thick] (3.5,-1.5) -- (3.5,1.5);
\node[popblue, align=center] at (0,2.8) {Long solenoid\\$N_s=800$, $L=40$ cm};
% Magnetic field lines (horizontal, left to right)
\foreach \y in {-1.0, 0, 1.0} {
  \draw[->, popmuted, thick] (-3.0,\y) -- (3.0,\y);
}
\node[popmuted] at (3.5,1.0) {$\vec{B}$};
% Coil (rectangular loop, seen edge-on as a line? Better: front view? Let's do a 3D-ish side view)
% Actually, side view: coil is a rectangle in vertical plane? The problem says coil plane contains B when horizontal.
% Let's draw the coil as a rectangle in the vertical plane (pivot horizontal).
% Pivot at (0,0). Coil extends horizontally (width l) and vertically (length d)? Wait.
% "horizontal width l = 4.0 cm (the side perpendicular to the field)" and "length d = 6.0 cm along the pivot axis".
% So the coil is vertical? No: "coil is pivoted about a horizontal axis passing through its centre, so that it can tilt like a balance beam."
% The coil plane is vertical when horizontal? Let's interpret: The coil is a rectangular frame. The pivot axis is horizontal along the length d. The coil can tilt in the vertical plane. The magnetic field is horizontal and perpendicular to the sides of length l.
% So in side view (looking along the pivot axis), we see the coil as a line of length l? Actually we see the two sides of length l as two horizontal lines separated by distance d? No, d is along the pivot axis, so in side view we see the coil as a rectangle of width l (horizontal) and height? The coil has thickness? Let's simplify.
% We'll draw a schematic side view showing the solenoid, the coil as a rectangle (width l, height representing the vertical dimension? Actually the coil is vertical? The problem says "plane of the coil contains B when the coil is horizontal". That means the coil plane is horizontal? Wait: "solenoid's axis is horizontal and perpendicular to the plane of the coil when the coil is horizontal". If the coil is horizontal, its plane is horizontal. The solenoid axis is horizontal and perpendicular to that plane. So B is horizontal and lies in the coil plane? "so that the field B inside the solenoid lies in the plane of the coil and is perpendicular to the sides of length l." So the coil is horizontal (like a flat loop), B is horizontal and perpendicular to the sides of length l. The coil is pivoted about a horizontal axis through its centre, along the length d. So the coil can tilt up/down. The arm is fixed to the coil frame on the opposite side of the pivot.
% This is hard to draw in 2D. Let's draw a 3D-ish schematic using TikZ with perspective.
% Better: draw a top view and a side view? The poster will have a diagram. For the course, a clear schematic is enough.
% I'll draw a simplified side view showing the solenoid as a cylinder, the coil as a loop with pivot, arm and rider.
\end{tikzpicture}
\legende{Schematic of the current balance (side view along the pivot axis). The solenoid produces a horizontal magnetic field $\vec{B}$ perpendicular to the sides of length $l$ of the rectangular coil. The coil is pivoted at its centre; a rider on the balance arm provides a restoring torque.}
\end{popfigure}

\courssub{Tasks}

\begin{enumerate}
\item \textbf{Field of the solenoid.} Calculate the turn density $n$ of the solenoid in turns per metre. Show that when the solenoid carries $I_s = 2.0\ \mathrm{A}$, the flux density at its centre is $B \approx 5.0\ \mathrm{mT}$. State the assumptions behind the formula $B = \mu_0 n I_s$.

\item \textbf{Torque on the coil.} The coil carries $I_c = 0.50\ \mathrm{A}$. Using $\tau = N A I_c B$, where $A = l \times d$, calculate the magnetic torque on the coil when its plane contains $\vec{B}$. Indicate, with the help of Fleming's left-hand rule, why the forces on the two sides of length $l$ form a couple, and why the forces on the other two sides produce no torque about the pivot.

\item \textbf{Balance condition.} Write the equilibrium condition (principle of moments) between the magnetic torque and the torque of the rider's weight $mgb$. Deduce the coil current $I_c$ required for balance when the rider sits at the end of the arm ($b = 10.0\ \mathrm{cm}$), with $I_s = 2.0\ \mathrm{A}$.

\item \textbf{Inverse use: measuring an unknown current.} The device is now used to measure an \emph{unknown} coil current $I_c$. With $I_s = 2.0\ \mathrm{A}$, balance is obtained when the rider is placed at $b = 7.5\ \mathrm{cm}$. Find $I_c$.

\item \textbf{Disturbance by the power cable.} Calculate the flux density $B_{\text{cable}}$ produced by the photocopier cable at the position of the coil. Compare it with the solenoid field (express it as a percentage) and with the Earth's field $B_E$. Should the cable be worried about? Justify.

\item \textbf{Disturbance by the Earth's field.} Estimate the extra torque that the Earth's horizontal field could exert on the coil in the worst orientation, and the corresponding error on the measured current in question 4. Suggest an experimental trick (reversing currents and averaging) to cancel this systematic effect.

\item \textbf{Shielding and design improvements.} Propose a magnetic shielding solution for the apparatus, naming a suitable material and explaining, in terms of ferromagnetism, how it works. Give two other improvements that would increase the sensitivity of the balance.

\item \textbf{Poster.} Prepare your poster: diagram of the device, the three key formulas with their derivations sketched, the numerical results, and an honest error discussion.
\end{enumerate}

\courssub{Model Answers}

\begin{exoresolu}[Task 1: Field of the solenoid]
\textbf{Statement:} Calculate the turn density $n$ of the solenoid in turns per metre. Show that when the solenoid carries $I_s = 2.0\ \mathrm{A}$, the flux density at its centre is $B \approx 5.0\ \mathrm{mT}$. State the assumptions behind the formula $B = \mu_0 n I_s$.
\tcblower
\textbf{Solution:}
The turn density (turns per unit length) is
\[
n = \frac{N_s}{L} = \frac{800}{0.40\ \mathrm{m}} = 2000\ \mathrm{turns\,m^{-1}}.
\]
For a long solenoid carrying current $I_s$, the magnetic flux density at its centre is
\[
B = \mu_0 n I_s = (4\pi \times 10^{-7}\ \mathrm{H\,m^{-1}}) \times 2000\ \mathrm{m^{-1}} \times 2.0\ \mathrm{A}
= 5.03 \times 10^{-3}\ \mathrm{T} \approx 5.0\ \mathrm{mT}.
\]
\textbf{Assumptions for $B = \mu_0 n I_s$:}
\begin{itemize}
\item The solenoid is \emph{infinitely long} (length $\gg$ diameter), so the field is uniform inside and zero outside.
\item The windings are tightly packed and perfectly helical (no pitch angle).
\item The current is steady (magnetostatics).
\item The observation point is on the axis, far from the ends.
\end{itemize}
Our solenoid has $L = 40\ \mathrm{cm}$; a typical diameter might be $\sim 4\ \mathrm{cm}$, so $L/d \approx 10$. The formula gives a good approximation at the centre (error $< 2\%$).
\end{exoresolu}

\begin{exoresolu}[Task 2: Torque on the coil]
\textbf{Statement:} The coil carries $I_c = 0.50\ \mathrm{A}$. Using $\tau = N A I_c B$, where $A = l \times d$, calculate the magnetic torque on the coil when its plane contains $\vec{B}$. Indicate, with the help of Fleming's left-hand rule, why the forces on the two sides of length $l$ form a couple, and why the forces on the other two sides produce no torque about the pivot.
\tcblower
\textbf{Solution:}
The area of the coil is
\[
A = l \times d = 0.040\ \mathrm{m} \times 0.060\ \mathrm{m} = 2.4 \times 10^{-3}\ \mathrm{m^2}.
\]
With the plane of the coil containing $\vec{B}$ (maximum torque position),
\[
\tau = N A I_c B = 50 \times (2.4 \times 10^{-3}\ \mathrm{m^2}) \times 0.50\ \mathrm{A} \times (5.03 \times 10^{-3}\ \mathrm{T})
= 3.02 \times 10^{-4}\ \mathrm{N\,m} \approx 3.0 \times 10^{-4}\ \mathrm{N\,m}.
\]

\textbf{Origin of the torque (Fleming's left-hand rule):}
\begin{itemize}
\item On each of the two sides of length $l$, the current is perpendicular to $\vec{B}$. The force per turn is $F = B I_c l$. By Fleming's left-hand rule (Field $\to$ index finger, Current $\to$ middle finger, Force $\to$ thumb), the forces on the two sides are equal in magnitude, opposite in direction, and perpendicular to the coil plane. They are separated by the distance $d$ (the length of the coil along the pivot axis). These two forces form a \cle{couple} with torque $\tau_{\text{couple}} = F \times d = B I_c l d = B I_c A$ per turn. For $N$ turns, $\tau = N B I_c A$.
\item On the two sides of length $d$, the current is parallel to $\vec{B}$ (since $\vec{B}$ lies in the coil plane and is perpendicular to $l$). Hence the magnetic force on these sides is zero ($F = B I d \sin 0 = 0$). Even if there were a small component, their lines of action pass through the pivot axis, so their moment about the pivot is zero.
\end{itemize}
\end{exoresolu}

\begin{exoresolu}[Task 3: Balance condition]
\textbf{Statement:} Write the equilibrium condition (principle of moments) between the magnetic torque and the torque of the rider's weight $mgb$. Deduce the coil current $I_c$ required for balance when the rider sits at the end of the arm ($b = 10.0\ \mathrm{cm}$), with $I_s = 2.0\ \mathrm{A}$.
\tcblower
\textbf{Solution:}
At equilibrium, the sum of moments about the pivot is zero:
\[
\tau_{\text{magnetic}} = \tau_{\text{gravitational}} \quad \Longrightarrow \quad N A I_c B = m g b.
\]
Solving for $I_c$:
\[
I_c = \frac{m g b}{N A B}
= \frac{(2.0 \times 10^{-3}\ \mathrm{kg}) \times 9.81\ \mathrm{m\,s^{-2}} \times 0.100\ \mathrm{m}}
{50 \times (2.4 \times 10^{-3}\ \mathrm{m^2}) \times (5.03 \times 10^{-3}\ \mathrm{T})}
= 3.25\ \mathrm{A} \approx 3.2\ \mathrm{A}.
\]
\textbf{Note:} This current is relatively large for thin copper wire (risk of heating). In practice, one would increase the rider mass $m$, lengthen the arm $b$, increase $N$ or $A$, or increase $B$ (larger $I_s$) to reduce the operating current.
\end{exoresolu}

\begin{exoresolu}[Task 4: Measuring an unknown current]
\textbf{Statement:} The device is now used to measure an \emph{unknown} coil current $I_c$. With $I_s = 2.0\ \mathrm{A}$, balance is obtained when the rider is placed at $b = 7.5\ \mathrm{cm}$. Find $I_c$.
\tcblower
\textbf{Solution:}
The same balance condition applies. With $b = 7.5\ \mathrm{cm} = 0.075\ \mathrm{m}$:
\[
I_c = \frac{m g b}{N A B}
= \frac{(2.0 \times 10^{-3}) \times 9.81 \times 0.075}
{50 \times 2.4 \times 10^{-3} \times 5.03 \times 10^{-3}}
= 2.44\ \mathrm{A} \approx 2.4\ \mathrm{A}.
\]
The linear relationship $I_c \propto b$ means the rider position acts as a direct, calibrated scale for the current. No electronic meter is needed.
\end{exoresolu}

\begin{exoresolu}[Task 5: Disturbance by the power cable]
\textbf{Statement:} Calculate the flux density $B_{\text{cable}}$ produced by the photocopier cable at the position of the coil. Compare it with the solenoid field (express it as a percentage) and with the Earth's field $B_E$. Should the cable be worried about? Justify.
\tcblower
\textbf{Solution:}
For a long straight conductor carrying current $I = 20\ \mathrm{A}$ at perpendicular distance $r = 0.50\ \mathrm{m}$:
\[
B_{\text{cable}} = \frac{\mu_0 I}{2\pi r}
= \frac{4\pi \times 10^{-7} \times 20}{2\pi \times 0.50}
= 8.0 \times 10^{-6}\ \mathrm{T}.
\]
Comparison with the solenoid field $B = 5.03 \times 10^{-3}\ \mathrm{T}$:
\[
\frac{B_{\text{cable}}}{B} \times 100\% = \frac{8.0 \times 10^{-6}}{5.03 \times 10^{-3}} \times 100\% \approx 0.16\%.
\]
Comparison with Earth's horizontal field $B_E = 3.0 \times 10^{-5}\ \mathrm{T}$:
\[
B_{\text{cable}} = 8.0 \times 10^{-6}\ \mathrm{T} < B_E = 3.0 \times 10^{-5}\ \mathrm{T}.
\]
\textbf{Conclusion:} The cable's field is only $0.16\%$ of the solenoid field and smaller than the Earth's field. At the precision of a school balance (typically $1\%$--$2\%$), this disturbance is negligible. However, if the cable were at $5\ \mathrm{cm}$ instead of $50\ \mathrm{cm}$, the field would be ten times larger ($8.0 \times 10^{-5}\ \mathrm{T}$), reaching $1.6\%$ of $B$ and exceeding $B_E$ — a significant systematic error. \textbf{Keep the apparatus away from walls carrying large currents.}
\end{exoresolu}

\begin{exoresolu}[Task 6: Disturbance by the Earth's field]
\textbf{Statement:} Estimate the extra torque that the Earth's horizontal field could exert on the coil in the worst orientation, and the corresponding error on the measured current in question 4. Suggest an experimental trick (reversing currents and averaging) to cancel this systematic effect.
\tcblower
\textbf{Solution:}
The Earth's horizontal field $B_E = 3.0 \times 10^{-5}\ \mathrm{T}$ is fixed in direction (magnetic North). The worst case occurs when $B_E$ is perpendicular to the coil plane (i.e. parallel to the coil's magnetic moment), adding directly to the solenoid field $B$ (or opposing it). The extra torque is
\[
\tau_E = N A I_c B_E.
\]
Using the unknown current from Task 4 ($I_c = 2.44\ \mathrm{A}$):
\[
\tau_E = 50 \times (2.4 \times 10^{-3}) \times 2.44 \times (3.0 \times 10^{-5})
= 8.8 \times 10^{-6}\ \mathrm{N\,m}.
\]
The useful torque (from the solenoid) is
\[
\tau = N A I_c B = 50 \times 2.4 \times 10^{-3} \times 2.44 \times 5.03 \times 10^{-3}
= 1.47 \times 10^{-3}\ \mathrm{N\,m}.
\]
The relative systematic error on the torque (and hence on the measured current, since $I_c \propto \tau$) is
\[
\frac{\Delta I_c}{I_c} = \frac{\tau_E}{\tau} = \frac{B_E}{B}
= \frac{3.0 \times 10^{-5}}{5.03 \times 10^{-3}} \approx 0.60\%.
\]

\textbf{Cancellation trick (reversal method):}
Perform the measurement twice:
\begin{enumerate}
\item First measurement: $I_s$ and $I_c$ in given directions $\to$ balance at $b_1$.
\item Second measurement: Reverse \emph{both} $I_s$ and $I_c$ (so the magnetic torque $\tau = NAIB$ keeps the same sign because $B \propto I_s$ and $\tau \propto I_c B$). The Earth's torque $\tau_E \propto I_c B_E$ changes sign because $I_c$ is reversed while $B_E$ is fixed.
\end{enumerate}
The two balance positions $b_1$ and $b_2$ satisfy:
\[
NAI_c B \pm NAI_c B_E = mgb_{1,2}.
\]
Averaging: $\frac{b_1 + b_2}{2} = \frac{NAI_c B}{mg}$, cancelling the Earth's field contribution to first order.
\end{exoresolu}

\begin{exoresolu}[Task 7: Shielding and design improvements]
\textbf{Statement:} Propose a magnetic shielding solution for the apparatus, naming a suitable material and explaining, in terms of ferromagnetism, how it works. Give two other improvements that would increase the sensitivity of the balance.
\tcblower
\textbf{Solution:}

\textbf{Magnetic shielding:}
Enclose the coil region in an open-ended cylinder made of \cle{mumetal} (a nickel--iron alloy with relative permeability $\mu_r \sim 10^5$) or, more cheaply, \cle{mild steel} ($\mu_r \sim 10^3$). 
\begin{propriete}[Ferromagnetic shielding]
Ferromagnetic materials have a very high magnetic permeability $\mu = \mu_0 \mu_r$. They offer a path of very low \cle{magnetic reluctance} ($\mathcal{R} = l/(\mu A)$). External magnetic field lines prefer to travel through the ferromagnetic walls rather than through the air inside the shield. The field inside the cavity is strongly attenuated (by a factor $\sim 1/\mu_r$ for a closed shell). The shield must not form a closed loop around the solenoid's own field lines, or it would short-circuit the useful field. In practice, place the \emph{entire balance} (coil + pivot + arm) inside the shield, with the solenoid outside or with the shield slotted to allow the solenoid's field to enter.
\end{propriete}

\textbf{Two improvements for increased sensitivity:}
\begin{enumerate}
\item \textbf{Increase the magnetic torque per unit current:} Use more turns $N$ on the coil, increase the coil area $A$ (larger $l$ or $d$), or increase the solenoid field $B$ (more turns $N_s$, higher $I_s$, or a ferromagnetic core inside the solenoid). This gives a larger deflection for a given current.
\item \textbf{Increase the mechanical leverage:} Lengthen the balance arm $b$ or use a lighter rider $m$. Since $I_c \propto mgb/(NAB)$, a larger $b$ or smaller $m$ means a smaller current produces a measurable rider displacement.
\item \textbf{Better detection of balance:} Replace the rider with a pointer and calibrated scale, or use an optical lever (a small mirror on the pivot reflecting a light spot onto a distant scale) to detect tiny angular displacements. This allows measurement of very small currents without moving a rider.
\end{enumerate}
\end{exoresolu}

\begin{attention}[Common mistakes to avoid on the poster]
\begin{itemize}
\item Using $n = N_s$ instead of $n = N_s/L$: $n$ is a turn \emph{density} (turns per metre), not the total number of turns.
\item Forgetting that $\tau = NAIB$ is the \emph{maximum} torque (plane of coil parallel to $\vec{B}$); in general $\tau = NAIB\cos\theta$ where $\theta$ is the angle between the coil normal and $\vec{B}$.
\item Mixing units: always convert centimetres to metres and grams to kilograms before substituting into formulas.
\item Comparing fields without comparing \emph{directions}: two fields of similar magnitude may add or cancel depending on their relative orientation. Always consider vector addition.
\item Neglecting the reversal method for Earth's field: a single measurement always contains a systematic error of order $B_E/B$.
\end{itemize}
\end{attention}

\begin{aretenir}[What this activity consolidates]
One coherent chain of physics: a current \emph{creates} a field ($B=\mu_0 n I_s$), the field \emph{exerts forces} on another current ($F = BIl$), the forces form a \emph{couple} ($\tau = NAIB$), and the couple is \emph{weighed} against gravity ($\tau = mgb$). Add the disturbances ($B = \mu_0 I/2\pi r$, Earth's field) and their cure (reversal method, ferromagnetic shielding), and you have revised the whole chapter through one real instrument.
\end{aretenir}

\begin{savaistu}[From the laboratory to the definition of the ampere]
The current balance is not just a student experiment. Historically, the \cle{ampere} was defined (until 2019) via the force between two parallel current-carrying conductors: ``The ampere is that constant current which, if maintained in two straight parallel conductors of infinite length, of negligible circular cross-section, and placed 1 metre apart in vacuum, would produce between these conductors a force equal to $2 \times 10^{-7}$ newton per metre of length.'' This is a direct application of $F/L = \mu_0 I_1 I_2 / (2\pi r)$. The modern definition fixes the elementary charge $e = 1.602176634 \times 10^{-19}\ \mathrm{C}$, but the current balance remains a primary method for realising the ampere with high precision in national metrology institutes.
\end{savaistu}

\newpage

\coursec{Methods and worked examples}

\courssub{Method 1 — Predicting the Direction of the Force on a Current-Carrying Conductor}

A current-carrying conductor placed in a magnetic field experiences a force: this is the \emph{motor effect}, observed whenever a wire between the poles of a magnet jumps when the current is switched on. Before computing magnitudes, the first skill required at the GCE is to \textbf{predict the direction} of this force unambiguously.

\begin{methode}[Using Fleming's Left-Hand Rule]
\begin{enumerate}
\item \textbf{Identify the field direction.} Draw or read the direction of the magnetic flux density $\vec{B}$ (from North pole to South pole outside a magnet, or use the right-hand grip rule if the field is produced by a current).
\item \textbf{Identify the current direction.} The current $I$ is the \emph{conventional} current: from the positive terminal to the negative terminal of the supply. If the moving charges are electrons, the conventional current is \textbf{opposite} to their velocity.
\item \textbf{Apply Fleming's left-hand rule.} Hold the thumb, first finger and second finger of the \textbf{left} hand mutually at right angles:
\begin{itemize}
\item \textbf{F}irst finger $\rightarrow$ \textbf{F}ield ($\vec{B}$);
\item se\textbf{C}ond finger $\rightarrow$ \textbf{C}urrent ($I$);
\item thu\textbf{M}b $\rightarrow$ \textbf{M}otion (the force $\vec{F}$).
\end{itemize}
\item \textbf{Check with the vector product.} The force is $\vec{F}=I\vec{L}\wedge\vec{B}$, with magnitude
\[ F = BIL\sin\theta, \]
where $\theta$ is the angle between the conductor and $\vec{B}$. The force is maximum when the conductor is perpendicular to the field ($F=BIL$) and zero when it is parallel to the field.
\item \textbf{Conclude in words and on the diagram}: state the direction of the force (e.g. ``vertically upwards'') and mark it clearly.
\end{enumerate}
\end{methode}

\begin{popfigure}
\begin{tikzpicture}[scale=1.0, >=stealth]
% Field lines N to S
\draw[thick, popblue, ->] (-2.2,0) -- (2.2,0) node[right, popblue]{$\vec{B}$ (N $\to$ S)};
\fill[popblueL] (-2.9,-0.6) rectangle (-2.2,0.6);
\node[popblue] at (-2.55,0) {\textbf{N}};
\fill[popblueL] (2.9,-0.6) rectangle (2.2,0.6);
\node[popblue] at (2.55,0) {\textbf{S}};
% Wire
\draw[very thick, popdark] (0,-1.1) -- (0,1.1);
\node[popdark] at (0.35,1.15) {wire};
% Current into page
\draw[poporange, thick] (0,0) circle (0.18);
\node[poporange] at (0,0) {\small $\times$};
\node[poporange, align=center] at (1.5,-1.0) {$I$ into\\ the page};
% Force up
\draw[very thick, poppink, ->] (0,0.25) -- (0,1.9) node[above, poppink]{$\vec{F}$};
\end{tikzpicture}
\legende{Fleming's left-hand rule applied to a wire perpendicular to a uniform field: field to the right, current into the page, force vertically upwards.}
\end{popfigure}

\begin{attention}
Never use the \emph{right} hand for the force on a conductor: the right hand is reserved for finding the field produced by a current (grip rule). Mixing the two hands is the most frequent error at the GCE. Also remember that $F=BIL\sin\theta$ requires $\theta$ to be the angle between the \emph{conductor} and $\vec{B}$, not between the force and $\vec{B}$: the force is always perpendicular to both.
\end{attention}

\begin{exoresolu}[Force on a horizontal wire between magnet poles]
A straight copper wire of length $L=\SI{15}{cm}$ is placed horizontally between the poles of a horseshoe magnet producing a uniform horizontal field of flux density $B=\SI{0.40}{T}$, directed from the North pole (on the left) to the South pole (on the right). The wire is perpendicular to the field and carries a current $I=\SI{5.0}{A}$ flowing away from the observer (into the page). Determine the magnitude and the direction of the force on the wire.
\tcblower
\textbf{Step 1 — Data and geometry.} $B=\SI{0.40}{T}$, $I=\SI{5.0}{A}$, $L=\SI{0.15}{m}$. The wire is perpendicular to $\vec{B}$, so $\theta=90^{\circ}$ and $\sin\theta=1$.

\textbf{Step 2 — Magnitude.}
\[ F = BIL\sin\theta = 0.40\times 5.0\times 0.15\times 1 = \SI{0.30}{N}. \]

\textbf{Step 3 — Direction (Fleming's left hand).}
\begin{itemize}
\item First finger (Field): points to the \textbf{right} (N $\rightarrow$ S);
\item Second finger (Current): points \textbf{into the page};
\item Thumb (Motion): then points \textbf{vertically upwards}.
\end{itemize}

\textbf{Conclusion.} The wire experiences a force of magnitude $\SI{0.30}{N}$ directed \textbf{vertically upwards}.

\textbf{Remark.} If the current were reversed, the force would point downwards; if the wire were turned parallel to $\vec{B}$, the force would vanish. This up/down switching as the current is reversed is exactly what makes a loudspeaker cone vibrate.
\end{exoresolu}

\courssub{Method 2 — Computing the Force on a Conductor and Using It to Measure B}

The relation $F=BIL\sin\theta$ is not only a calculating tool: because it is \emph{linear}, it underpins the standard laboratory measurement of magnetic flux density with a \cle{current balance}.

\begin{methode}[The current balance: $F=BIL\sin\theta$]
\begin{enumerate}
\item \textbf{Write the equilibrium condition.} If the conductor rests in equilibrium (e.g. on a balance), equate the magnetic force to the restoring force (weight, spring force): $BIL = mg$ for a perpendicular arrangement.
\item \textbf{Convert all data to SI units} before substituting: lengths in metres, currents in amperes, masses in kilograms.
\item \textbf{Solve for the unknown} ($B$, $I$ or $L$) and give the result with its unit and a sensible number of significant figures.
\item \textbf{Exploit linearity.} Since $F=BIL$, a graph of $F$ against $I$ is a straight line through the origin of gradient $BL$; a graph of $F$ against $L$ has gradient $BI$. This is the standard way of measuring $B$ experimentally.
\item \textbf{Check the angle.} If the conductor is not perpendicular to the field, keep the factor $\sin\theta$; never drop it silently.
\end{enumerate}
\end{methode}

\begin{exoresolu}[Measuring B with a current balance]
A rectangular frame of width $L=\SI{5.0}{cm}$ carries a current $I$ and is placed so that one horizontal side lies in a uniform magnetic field perpendicular to it. The frame is balanced on a knife edge; when the current is switched on, a rider of mass $m=\SI{0.20}{g}$ must be added to restore equilibrium. The current is $I=\SI{2.5}{A}$. Take $g=\SI{9.81}{m.s^{-2}}$.
\begin{enumerate}
\item Calculate the magnetic flux density $B$.
\item The experiment is repeated for several currents. Sketch the graph of the magnetic force $F$ against $I$ and state how $B$ is obtained from it.
\end{enumerate}
\tcblower
\textbf{1. Calculation of $B$.}

At equilibrium, the magnetic force on the side of length $L$ balances the weight of the rider:
\[ BIL = mg \quad\Longrightarrow\quad B=\frac{mg}{IL}. \]
Converting to SI: $m=0.20\times10^{-3}\ \mathrm{kg}$, $L=5.0\times10^{-2}\ \mathrm{m}$.
\[ B=\frac{0.20\times10^{-3}\times 9.81}{2.5\times 5.0\times10^{-2}}
=\frac{1.962\times10^{-3}}{0.125}
=1.57\times10^{-2}\ \mathrm{T}. \]
\[ \boxed{B\approx \SI{1.6e-2}{T} = \SI{16}{mT}} \]

\textbf{2. Graph of $F$ against $I$.}

Since $F=BIL$ with $B$ and $L$ constant, $F\propto I$: the graph is a \textbf{straight line through the origin}. Its gradient is
\[ \text{gradient}=BL \quad\Longrightarrow\quad B=\frac{\text{gradient}}{L}. \]
This method averages out random errors and reveals any systematic zero error of the balance (the line would then miss the origin).
\end{exoresolu}

\courssub{Method 3 — Torque on a Rectangular Coil and the Electric Motor}

A single straight wire in a field is merely pushed; but mount a \emph{coil} on an axle and the forces on its two sides form a \textbf{couple}: the coil rotates. This is the working principle of the moving-coil galvanometer and of the electric motor.

\begin{methode}[Applying $\tau = BAIN\cos\theta$]
\begin{enumerate}
\item \textbf{Identify the geometry}: number of turns $N$, area of one turn $A$ (for a rectangle $A=ab$), current $I$, flux density $B$.
\item \textbf{Locate the angle $\theta$}: it is the angle between the \textbf{normal to the plane of the coil} and $\vec{B}$. (Some texts define $\theta$ from the plane of the coil; then $\sin\theta$ appears. Read the question carefully.)
\item \textbf{Write the torque}: each side of the coil perpendicular to $\vec{B}$ experiences $F=NBIL$; the two forces form a couple whose torque is
\[ \tau = BAIN\cos\theta. \]
\item \textbf{Extreme positions}:
\begin{itemize}
\item coil plane \emph{parallel} to $\vec{B}$ ($\theta=0$): maximum torque $\tau_{\max}=BAIN$;
\item coil plane \emph{perpendicular} to $\vec{B}$ ($\theta=90^{\circ}$): zero torque (equilibrium position).
\end{itemize}
\item \textbf{For the motor}: explain the role of the split-ring commutator (dc) — it reverses the current every half-turn so that the torque keeps the same sense of rotation. In an ac motor, slip rings feed the coil and the torque alternates.
\end{enumerate}
\end{methode}

\begin{popfigure}
\begin{tikzpicture}[scale=1.0, >=stealth]
% Field
\foreach \y in {-1.0,-0.35,0.3,0.95} {
  \draw[popblue, ->] (-2.6,\y) -- (2.6,\y);
}
\node[popblue] at (2.9,0.95) {$\vec{B}$};
% Coil seen in perspective: rectangle tilted
\draw[very thick, popdark] (-0.9,-0.75) -- (0.9,-0.45) -- (0.9,0.75) -- (-0.9,0.45) -- cycle;
% Forces on the two active sides
\draw[very thick, poppink, ->] (-0.9,0.45) -- (-0.9,1.5) node[above, poppink]{$\vec{F}$};
\draw[very thick, poppink, ->] (0.9,-0.45) -- (0.9,-1.35) node[below, poppink]{$\vec{F}$};
% Axis
\draw[dashed, popmuted] (0,-1.5) -- (0,1.7) node[above, popmuted]{axle};
% Current arrows on active sides
\draw[poporange, ->] (-0.9,-0.5) -- (-0.9,0.1);
\draw[poporange, ->] (0.9,0.5) -- (0.9,-0.1);
\node[poporange] at (-1.35,-0.2) {$I$};
\node[poporange] at (1.35,0.2) {$I$};
\end{tikzpicture}
\legende{A rectangular coil in a uniform field: the forces on the two current-carrying sides form a couple of torque $\tau=BAIN\cos\theta$, rotating the coil about its axle.}
\end{popfigure}

\begin{exoresolu}[Torque on the coil of a small dc motor]
The armature of a small dc motor consists of a rectangular coil of $N=200$ turns, of dimensions $\SI{4.0}{cm}\times\SI{3.0}{cm}$, carrying a current $I=\SI{1.5}{A}$ in a field that is effectively uniform, of magnitude $B=\SI{0.25}{T}$, parallel to the plane of the coil.
\begin{enumerate}
\item Calculate the maximum torque on the coil.
\item Calculate the torque when the normal to the coil makes an angle of $30^{\circ}$ with the field.
\item Explain why the torque does not average to zero over a complete revolution in a dc motor.
\end{enumerate}
\tcblower
\textbf{Data.} $N=200$; $A=4.0\times10^{-2}\times3.0\times10^{-2}=1.2\times10^{-3}\ \mathrm{m^2}$; $I=\SI{1.5}{A}$; $B=\SI{0.25}{T}$.

\textbf{1. Maximum torque} (coil plane parallel to $\vec{B}$, $\theta=0$, $\cos\theta=1$):
\[ \tau_{\max}=BAIN=0.25\times1.2\times10^{-3}\times1.5\times200
=9.0\times10^{-2}\ \mathrm{N\,m}. \]
\[ \boxed{\tau_{\max}=\SI{0.090}{N.m}} \]

\textbf{2. Torque at $\theta=30^{\circ}$:}
\[ \tau=\tau_{\max}\cos 30^{\circ}=0.090\times0.866=7.8\times10^{-2}\ \mathrm{N\,m}. \]

\textbf{3. Role of the commutator.} After half a revolution, the sides of the coil have swapped positions relative to the poles; without any change, the forces would then produce a torque in the \emph{opposite} sense, and the coil would simply oscillate about the equilibrium position. The \textbf{split-ring commutator} reverses the current in the coil every half-turn, exactly when the coil passes through the equilibrium plane. The forces on each side then always act so as to produce a torque in the \textbf{same sense of rotation}, and the motor keeps turning. (In a practical motor, many coils at different angles smooth the torque.)
\end{exoresolu}

\begin{savaistu}[Electric motors around us]
The same physics turns the fan in a classroom in Yaound\'e, the water pump of a village borehole and the starter motor of a \emph{bendskin} (motorbike taxi). Conversely, a moving-coil meter uses $\tau=BAIN\cos\theta$ opposed by a hairspring: the deflection is proportional to the current, which is why the scale of an analogue ammeter is uniform.
\end{savaistu}

\courssub{Method 4 — Calculating B from the Biot–Savart / Ampère Results}

So far $\vec{B}$ was \emph{given}. The Biot–Savart law and Ampère's circuital law tell us how a current \emph{creates} the field; for the three standard geometries of the syllabus, the results below must be known and used directly.

\begin{methode}[Choosing the right formula for the right conductor]
\begin{enumerate}
\item \textbf{Recognise the geometry} and pick the corresponding result (obtained from the Biot–Savart law or Ampère's circuital law):
\begin{center}
\begin{tabularx}{\linewidth}{|l|X|}\hline
\rowcolor{popblueL} \textbf{Conductor} & \textbf{Flux density} \\ \hline
Long straight wire & $B=\dfrac{\mu_0 I}{2\pi r}$ at distance $r$ \\ \hline
Centre of a plane circular coil ($N$ turns, radius $r$) & $B=\dfrac{\mu_0 N I}{2r}$ \\ \hline
Inside a long solenoid ($n$ turns per metre) & $B=\mu_0 n I$ (uniform) \\ \hline
\end{tabularx}
\end{center}
with $\mu_0=4\pi\times10^{-7}\ \mathrm{H\,m^{-1}}$ (the permeability of free space).
\item \textbf{Find the direction} with the right-hand grip rule: thumb along the current, curled fingers give the sense of $\vec{B}$ around the wire; for a coil or solenoid, curled fingers follow the current and the thumb points to the North face.
\item \textbf{Convert units}: $n$ must be in turns per \emph{metre} (e.g. 800 turns over 40 cm gives $n=2000\ \mathrm{m^{-1}}$).
\item \textbf{Superpose if needed}: if several conductors contribute, add the fields \emph{vectorially}, taking directions into account.
\end{enumerate}
\end{methode}

\begin{attention}
Do not confuse the three formulae: the wire's field decreases as $1/r$, the coil's field is evaluated \emph{at its centre only}, and the solenoid's field is \emph{uniform inside} and practically zero outside. Writing $B=\mu_0 nI$ for a short coil is a classic error.
\end{attention}

\begin{exoresolu}[Field inside a solenoid and at the centre of a coil]
\begin{enumerate}
\item A solenoid of length $\SI{50}{cm}$ has $1000$ turns and carries a current of $\SI{2.0}{A}$. Calculate the flux density at its centre (assumed long).
\item The same current flows through a plane circular coil of $20$ turns and radius $\SI{10}{cm}$. Calculate the flux density at its centre.
\item The coil is placed inside the solenoid with its plane perpendicular to the solenoid's axis, the currents aiding. Find the total field at the centre.
\end{enumerate}
\tcblower
\textbf{1. Solenoid.} Number of turns per metre:
\[ n=\frac{N}{\ell}=\frac{1000}{0.50}=2000\ \mathrm{turns\,m^{-1}}. \]
\[ B_{\text{sol}}=\mu_0 n I = 4\pi\times10^{-7}\times 2000\times 2.0
= 5.03\times10^{-3}\ \mathrm{T}\approx \SI{5.0}{mT}. \]

\textbf{2. Circular coil.}
\[ B_{\text{coil}}=\frac{\mu_0 N I}{2r}
=\frac{4\pi\times10^{-7}\times 20\times 2.0}{2\times 0.10}
=2.51\times10^{-4}\ \mathrm{T}\approx \SI{0.25}{mT}. \]

\textbf{3. Superposition.} The plane of the coil is perpendicular to the solenoid axis, so the field at its centre is along that same axis. Since the currents aid, the fields add:
\[ B_{\text{tot}}=B_{\text{sol}}+B_{\text{coil}}=5.03\times10^{-3}+0.25\times10^{-3}
=5.28\times10^{-3}\ \mathrm{T}\approx \SI{5.3}{mT}, \]
directed along the axis of the solenoid, towards the North end given by the grip rule.

\textbf{Remark.} Note that the solenoid's field is about 20 times larger: multiplying turns along an axis is far more efficient than a single flat coil, which is why electromagnets are built as solenoids.
\end{exoresolu}

\courssub{Method 5 — Force Between Two Parallel Current-Carrying Wires}

Two neighbouring currents interact because each wire sits in the magnetic field \emph{created by the other}. Nothing new is needed: combine the field formula of Method 4 with the force formula of Method 1.

\begin{methode}[Combining $B=\mu_0 I/2\pi r$ and $F=BIL$]
\begin{enumerate}
\item \textbf{Field created by wire 1} at the position of wire 2 (separation $d$): $B_1=\dfrac{\mu_0 I_1}{2\pi d}$.
\item \textbf{Force on wire 2} (length $\ell$, perpendicular to $B_1$):
\[ F = B_1 I_2 \ell = \frac{\mu_0 I_1 I_2 \ell}{2\pi d}. \]
By Newton's third law, wire 1 feels an equal and opposite force.
\item \textbf{Direction}: use the grip rule for $B_1$, then Fleming's left hand for the force. Memorise the result: \textbf{currents in the same direction attract; opposite directions repel}.
\item \textbf{Force per unit length}: $F/\ell=\dfrac{\mu_0 I_1 I_2}{2\pi d}$ — this relation historically defines the ampere.
\end{enumerate}
\end{methode}

\begin{exoresolu}[Two bus-bars in a power installation]
Two long parallel copper bars in an electrical installation at a factory in Douala are separated by $d=\SI{8.0}{cm}$ and carry currents $I_1=\SI{400}{A}$ and $I_2=\SI{600}{A}$ in the \textbf{same direction}. Calculate the force per unit length between the bars and state whether it is attractive or repulsive. Comment on the practical consequence.
\tcblower
\textbf{Step 1 — Field of bar 1 at bar 2:}
\[ B_1=\frac{\mu_0 I_1}{2\pi d}=\frac{4\pi\times10^{-7}\times 400}{2\pi\times 0.080}
=\frac{2\times10^{-7}\times400}{0.080}=1.0\times10^{-3}\ \mathrm{T}. \]

\textbf{Step 2 — Force per unit length on bar 2:}
\[ \frac{F}{\ell}=B_1 I_2 = 1.0\times10^{-3}\times 600 = \SI{0.60}{N.m^{-1}}. \]
(Equivalently, $\dfrac{F}{\ell}=\dfrac{\mu_0 I_1 I_2}{2\pi d}=\dfrac{2\times10^{-7}\times400\times600}{0.080}=\SI{0.60}{N.m^{-1}}$.)

\textbf{Step 3 — Direction.} The currents are in the same direction, so the force is \textbf{attractive}: each bar is pulled towards the other.

\textbf{Comment.} $\SI{0.60}{N}$ per metre seems modest, but during a short-circuit the current can rise to tens of kiloamperes; since $F\propto I_1 I_2$, the force is then multiplied by thousands and can bend or tear the bars. This is why bus-bars must be rigidly clamped — a direct application of this formula in electrical engineering.
\end{exoresolu}

\courssub{Method 6 — A Moving Charge in a Uniform Field: Circular Motion and $e/m$}

A current is a flow of charges, so the force on a wire is the sum of the forces on the moving charges inside it. Isolating one charge leads to the magnetic part of the \cle{Lorentz force}, $\vec{F}=q\vec{v}\wedge\vec{B}$, and to one of the most elegant motions in physics: the circle.

\begin{methode}[Magnetic force as centripetal force]
\begin{enumerate}
\item \textbf{Magnitude of the force}: $F=qvB\sin\theta$. For a velocity perpendicular to $\vec{B}$, $F=qvB$, constant in magnitude and always perpendicular to $\vec{v}$.
\item \textbf{Consequence}: the force does no work ($\vec{F}\perp\vec{v}$), so the \textbf{speed and kinetic energy are constant}; the path is a \textbf{circle} (helix if $\vec{v}$ has a component parallel to $\vec{B}$).
\item \textbf{Radius}: equate the magnetic force to the centripetal force:
\[ qvB=\frac{mv^2}{r}\quad\Longrightarrow\quad r=\frac{mv}{qB}. \]
\item \textbf{Period}: $T=\dfrac{2\pi r}{v}=\dfrac{2\pi m}{qB}$ — independent of $v$ (basis of the cyclotron).
\item \textbf{Specific charge}: if the particle was accelerated from rest through a voltage $U$, $\tfrac12 mv^2=qU$; combining with $r=mv/qB$ gives
\[ \frac{q}{m}=\frac{2U}{B^2 r^2}. \]
\item \textbf{Direction}: use $\vec{F}=q\vec{v}\wedge\vec{B}$; for an \textbf{electron} ($q<0$) the force is \emph{opposite} to $\vec{v}\wedge\vec{B}$.
\end{enumerate}
\end{methode}

\begin{exoresolu}[Specific charge of the electron]
In a fine-beam tube, electrons are accelerated from rest through a potential difference $U=\SI{250}{V}$ and enter a uniform magnetic field $B=\SI{1.2}{mT}$ perpendicular to their velocity. The circular path has radius $r=\SI{7.7}{cm}$.
\begin{enumerate}
\item Show that the speed of the electrons is $v=\sqrt{2eU/m_e}$ and explain why it stays constant.
\item Derive $e/m_e=2U/(B^2r^2)$ and compute its value.
\item Calculate the period of revolution.
\end{enumerate}
\tcblower
\textbf{1. Speed.} The work done by the electric field becomes kinetic energy:
\[ eU=\tfrac12 m_e v^2 \quad\Longrightarrow\quad v=\sqrt{\frac{2eU}{m_e}}. \]
In the magnetic region the Lorentz force is always perpendicular to the velocity, so its power $\vec{F}\cdot\vec{v}$ is zero: no work is done and the speed (and kinetic energy) remain constant. Only the \emph{direction} of $\vec{v}$ changes.

\textbf{2. Specific charge.} In the field, the magnetic force provides the centripetal force:
\[ evB=\frac{m_e v^2}{r}\quad\Longrightarrow\quad v=\frac{eBr}{m_e}. \]
Substituting into $eU=\tfrac12 m_e v^2$:
\[ eU=\frac12 m_e\left(\frac{eBr}{m_e}\right)^2=\frac{e^2B^2r^2}{2m_e}
\quad\Longrightarrow\quad \boxed{\frac{e}{m_e}=\frac{2U}{B^2r^2}}. \]
Numerically:
\[ \frac{e}{m_e}=\frac{2\times250}{(1.2\times10^{-3})^2\times(7.7\times10^{-2})^2}
=\frac{500}{1.44\times10^{-6}\times5.93\times10^{-3}}
=\frac{500}{8.54\times10^{-9}}. \]
\[ \boxed{\frac{e}{m_e}\approx 5.9\times10^{10}\ \mathrm{C\,kg^{-1}}} \]
(Order of magnitude consistent with the accepted value $1.76\times10^{11}\ \mathrm{C\,kg^{-1}}$; the difference comes from the rounded data of this exercise.)

\textbf{3. Period.}
\[ T=\frac{2\pi r}{v}=\frac{2\pi m_e}{eB}
=\frac{2\pi}{(e/m_e)B}
=\frac{2\pi}{5.9\times10^{10}\times1.2\times10^{-3}}
\approx 8.9\times10^{-8}\ \mathrm{s}. \]
Note that $T$ does not depend on $r$ or $v$: faster electrons describe larger circles in the same time.
\end{exoresolu}

\courssub{Method 7 — The Hall Effect and Magnetic Materials}

When the moving charges are confined in a solid slab, the magnetic force pushes them sideways until an internal electric field stops them: the resulting transverse voltage is the \cle{Hall voltage}, a remarkably practical way of measuring $\vec{B}$.

\begin{methode}[Analysing a Hall-effect problem]
\begin{enumerate}
\item \textbf{Set up the geometry}: current $I$ along the slab, field $\vec{B}$ perpendicular to the slab, Hall voltage $U_H$ measured across the faces perpendicular to both.
\item \textbf{Force on the carriers}: $\vec{F}=q\vec{v}\wedge\vec{B}$ pushes the charge carriers to one face until the transverse electric field $E_H$ balances it:
\[ qE_H=qvB \quad\Longrightarrow\quad E_H=vB. \]
\item \textbf{Hall voltage}: $U_H=E_H\,d=vBd$, where $d$ is the distance between the faces.
\item \textbf{Link to the current}: $I=nqAv$ with $A$ the cross-section, so $v=I/(nqA)$ and
\[ U_H=\frac{IB}{nqt} \qquad (t=\text{thickness along }\vec{B}), \]
where $R_H=1/(nq)$ is the Hall coefficient.
\item \textbf{Uses}: measuring $B$ (Hall probe), determining the \emph{sign} and density $n$ of the carriers (negative for electrons in metals and n-type semiconductors, positive for holes in p-type).
\item \textbf{Materials}: classify the substance — \textbf{diamagnetic} (weakly repelled by a magnet, $\mu_r$ slightly $<1$), \textbf{paramagnetic} (weakly attracted, $\mu_r$ slightly $>1$), \textbf{ferromagnetic} (strongly magnetised, $\mu_r\gg1$, hysteresis, permanent magnets; e.g. iron, cobalt, nickel). For \textbf{magnetic shielding}, enclose the region in a high-permeability ferromagnetic shell which channels the field lines around the protected volume.
\end{enumerate}
\end{methode}

\begin{exoresolu}[Hall probe measuring a magnetic field]
A Hall probe uses a semiconductor slice of thickness $t=\SI{0.50}{mm}$ containing $n=1.0\times10^{23}\ \mathrm{m^{-3}}$ charge carriers (electrons, $e=1.60\times10^{-19}\ \mathrm{C}$). It carries a current $I=\SI{20}{mA}$. Placed perpendicular to an unknown magnetic field, it gives a Hall voltage $U_H=\SI{0.40}{mV}$.
\begin{enumerate}
\item Explain, with the aid of the Lorentz force, how the Hall voltage arises.
\item Calculate the flux density $B$.
\item State two advantages of determining the sign of $U_H$.
\end{enumerate}
\tcblower
\textbf{1. Origin of the Hall voltage.} The electrons drifting through the slice with velocity $\vec{v}$ experience the magnetic part of the Lorentz force $\vec{F}=-e\,\vec{v}\wedge\vec{B}$, which deflects them towards one face of the slice. Negative charge accumulates on that face, leaving the opposite face positive: a transverse electric field $E_H$ builds up. Equilibrium is reached when the electric force $eE_H$ exactly cancels the magnetic force $evB$; the resulting potential difference between the faces is the Hall voltage $U_H=E_H d$.

\textbf{2. Calculation of $B$.} At equilibrium $E_H=vB$, and $v=I/(neA)$ with $A$ the cross-section traversed by the current. Combining:
\[ U_H=\frac{IB}{net}\quad\Longrightarrow\quad B=\frac{net\,U_H}{I}. \]
\[ B=\frac{1.0\times10^{23}\times1.60\times10^{-19}\times0.50\times10^{-3}\times0.40\times10^{-3}}{20\times10^{-3}}. \]
Numerator: $1.0\times10^{23}\times1.60\times10^{-19}=1.6\times10^{4}$; then $1.6\times10^{4}\times0.50\times10^{-3}=8.0$; then $8.0\times0.40\times10^{-3}=3.2\times10^{-3}$.
\[ B=\frac{3.2\times10^{-3}}{20\times10^{-3}}=\SI{0.16}{T}. \]
\[ \boxed{B=\SI{0.16}{T}} \]

\textbf{3. Sign of $U_H$.}
\begin{itemize}
\item It reveals the \textbf{sign of the charge carriers}: negative (electrons) in metals and n-type semiconductors, positive (holes) in p-type semiconductors.
\item It gives the \textbf{direction (sense) of $\vec{B}$}, not just its magnitude, which a simple current balance cannot do directly.
\end{itemize}
\end{exoresolu}

\begin{aretenir}[Reflexes to remember for the examination]
\begin{itemize}
\item Force on a conductor: $F=BIL\sin\theta$; direction from Fleming's \textbf{left} hand.
\item Torque on a coil: $\tau=BAIN\cos\theta$; the split-ring commutator keeps a dc motor turning.
\item Fields: wire $\mu_0 I/2\pi r$; coil centre $\mu_0 NI/2r$; solenoid $\mu_0 nI$; direction from the \textbf{right}-hand grip rule.
\item Two parallel wires: $F/\ell=\mu_0 I_1 I_2/2\pi d$; like currents attract.
\item Moving charge: $r=mv/qB$, $T=2\pi m/qB$; the magnetic force never changes the speed.
\item Hall effect: $U_H=IB/(nqt)$; measures $B$ and identifies the carriers.
\end{itemize}
\end{aretenir}

\newpage

\coursec{Exercises}

\courssub{I check my knowledge}

\begin{exercice}[Multiple choice: magnetic flux density]
A straight conductor of length $\ell = \SI{0.20}{m}$ carrying a current $I = \SI{5.0}{A}$ is placed perpendicular to a uniform magnetic field and experiences a force $F = \SI{0.10}{N}$. The magnetic flux density is:
\begin{enumerate}
\item $B = \SI{0.010}{T}$
\item $B = \SI{0.10}{T}$
\item $B = \SI{1.0}{T}$
\item $B = \SI{10}{T}$
\end{enumerate}
Justify your choice with a short calculation.
\end{exercice}

\begin{exercice}[True or false? Justify]
For each of the following statements, say whether it is \textbf{true} or \textbf{false} and justify your answer in one or two sentences.
\begin{enumerate}
\item The force on a current-carrying conductor in a magnetic field is greatest when the conductor is parallel to the field.
\item The magnetic flux density at the centre of a plane circular coil of radius $r$ carrying a current $I$ is $B = \dfrac{\mu_0 I}{2r}$.
\item Inside a long solenoid, the magnetic field is uniform and proportional to the current.
\item Two parallel conductors carrying currents in the \emph{same} direction repel each other.
\item The magnetic force on a moving charge can change the kinetic energy of the charge.
\end{enumerate}
\end{exercice}

\begin{exercice}[Definitions and units]
\begin{enumerate}
\item Define the \cle{tesla} in terms of the force on a current-carrying conductor.
\item State the \cle{Biot--Savart law} in words, and give the expression for the field at the centre of a circular loop that follows from it.
\item Write down the expression for the \cle{Lorentz force} on a charge $q$ moving with velocity $\vec{v}$ in a region containing both an electric field $\vec{E}$ and a magnetic field $\vec{B}$.
\item Give the SI unit of the permeability of free space $\mu_0$ and its numerical value.
\end{enumerate}
\end{exercice}

\begin{exercice}[Direct calculation: field of a solenoid]
A long solenoid has $1200$ turns wound uniformly over a length of $\SI{60}{cm}$ and carries a current of $\SI{3.0}{A}$.
\begin{enumerate}
\item Calculate the magnetic flux density $B$ inside the solenoid.
\item Sketch the field pattern inside and outside the solenoid, indicating the direction of the field lines.
\item State one practical way of doubling $B$ without changing the current.
\end{enumerate}
Take $\mu_0 = 4\pi \times 10^{-7}\ \mathrm{H\,m^{-1}}$.
\end{exercice}

\courssub{I practise}

\begin{exercice}[Force on a conductor at an angle]
A straight conductor of length $\SI{25}{cm}$ carrying a current of $\SI{8.0}{A}$ is placed in a uniform magnetic field of flux density $\SI{0.40}{T}$, the conductor making an angle of $30^\circ$ with the field direction.
\begin{enumerate}
\item Calculate the magnitude of the force on the conductor.
\item The current flows horizontally from west to east and the field points horizontally from south to north in the plane containing the conductor. Using Fleming's left-hand rule, state the direction of the force.
\item Determine the orientation of the conductor, relative to the field, for which the force is zero. Explain your answer.
\end{enumerate}
\end{exercice}

\begin{exercice}[Comparing a solenoid and a straight wire]
A solenoid $\SI{40}{cm}$ long has $800$ turns and carries a current of $\SI{2.5}{A}$.
\begin{enumerate}
\item Calculate the magnetic flux density $B$ at the centre of the solenoid.
\item A long straight wire is placed $\SI{5.0}{cm}$ from a point $P$. Calculate the current that must flow in the wire so that it produces at $P$ the same flux density as found in (a).
\item Comment on the order of magnitude of this current and on the practical difficulty of producing large fields with a single straight wire.
\end{enumerate}
\end{exercice}

\begin{exercice}[Force between two parallel conductors]
Two long parallel conductors are placed $\SI{4.0}{cm}$ apart in air and each carries a current of $\SI{15}{A}$ in the same direction.
\begin{enumerate}
\item Calculate the magnetic flux density produced by one conductor at the position of the other.
\item Hence calculate the force per unit length between the conductors.
\item State, with a reason, whether the force is attractive or repulsive.
\item Explain briefly how this phenomenon is used in the definition of the \cle{ampere}.
\end{enumerate}
\end{exercice}

\begin{exercice}[Hall effect in a copper strip]
A copper strip of thickness $\SI{0.10}{mm}$ carries a current of $\SI{10}{A}$ along its length. A uniform magnetic field of flux density $\SI{1.5}{T}$ is applied perpendicular to the plane of the strip, and a Hall voltage of $\SI{22}{\micro V}$ is measured across its thickness.
\begin{enumerate}
\item Write down the expression relating the Hall voltage $V_H$ to $B$, $I$, $n$, $e$ and the thickness $t$, explaining each symbol.
\item Calculate the charge carrier density $n$ in copper.
\item Explain how the \emph{sign} of the Hall voltage allows one to identify the sign of the charge carriers, and state what would be observed for an $n$-type semiconductor.
\end{enumerate}
\end{exercice}

\courssub{I prepare for the GCE}

\begin{exercice}[Torque on a coil and the moving-coil galvanometer \hfill (14 marks)]
\textbf{(a)} A rectangular coil of $N$ turns, of sides $a$ and $b$, carries a current $I$ and is suspended in a uniform magnetic field of flux density $B$, the plane of the coil making an angle $\theta$ with the field direction.
\begin{enumerate}
\item[(i)] Draw a clear labelled diagram of the coil in the field and indicate the forces acting on each of its four sides. \hfill (3 marks)
\item[(ii)] Starting from the force on each side, show that the couple (torque) on the coil is
\[ T = N A I B \cos\theta \]
where $A$ is the area of the coil. State clearly what happens when the plane of the coil is perpendicular to the field. \hfill (5 marks)
\end{enumerate}
\textbf{(b)} In a moving-coil galvanometer the magnet is shaped so that the field is \emph{radial}.
\begin{enumerate}
\item[(i)] Explain what is meant by a radial field and why it makes the deflection of the coil directly proportional to the current. \hfill (3 marks)
\item[(ii)] A galvanometer coil has $200$ turns, an area of $2.0 \times 10^{-4}\ \mathrm{m^2}$ and lies in a radial field of $\SI{0.15}{T}$. The control spring exerts a restoring couple of $1.2 \times 10^{-6}\ \mathrm{N\,m}$ per radian. Calculate the steady deflection, in degrees, produced by a current of $\SI{50}{\micro A}$. \hfill (3 marks)
\end{enumerate}
\end{exercice}

\begin{exercice}[Specific charge of the electron \hfill (12 marks)]
In a fine-beam tube, electrons are accelerated from rest through a potential difference $V = \SI{250}{V}$ and then enter a region of uniform magnetic field of flux density $B$ perpendicular to their velocity. They describe a circular path of radius $r = \SI{4.4}{cm}$.
\begin{enumerate}
\item[(a)] Write down the equation expressing the kinetic energy gained by an electron in terms of $e$ and $V$. \hfill (2 marks)
\item[(b)] By equating the magnetic force to the centripetal force, show that
\[ \frac{e}{m_e} = \frac{2V}{B^2 r^2}. \]
\hfill (4 marks)
\item[(c)] The measured value of the flux density is $B = \SI{1.2}{mT}$. Calculate the value of $e/m_e$ obtained from this experiment. \hfill (3 marks)
\item[(d)] Compare your result with the accepted value $e/m_e = 1.76 \times 10^{11}\ \mathrm{C\,kg^{-1}}$ and suggest one possible source of systematic error in the experiment. \hfill (3 marks)
\end{enumerate}
\end{exercice}

\begin{exercice}[Crossed fields and the velocity selector \hfill (12 marks)]
A beam of ions of charge $+q$ and various speeds enters a region where a uniform electric field $E$ (vertical) and a uniform magnetic field $B$ (horizontal, into the page) are applied at right angles to each other and to the beam.
\begin{enumerate}
\item[(a)] Draw a labelled diagram showing the directions of $\vec{E}$, $\vec{B}$, the velocity $\vec{v}$ of the ions, and the electric and magnetic forces on a positive ion. \hfill (3 marks)
\item[(b)] Show that ions passing through the region undeflected satisfy
\[ v = \frac{E}{B}. \]
Explain what happens to ions that are (i) faster and (ii) slower than this value. \hfill (4 marks)
\item[(c)] The selected ions then enter a second region where only the magnetic field $B$ acts. Show that they follow a circular path and derive the expression for its radius in terms of $m$, $q$, $E$ and $B$. \hfill (3 marks)
\item[(d)] Explain briefly how this arrangement is used in a mass spectrometer to separate isotopes of the same element. \hfill (2 marks)
\end{enumerate}
\end{exercice}

\newpage

\coursec{Solutions to the exercises}

\coursec{Magnetic Fields}

\courssub{I check my knowledge}

\begin{corrige}[Exercise 1 --- Multiple choice: magnetic flux density]
For a conductor \emph{perpendicular} to the field, $F = B I \ell$, so
\[ B = \frac{F}{I \ell} = \frac{0.10}{5.0 \times 0.20} = 0.10\ \mathrm{T}. \]
The correct answer is \textbf{(b)}.
\end{corrige}

\begin{corrige}[Exercise 2 --- True or false? Justify]
\begin{enumerate}
\item \textbf{False.} $F = B I \ell \sin\theta$ is greatest when the conductor is \emph{perpendicular} to the field ($\theta = 90^\circ$); it is zero when the conductor is parallel to the field.
\item \textbf{True.} This is the result of integrating the Biot--Savart law round the loop (for one turn; for $N$ turns, $B = \mu_0 N I / 2r$).
\item \textbf{True.} $B = \mu_0 n I$ inside a long solenoid: the field is uniform (independent of position) and directly proportional to $I$.
\item \textbf{False.} Currents in the same direction \emph{attract}; currents in opposite directions repel (check with Fleming's left-hand rule on one wire placed in the field of the other).
\item \textbf{False.} The magnetic force $q\,\vec{v} \times \vec{B}$ is always perpendicular to the velocity, so it does no work: it changes the \emph{direction} of the velocity but never the kinetic energy.
\end{enumerate}
\end{corrige}

\begin{corrige}[Exercise 3 --- Definitions and units]
\begin{enumerate}
\item One tesla is the magnetic flux density of a field that exerts a force of one newton on a conductor of length one metre carrying a current of one ampere placed perpendicular to the field: $1\ \mathrm{T} = 1\ \mathrm{N\,A^{-1}\,m^{-1}}$.
\item The Biot--Savart law: the contribution $d\vec{B}$ of a current element $I\,d\vec{\ell}$ to the field at a point $P$ is proportional to $I\,d\ell \sin\theta / r^2$, where $r$ is the distance from the element to $P$ and $\theta$ the angle between $d\vec{\ell}$ and the line joining the element to $P$; the direction of $d\vec{B}$ is given by the right-hand (corkscrew) rule. For a circular loop of radius $r$ and $N$ turns, integration gives $B = \dfrac{\mu_0 N I}{2r}$ at the centre.
\item $\vec{F} = q\left(\vec{E} + \vec{v} \times \vec{B}\right)$ (Lorentz force).
\item $\mu_0 = 4\pi \times 10^{-7}\ \mathrm{H\,m^{-1}}$ (equivalently $\mathrm{T\,m\,A^{-1}}$ or $\mathrm{N\,A^{-2}}$).
\end{enumerate}
\end{corrige}

\begin{corrige}[Exercise 4 --- Direct calculation: field of a solenoid]
\begin{enumerate}
\item The number of turns per unit length is
\[ n = \frac{N}{L} = \frac{1200}{0.60} = 2.0 \times 10^{3}\ \mathrm{turns\,m^{-1}}. \]
Hence
\[ B = \mu_0 n I = 4\pi \times 10^{-7} \times 2.0 \times 10^{3} \times 3.0 = 7.5 \times 10^{-3}\ \mathrm{T} = 7.5\ \mathrm{mT}. \]
\item Inside: straight, parallel, equally spaced lines running along the axis (uniform field), directed from the south end to the north end \emph{inside} the solenoid. Outside: lines loop back from the north end to the south end, widely spaced (weak field), like the field of a bar magnet. The direction is given by the right-hand grip rule (curl fingers in current direction, thumb points to North pole).
\item Double the number of turns per unit length, e.g.\ wind $2400$ turns on the same $60\ \mathrm{cm}$ (or insert a soft-iron core, which increases $B$ by a large factor $\mu_r$).
\end{enumerate}
\end{corrige}

\courssub{I practise}

\begin{corrige}[Exercise 5 --- Force on a conductor at an angle]
\begin{enumerate}
\item The force on a conductor at angle $\theta$ to the field is
\[ F = B I \ell \sin\theta = 0.40 \times 8.0 \times 0.25 \times \sin 30^\circ = 0.80 \times 0.50 = 0.40\ \mathrm{N}. \]
\item Fleming's left-hand rule: first finger (Field) points north, second finger (Current) points east; the thumb (Motion/force) then points \textbf{vertically upwards}. The force is perpendicular to the plane containing $\vec{B}$ and the conductor.
\item The force is zero when $\sin\theta = 0$, i.e.\ when the conductor is \textbf{parallel (or antiparallel) to the field} ($\theta = 0^\circ$ or $180^\circ$). Then the current has no component perpendicular to $\vec{B}$, so no sideways push exists.
\end{enumerate}
\end{corrige}

\begin{corrige}[Exercise 6 --- Comparing a solenoid and a straight wire]
\begin{enumerate}
\item Turns per unit length:
\[ n = \frac{N}{L} = \frac{800}{0.40} = 2.0 \times 10^{3}\ \mathrm{turns\,m^{-1}}. \]
\[ B = \mu_0 n I = 4\pi \times 10^{-7} \times 2.0 \times 10^{3} \times 2.5 = 6.3 \times 10^{-3}\ \mathrm{T} = 6.3\ \mathrm{mT}. \]
\item For a long straight wire, $B = \dfrac{\mu_0 I}{2\pi r}$, so
\[ I = \frac{2\pi r B}{\mu_0} = \frac{2\pi \times 0.050 \times 6.3 \times 10^{-3}}{4\pi \times 10^{-7}} = \frac{0.050 \times 6.3 \times 10^{-3}}{2 \times 10^{-7}} \approx 1.6 \times 10^{3}\ \mathrm{A}. \]
\item A current of about $1600\ \mathrm{A}$ is enormous --- far beyond what ordinary laboratory wiring can carry without overheating. A single straight wire is therefore a very inefficient source of strong fields; a solenoid multiplies the effect of each ampere by the number of turns per metre, which is why practical electromagnets always use coils.
\end{enumerate}
\end{corrige}

\begin{corrige}[Exercise 7 --- Force between two parallel conductors]
\begin{enumerate}
\item Field produced by wire 1 at distance $d = 0.040\ \mathrm{m}$:
\[ B_1 = \frac{\mu_0 I_1}{2\pi d} = \frac{4\pi \times 10^{-7} \times 15}{2\pi \times 0.040} = \frac{2 \times 10^{-7} \times 15}{0.040} = 7.5 \times 10^{-5}\ \mathrm{T}. \]
\item The force per unit length on wire 2 carrying $I_2$ in this field (the wires are perpendicular to $B_1$):
\[ \frac{F}{\ell} = B_1 I_2 = 7.5 \times 10^{-5} \times 15 = 1.1 \times 10^{-3}\ \mathrm{N\,m^{-1}}. \]
Equivalently, $\dfrac{F}{\ell} = \dfrac{\mu_0 I_1 I_2}{2\pi d}$.
\item \textbf{Attractive.} Each wire sits in the field of the other; applying Fleming's left-hand rule (or the right-hand grip rule for the field direction) shows the force on each wire is directed \emph{towards} the other when the currents are in the same direction.
\item The ampere is defined through this force: one ampere is the constant current which, maintained in two straight parallel conductors of infinite length and negligible cross-section placed one metre apart in vacuum, produces a force of $2 \times 10^{-7}\ \mathrm{N}$ per metre of length between them (this follows directly from $\frac{F}{\ell} = \frac{\mu_0 I^2}{2\pi d}$ with $I = 1\ \mathrm{A}$, $d = 1\ \mathrm{m}$).
\end{enumerate}
\end{corrige}

\begin{corrige}[Exercise 8 --- Hall effect in a copper strip]
\begin{enumerate}
\item At equilibrium the magnetic force on the carriers is balanced by the electric force due to the Hall field:
\[ e\,\frac{V_H}{w} = e v_d B \quad\Longrightarrow\quad V_H = \frac{B I}{n e t}, \]
where $V_H$ is the Hall voltage, $B$ the flux density (perpendicular to the strip), $I$ the current, $n$ the charge-carrier density, $e$ the elementary charge and $t$ the thickness of the strip measured \emph{along the direction of} $\vec{B}$ (width $w$ cancels out).
\item Solving for $n$:
\[ n = \frac{B I}{e t V_H} = \frac{1.5 \times 10}{1.60 \times 10^{-19} \times 1.0 \times 10^{-4} \times 22 \times 10^{-6}}. \]
\[ n = \frac{15}{3.52 \times 10^{-28}} \approx 4.3 \times 10^{28}\ \mathrm{m^{-3}}. \]
This is of the expected order of magnitude for a metal (about one free electron per atom).
\item The sign of $V_H$ tells which side of the strip accumulates negative charge. For negative carriers (electrons) drifting opposite to the conventional current, the magnetic force pushes them to the same edge as positive carriers would be pushed \emph{away} from; the edge that becomes negative for electrons would become positive for positive carriers. Hence measuring the polarity of $V_H$ reveals the sign of the carriers. For an $n$-type semiconductor the carriers are electrons, so the Hall voltage has the \emph{same polarity} as for copper (but it is much larger, because $n$ is much smaller).
\end{enumerate}
\end{corrige}

\courssub{I prepare for the GCE}

\begin{corrige}[Exercise 9 --- Torque on a coil and the moving-coil galvanometer \hfill (14 marks)]
\textbf{(a)(i)} Diagram (3 marks): rectangle $ABCD$ with sides $a$ (horizontal, length of coil) and $b$ (width), field $\vec{B}$ horizontal and uniform, current $I$ marked round the coil, forces $F = B I a N$ shown on the two sides of length $a$ (perpendicular to $\vec{B}$), one up and one down, forming a couple; zero force on the sides of length $b$ when parallel to $\vec{B}$ component. Award: 1 mark for coil + field directions, 1 mark for forces on the two active sides with correct directions (Fleming's left-hand rule), 1 mark for showing the forces on the other two sides are zero or collinear (no moment).

\textbf{(a)(ii)} Proof (5 marks):
Let $\theta$ be the angle between the \textbf{normal to the coil} and the magnetic field $\vec{B}$.
\begin{itemize}
\item Force on each side of length $a$ (perpendicular to $\vec{B}$): $F = N B I a$ \hfill (1)
\item The two forces are equal, opposite and displaced: they form a couple of arm $b\sin\theta$, the perpendicular distance between the lines of action \hfill (1)
\item Torque of a couple $=$ one force $\times$ perpendicular separation:
\[ T = F \times b\sin\theta = N B I a\, b \sin\theta \hfill (1) \]
\item Since $A = ab$ is the area of the coil: $\boxed{T = N A I B \sin\theta}$ \hfill (1)
\item When the plane of the coil is perpendicular to the field, the normal is parallel to $\vec{B}$ ($\theta = 0^\circ$) and $T = 0$: the forces are collinear and the coil is in equilibrium (this is the neutral position) \hfill (1)
\end{itemize}
\textit{Note: If $\theta$ is defined as the angle between the \emph{plane of the coil} and $\vec{B}$, the torque is $T = N A I B \cos\theta$; both conventions are acceptable provided the definition is clearly stated.}

\textbf{(b)(i)} (3 marks): In a radial field the pole pieces are curved (with a cylindrical soft-iron core) so that the field lines always lie \emph{in the plane of the coil} whatever its angular position \hfill (1). Then the normal to the coil is always perpendicular to $\vec{B}$ ($\theta = 90^\circ$), so $\sin\theta = 1$ and $T = N A I B$ is independent of position \hfill (1). Against a spring of stiffness $c$, equilibrium gives $N A I B = c\alpha$, hence $\alpha = \dfrac{NAB}{c}\,I \propto I$: the scale is linear \hfill (1).

\textbf{(b)(ii)} (3 marks): At equilibrium, deflecting torque $=$ restoring torque:
\[ N A I B = c\,\alpha \quad\Longrightarrow\quad \alpha = \frac{N A I B}{c} = \frac{200 \times 2.0 \times 10^{-4} \times 50 \times 10^{-6} \times 0.15}{1.2 \times 10^{-6}}\ \mathrm{rad}. \hfill (1) \]
\[ \alpha = \frac{3.0 \times 10^{-7}}{1.2 \times 10^{-6}} = 0.25\ \mathrm{rad} \hfill (1) \qquad \alpha = 0.25 \times \frac{180}{\pi} \approx 14^\circ. \hfill (1) \]
\end{corrige}

\begin{corrige}[Exercise 10 --- Specific charge of the electron \hfill (12 marks)]
\textbf{(a)} (2 marks): The electron, charge $e$, accelerated from rest through $V$, gains kinetic energy equal to the work done by the field:
\[ \tfrac{1}{2} m_e v^2 = eV. \hfill (2) \]
\textbf{(b)} (4 marks): From (a), $v = \sqrt{2eV/m_e}$ \hfill (1). In the field, the magnetic force supplies the centripetal force:
\[ e v B = \frac{m_e v^2}{r} \quad\Longrightarrow\quad r = \frac{m_e v}{e B}. \hfill (1) \]
Substituting $v$:
\[ r = \frac{m_e}{eB}\sqrt{\frac{2eV}{m_e}} = \frac{1}{B}\sqrt{\frac{2 m_e V}{e}} \quad\Longrightarrow\quad r^2 = \frac{2 m_e V}{e B^2}, \hfill (1) \]
\[ \boxed{\frac{e}{m_e} = \frac{2V}{B^2 r^2}}. \hfill (1) \]
\textbf{(c)} (3 marks):
\[ \frac{e}{m_e} = \frac{2 \times 250}{(1.2 \times 10^{-3})^2 \times (4.4 \times 10^{-2})^2} = \frac{500}{1.44 \times 10^{-6} \times 1.936 \times 10^{-3}}. \hfill (1) \]
\[ \frac{e}{m_e} = \frac{500}{2.79 \times 10^{-9}} \approx 1.8 \times 10^{11}\ \mathrm{C\,kg^{-1}}. \hfill (2) \]
\textbf{(d)} (3 marks): The result $1.8 \times 10^{11}\ \mathrm{C\,kg^{-1}}$ agrees with the accepted value $1.76 \times 10^{11}\ \mathrm{C\,kg^{-1}}$ to within about $2\%$ \hfill (1). Possible systematic errors (any one, 2 marks): the radius is measured on a faint beam with parallax error, tending to overestimate $r$; the field may not be perfectly uniform over the path, or the Helmholtz-coil calibration may be off; the accelerating voltage may drop slightly under load; the electrons may retain a small residual gas collision loss, reducing $v$.
\end{corrige}

\begin{corrige}[Exercise 11 --- Crossed fields and the velocity selector \hfill (12 marks)]
\textbf{(a)} (3 marks): Diagram: beam travelling left to right ($\vec{v}$ horizontal); $\vec{B}$ into the page ($\times$ symbols); $\vec{E}$ vertical, say downwards. On a positive ion: electric force $\vec{F}_E = q\vec{E}$ parallel to $\vec{E}$ (downwards); magnetic force $\vec{F}_B = q\vec{v}\times\vec{B}$ opposite (upwards). Award: 1 mark for the three vector directions mutually perpendicular, 1 mark for $F_E$ along $\vec{E}$, 1 mark for $F_B$ opposite to $F_E$ (use of right-hand/left-hand rule correctly stated).

\textbf{(b)} (4 marks): Undeflected motion requires zero net transverse force:
\[ qE = qvB \quad\Longrightarrow\quad \boxed{v = \frac{E}{B}}. \hfill (2) \]
(i) Ions faster than $E/B$: $qvB > qE$, the magnetic force wins, so they are deflected towards the magnetic-force side and strike the slit plate \hfill (1). (ii) Slower ions: $qE > qvB$, they are deflected the other way and are also stopped \hfill (1). Only ions with $v = E/B$ emerge, independent of $q$ and $m$.

\textbf{(c)} (3 marks): In the second region only $\vec{B}$ acts; the force $qvB$ is always perpendicular to $\vec{v}$, so its magnitude is constant and it acts as a centripetal force: the path is a circle \hfill (1). Then
\[ qvB = \frac{mv^2}{r} \quad\Longrightarrow\quad r = \frac{mv}{qB}, \hfill (1) \]
and with $v = E/B$:
\[ \boxed{r = \frac{mE}{qB^2}}. \hfill (1) \]

\textbf{(d)} (2 marks): Isotopes of the same element have the same charge $q$ but different masses $m$. Since $r = mE/(qB^2)$ with $E$, $B$, $q$ identical for all selected ions, each isotope follows a semicircle of different radius and strikes the detector (photographic plate or collector) at a different position \hfill (1); the relative intensities of the traces give the relative abundances \hfill (1).
\end{corrige}

\newpage

\coursec{Summary sheet}

\begin{popfigure}
\begin{tikzpicture}[mindmap, grow cyclic,
  every node/.style={concept, text=white, font=\small},
  root concept/.append style={concept color=popink, font=\bfseries, minimum size=3.2cm},
  level 1/.append style={sibling angle=60, level distance=3.6cm, minimum size=2.4cm, font=\scriptsize}]
\node[root concept]{Magnetic Fields}
  child[concept color=popblue]{ node {Fields \& flux density $B$ (tesla)} }
  child[concept color=poporange]{ node {Force on conductor $F=BIL\sin\theta$} }
  child[concept color=popgreen]{ node {Torque $\tau=NAIB$ \& motors} }
  child[concept color=poppurple]{ node {Sources: Biot--Savart, Amp\`ere} }
  child[concept color=poppink]{ node {Moving charges $F=qvB$} }
  child[concept color=popgold]{ node {Hall effect \& materials} };
\end{tikzpicture}
\legende{Mind map of the chapter Magnetic Fields.}
\end{popfigure}

\begin{bilan}
\textbf{Key definitions}
\begin{itemize}
\item \cle{Magnetic flux density} $B$: strength of a magnetic field; unit the \cle{tesla} (\SI{}{T}), with $1\ \mathrm{T}=1\ \mathrm{N\,A^{-1}\,m^{-1}}$.
\item Field patterns: concentric circles around a \textbf{straight wire} (right-hand grip rule); a \textbf{plane circular coil} gives a dipole-like field; a \textbf{solenoid} gives a nearly uniform field inside, like a bar magnet outside.
\item \cle{Lorentz force}: total force on a charge in electric and magnetic fields, $\vec{F}=q\vec{E}+q\vec{v}\wedge\vec{B}$.
\item \cle{Magnetic shielding}: enclosing a region in high-permeability (soft iron / mu-metal) material so field lines pass through the material, not the region.
\end{itemize}

\textbf{Laws and formulas to know}
\begin{itemize}
\item Force on a current-carrying conductor: $F=BIL\sin\theta$; direction from \cle{Fleming's left-hand rule} (First finger Field, seCond finger Current, thuMb Motion).
\item Torque on a rectangular coil of $N$ turns, area $A$: $\tau=NAIB$ (maximum when the coil plane is parallel to $\vec B$). Principle of d.c.\ motors (split-ring commutator) and a.c.\ motors.
\item Biot--Savart and Amp\`ere's law $\oint\vec B\cdot d\vec l=\mu_0 I_{\text{enclosed}}$ give:
\[
B=\frac{\mu_0 I}{2\pi r}\ \text{(long wire)},\qquad
B=\frac{\mu_0 I}{2r}\ \text{(centre of coil)},\qquad
B=\mu_0 n I\ \text{(long solenoid)}.
\]
\item Force per unit length between two parallel wires distance $d$ apart: $\dfrac{F}{L}=\dfrac{\mu_0 I_1 I_2}{2\pi d}$; attractive if currents are parallel, repulsive if antiparallel (defines the ampere).
\item Force on a moving charge: $F=qvB\sin\theta$; perpendicular to $\vec v$, so it does \textbf{no work}; the path is a circle of radius $r=\dfrac{mv}{qB}$, period $T=\dfrac{2\pi m}{qB}$.
\item Specific charge measurement (fine-beam tube): $\dfrac{e}{m_e}=\dfrac{2V}{B^2 r^2}$; accepted value $\approx 1.76\times 10^{11}\ \mathrm{C\,kg^{-1}}$.
\item Hall effect: Hall voltage $V_H=\dfrac{IB}{nqt}$ (or $V_H=Bvd$); sign of $V_H$ reveals the sign of the charge carriers; used to measure $B$ and carrier density.
\item Magnetic materials: \textbf{diamagnetic} (weakly repelled), \textbf{paramagnetic} (weakly attracted), \textbf{ferromagnetic} (strongly magnetised, e.g.\ iron; used for cores and shielding).
\end{itemize}

\textbf{Methods}
\begin{itemize}
\item Always draw the field direction first, then apply the left-hand rule (currents) or the rule for $q\vec v\wedge\vec B$ (charges; reverse for negative charges).
\item For circular motion of charges: set $qvB=\dfrac{mv^2}{r}$ and solve.
\item For Amp\`ere's law: choose a loop respecting the symmetry (circle around a wire, rectangle through a solenoid).
\end{itemize}

\textbf{Common mistakes}
\begin{itemize}
\item Forgetting $\sin\theta$ when the conductor or velocity is not perpendicular to $\vec B$.
\item Using the left-hand rule for a \emph{negative} charge without reversing the direction.
\item Confusing $n$ (turns per metre, $\mathrm{m^{-1}}$) with $N$ (total turns) in $B=\mu_0 nI$.
\item Mixing up the three field formulas (wire, coil, solenoid) and forgetting $\mu_0=4\pi\times10^{-7}\ \mathrm{H\,m^{-1}}$.
\item Believing the magnetic force accelerates a charge along its path: it only changes the \emph{direction} of $\vec v$.
\end{itemize}

\textbf{GCE tips}
\begin{itemize}
\item State the rule you use (Fleming, right-hand grip) explicitly --- marks are awarded for it.
\item Carry SI units through every line; convert coil dimensions to metres and turns per \emph{metre}.
\item In structured questions, derive $r=mv/qB$ from $qvB=mv^2/r$ rather than quoting it.
\item Sketches of field patterns (solenoid, two wires) must show direction arrows and correct spacing (lines closer where $B$ is stronger).
\end{itemize}
\end{bilan}

\findecours

\end{document}
